% ============================================================================
% n2dbn.tex
%
% Nested non-stationary dynamic Bayesian networks: the model, the learning
% procedure, the benchmark and the results, in one document.
%
% Build:
%   docker compose run --rm paper
% or, from paper/:
%   latexmk -pdf n2dbn.tex
%
% Sections live in sec_*.tex and are input below in reading order.
% ============================================================================

\documentclass[letterpaper,10pt]{article}

\usepackage[submission]{probml}
\usepackage{algorithm}
\usepackage{algpseudocode}
\usepackage{booktabs}
\usepackage{longtable}
\usepackage{array}
\usepackage{amsmath,amssymb,amsthm}
\usepackage{graphicx}
\ShortHeadings{Nested Non-Stationary Dynamic Bayesian Networks}

%!TEX root = ./main.tex

\usepackage{bm,amssymb,amsmath}

% Make theorem titles bold
\makeatletter
\def\th@plain{%
  \thm@notefont{}% same as heading font
  \itshape % body font
}
\def\th@definition{%
  \thm@notefont{}% same as heading font
  \normalfont % body font
}
\makeatother

\newcommand{\overbar}[1]{\mkern 1.5mu\overline{\mkern-1.5mu#1\mkern-1.5mu}\mkern 1.5mu}


% Mark sections of captions for referring to divisions of figures
\newcommand{\figleft}{{\em (Left)}}
\newcommand{\figcenter}{{\em (Center)}}
\newcommand{\figright}{{\em (Right)}}
\newcommand{\figtop}{{\em (Top)}}
\newcommand{\figbottom}{{\em (Bottom)}}
\newcommand{\captiona}{{\em (a)}}
\newcommand{\captionb}{{\em (b)}}
\newcommand{\captionc}{{\em (c)}}
\newcommand{\captiond}{{\em (d)}}

% Highlight a newly defined term
\newcommand{\newterm}[1]{{\bf #1}}


% Figure reference, lower-case.
\def\figref#1{figure~\ref{#1}}
% Figure reference, capital. For start of sentence
\def\Figref#1{Figure~\ref{#1}}
\def\twofigref#1#2{figures \ref{#1} and \ref{#2}}
\def\quadfigref#1#2#3#4{figures \ref{#1}, \ref{#2}, \ref{#3} and \ref{#4}}
% Section reference, lower-case.
\def\secref#1{section~\ref{#1}}
% Section reference, capital.
\def\Secref#1{Section~\ref{#1}}
% Reference to two sections.
\def\twosecrefs#1#2{sections \ref{#1} and \ref{#2}}
% Reference to three sections.
\def\secrefs#1#2#3{sections \ref{#1}, \ref{#2} and \ref{#3}}
% Reference to a chapter, lower-case.
\def\chapref#1{chapter~\ref{#1}}
% Reference to an equation, upper case.
\def\Chapref#1{Chapter~\ref{#1}}
% Reference to a range of chapters
\def\rangechapref#1#2{chapters\ref{#1}--\ref{#2}}
% Reference to an algorithm, lower-case.
\def\algref#1{algorithm~\ref{#1}}
% Reference to an algorithm, upper case.
\def\Algref#1{Algorithm~\ref{#1}}
\def\twoalgref#1#2{algorithms \ref{#1} and \ref{#2}}
\def\Twoalgref#1#2{Algorithms \ref{#1} and \ref{#2}}
% Reference to a part, lower case
\def\partref#1{part~\ref{#1}}
% Reference to a part, upper case
\def\Partref#1{Part~\ref{#1}}
\def\twopartref#1#2{parts \ref{#1} and \ref{#2}}

\def\ceil#1{\lceil #1 \rceil}
\def\floor#1{\lfloor #1 \rfloor}
\def\1{\bm{1}}
\newcommand{\train}{\mathcal{D}}
\newcommand{\valid}{\mathcal{D_{\mathrm{valid}}}}
\newcommand{\test}{\mathcal{D_{\mathrm{test}}}}

\def\eps{{\epsilon}}


% Random variables
\def\reta{{\textnormal{$\eta$}}}
\def\ra{{\textnormal{a}}}
\def\rb{{\textnormal{b}}}
\def\rc{{\textnormal{c}}}
\def\rd{{\textnormal{d}}}
\def\re{{\textnormal{e}}}
\def\rf{{\textnormal{f}}}
\def\rg{{\textnormal{g}}}
\def\rh{{\textnormal{h}}}
\def\ri{{\textnormal{i}}}
\def\rj{{\textnormal{j}}}
\def\rk{{\textnormal{k}}}
\def\rl{{\textnormal{l}}}
% rm is already a command, just don't name any random variables m
\def\rn{{\textnormal{n}}}
\def\ro{{\textnormal{o}}}
\def\rp{{\textnormal{p}}}
\def\rq{{\textnormal{q}}}
\def\rr{{\textnormal{r}}}
\def\rs{{\textnormal{s}}}
\def\rt{{\textnormal{t}}}
\def\ru{{\textnormal{u}}}
\def\rv{{\textnormal{v}}}
\def\rw{{\textnormal{w}}}
\def\rx{{\textnormal{x}}}
\def\ry{{\textnormal{y}}}
\def\rz{{\textnormal{z}}}

% Random vectors
\def\rvepsilon{{\mathbf{\epsilon}}}
\def\rvtheta{{\mathbf{\theta}}}
\def\rvpsi{{\boldsymbol{\psi}}}
\def\rvvarrho{{\boldsymbol{\varrho}}}
\def\rva{{\mathbf{a}}}
\def\rvb{{\mathbf{b}}}
\def\rvc{{\mathbf{c}}}
\def\rvd{{\mathbf{d}}}
\def\rve{{\mathbf{e}}}
\def\rvf{{\mathbf{f}}}
\def\rvg{{\mathbf{g}}}
\def\rvh{{\mathbf{h}}}
\def\rvu{{\mathbf{i}}}
\def\rvj{{\mathbf{j}}}
\def\rvk{{\mathbf{k}}}
\def\rvl{{\mathbf{l}}}
\def\rvm{{\mathbf{m}}}
\def\rvn{{\mathbf{n}}}
\def\rvo{{\mathbf{o}}}
\def\rvp{{\mathbf{p}}}
\def\rvq{{\mathbf{q}}}
\def\rvr{{\mathbf{r}}}
\def\rvs{{\mathbf{s}}}
\def\rvt{{\mathbf{t}}}
\def\rvu{{\mathbf{u}}}
\def\rvv{{\mathbf{v}}}
\def\rvw{{\mathbf{w}}}
\def\rvx{{\mathbf{x}}}
\def\rvy{{\mathbf{y}}}
\def\rvz{{\mathbf{z}}}

% Elements of random vectors
\def\erva{{\textnormal{a}}}
\def\ervb{{\textnormal{b}}}
\def\ervc{{\textnormal{c}}}
\def\ervd{{\textnormal{d}}}
\def\erve{{\textnormal{e}}}
\def\ervf{{\textnormal{f}}}
\def\ervg{{\textnormal{g}}}
\def\ervh{{\textnormal{h}}}
\def\ervi{{\textnormal{i}}}
\def\ervj{{\textnormal{j}}}
\def\ervk{{\textnormal{k}}}
\def\ervl{{\textnormal{l}}}
\def\ervm{{\textnormal{m}}}
\def\ervn{{\textnormal{n}}}
\def\ervo{{\textnormal{o}}}
\def\ervp{{\textnormal{p}}}
\def\ervq{{\textnormal{q}}}
\def\ervr{{\textnormal{r}}}
\def\ervs{{\textnormal{s}}}
\def\ervt{{\textnormal{t}}}
\def\ervu{{\textnormal{u}}}
\def\ervv{{\textnormal{v}}}
\def\ervw{{\textnormal{w}}}
\def\ervx{{\textnormal{x}}}
\def\ervy{{\textnormal{y}}}
\def\ervz{{\textnormal{z}}}

% Random matrices
\def\rmA{{\mathbf{A}}}
\def\rmB{{\mathbf{B}}}
\def\rmC{{\mathbf{C}}}
\def\rmD{{\mathbf{D}}}
\def\rmE{{\mathbf{E}}}
\def\rmF{{\mathbf{F}}}
\def\rmG{{\mathbf{G}}}
\def\rmH{{\mathbf{H}}}
\def\rmI{{\mathbf{I}}}
\def\rmJ{{\mathbf{J}}}
\def\rmK{{\mathbf{K}}}
\def\rmL{{\mathbf{L}}}
\def\rmM{{\mathbf{M}}}
\def\rmN{{\mathbf{N}}}
\def\rmO{{\mathbf{O}}}
\def\rmP{{\mathbf{P}}}
\def\rmQ{{\mathbf{Q}}}
\def\rmR{{\mathbf{R}}}
\def\rmS{{\mathbf{S}}}
\def\rmT{{\mathbf{T}}}
\def\rmU{{\mathbf{U}}}
\def\rmV{{\mathbf{V}}}
\def\rmW{{\mathbf{W}}}
\def\rmX{{\mathbf{X}}}
\def\rmY{{\mathbf{Y}}}
\def\rmZ{{\mathbf{Z}}}

% Elements of random matrices
\def\ermA{{\textnormal{A}}}
\def\ermB{{\textnormal{B}}}
\def\ermC{{\textnormal{C}}}
\def\ermD{{\textnormal{D}}}
\def\ermE{{\textnormal{E}}}
\def\ermF{{\textnormal{F}}}
\def\ermG{{\textnormal{G}}}
\def\ermH{{\textnormal{H}}}
\def\ermI{{\textnormal{I}}}
\def\ermJ{{\textnormal{J}}}
\def\ermK{{\textnormal{K}}}
\def\ermL{{\textnormal{L}}}
\def\ermM{{\textnormal{M}}}
\def\ermN{{\textnormal{N}}}
\def\ermO{{\textnormal{O}}}
\def\ermP{{\textnormal{P}}}
\def\ermQ{{\textnormal{Q}}}
\def\ermR{{\textnormal{R}}}
\def\ermS{{\textnormal{S}}}
\def\ermT{{\textnormal{T}}}
\def\ermU{{\textnormal{U}}}
\def\ermV{{\textnormal{V}}}
\def\ermW{{\textnormal{W}}}
\def\ermX{{\textnormal{X}}}
\def\ermY{{\textnormal{Y}}}
\def\ermZ{{\textnormal{Z}}}

% Vectors
\def\vzero{{\bm{0}}}
\def\vone{{\bm{1}}}
% \def\vmu{{\mathbf{\mu}}}
\def\vmu{{\boldsymbol{\mu}}}
\def\valpha{{\boldsymbol{\alpha}}}
\def\vbeta{{\boldsymbol{\beta}}}
\def\vgamma{{\boldsymbol{\gamma}}}
\def\vepsilon{{\boldsymbol{\epsilon}}}
\def\vtheta{{\boldsymbol{\theta}}}
\def\va{{\bm{a}}}
\def\vb{{\bm{b}}}
\def\vc{{\bm{c}}}
\def\vd{{\bm{d}}}
\def\ve{{\bm{e}}}
\def\vf{{\bm{f}}}
\def\vg{{\bm{g}}}
\def\vh{{\bm{h}}}
\def\vi{{\bm{i}}}
\def\vj{{\bm{j}}}
\def\vk{{\bm{k}}}
\def\vl{{\bm{l}}}
\def\vm{{\bm{m}}}
\def\vn{{\bm{n}}}
\def\vo{{\bm{o}}}
\def\vp{{\bm{p}}}
\def\vq{{\bm{q}}}
\def\vr{{\bm{r}}}
\def\vs{{\bm{s}}}
\def\vt{{\bm{t}}}
\def\vu{{\bm{u}}}
\def\vv{{\bm{v}}}
\def\vw{{\bm{w}}}
\def\vx{{\bm{x}}}
\def\vy{{\bm{y}}}
\def\vz{{\bm{z}}}

% Elements of vectors
\def\evalpha{{\alpha}}
\def\evbeta{{\beta}}
\def\evepsilon{{\epsilon}}
\def\evlambda{{\lambda}}
\def\evomega{{\omega}}
\def\evmu{{\mu}}
\def\evpsi{{\psi}}
\def\evsigma{{\sigma}}
\def\evtheta{{\theta}}
\def\eva{{a}}
\def\evb{{b}}
\def\evc{{c}}
\def\evd{{d}}
\def\eve{{e}}
\def\evf{{f}}
\def\evg{{g}}
\def\evh{{h}}
\def\evi{{i}}
\def\evj{{j}}
\def\evk{{k}}
\def\evl{{l}}
\def\evm{{m}}
\def\evn{{n}}
\def\evo{{o}}
\def\evp{{p}}
\def\evq{{q}}
\def\evr{{r}}
\def\evs{{s}}
\def\evt{{t}}
\def\evu{{u}}
\def\evv{{v}}
\def\evw{{w}}
\def\evx{{x}}
\def\evy{{y}}
\def\evz{{z}}

% Matrix
\def\mA{{\bm{A}}}
\def\mB{{\bm{B}}}
\def\mC{{\bm{C}}}
\def\mD{{\bm{D}}}
\def\mE{{\bm{E}}}
\def\mF{{\bm{F}}}
\def\mG{{\bm{G}}}
\def\mH{{\bm{H}}}
\def\mI{{\bm{I}}}
\def\mJ{{\bm{J}}}
\def\mK{{\bm{K}}}
\def\mL{{\bm{L}}}
\def\mM{{\bm{M}}}
\def\mN{{\bm{N}}}
\def\mO{{\bm{O}}}
\def\mP{{\bm{P}}}
\def\mQ{{\bm{Q}}}
\def\mR{{\bm{R}}}
\def\mS{{\bm{S}}}
\def\mT{{\bm{T}}}
\def\mU{{\bm{U}}}
\def\mV{{\bm{V}}}
\def\mW{{\bm{W}}}
\def\mX{{\bm{X}}}
\def\mY{{\bm{Y}}}
\def\mZ{{\bm{Z}}}
\def\mBeta{{\boldsymbol{\beta}}}
\def\mPhi{{\boldsymbol{\Phi}}}
\def\mLambda{{\boldsymbol{\Lambda}}}
\def\mSigma{{\boldsymbol{\Sigma}}}

% Tensor
\DeclareMathAlphabet{\mathsfit}{\encodingdefault}{\sfdefault}{m}{sl}
\SetMathAlphabet{\mathsfit}{bold}{\encodingdefault}{\sfdefault}{bx}{n}
\newcommand{\tens}[1]{\bm{\mathsfit{#1}}}
\def\tA{{\tens{A}}}
\def\tB{{\tens{B}}}
\def\tC{{\tens{C}}}
\def\tD{{\tens{D}}}
\def\tE{{\tens{E}}}
\def\tF{{\tens{F}}}
\def\tG{{\tens{G}}}
\def\tH{{\tens{H}}}
\def\tI{{\tens{I}}}
\def\tJ{{\tens{J}}}
\def\tK{{\tens{K}}}
\def\tL{{\tens{L}}}
\def\tM{{\tens{M}}}
\def\tN{{\tens{N}}}
\def\tO{{\tens{O}}}
\def\tP{{\tens{P}}}
\def\tQ{{\tens{Q}}}
\def\tR{{\tens{R}}}
\def\tS{{\tens{S}}}
\def\tT{{\tens{T}}}
\def\tU{{\tens{U}}}
\def\tV{{\tens{V}}}
\def\tW{{\tens{W}}}
\def\tX{{\tens{X}}}
\def\tY{{\tens{Y}}}
\def\tZ{{\tens{Z}}}


% Graph
\def\gA{{\mathcal{A}}}
\def\gB{{\mathcal{B}}}
\def\gC{{\mathcal{C}}}
\def\gD{{\mathcal{D}}}
\def\gE{{\mathcal{E}}}
\def\gF{{\mathcal{F}}}
\def\gG{{\mathcal{G}}}
\def\gH{{\mathcal{H}}}
\def\gI{{\mathcal{I}}}
\def\gJ{{\mathcal{J}}}
\def\gK{{\mathcal{K}}}
\def\gL{{\mathcal{L}}}
\def\gM{{\mathcal{M}}}
\def\gN{{\mathcal{N}}}
\def\gO{{\mathcal{O}}}
\def\gP{{\mathcal{P}}}
\def\gQ{{\mathcal{Q}}}
\def\gR{{\mathcal{R}}}
\def\gS{{\mathcal{S}}}
\def\gT{{\mathcal{T}}}
\def\gU{{\mathcal{U}}}
\def\gV{{\mathcal{V}}}
\def\gW{{\mathcal{W}}}
\def\gX{{\mathcal{X}}}
\def\gY{{\mathcal{Y}}}
\def\gZ{{\mathcal{Z}}}

% Sets
\def\sA{{\mathbb{A}}}
\def\sB{{\mathbb{B}}}
\def\sC{{\mathbb{C}}}
\def\sD{{\mathbb{D}}}
% Don't use a set called E, because this would be the same as our symbol
% for expectation.
\def\sF{{\mathbb{F}}}
\def\sG{{\mathbb{G}}}
\def\sH{{\mathbb{H}}}
\def\sI{{\mathbb{I}}}
\def\sJ{{\mathbb{J}}}
\def\sK{{\mathbb{K}}}
\def\sL{{\mathbb{L}}}
\def\sM{{\mathbb{M}}}
\def\sN{{\mathbb{N}}}
\def\sO{{\mathbb{O}}}
\def\sP{{\mathbb{P}}}
\def\sQ{{\mathbb{Q}}}
\def\sR{{\mathbb{R}}}
\def\sS{{\mathbb{S}}}
\def\sT{{\mathbb{T}}}
\def\sU{{\mathbb{U}}}
\def\sV{{\mathbb{V}}}
\def\sW{{\mathbb{W}}}
\def\sX{{\mathbb{X}}}
\def\sY{{\mathbb{Y}}}
\def\sZ{{\mathbb{Z}}}

% Entries of a matrix
\def\emLambda{{\Lambda}}
\def\emA{{A}}
\def\emB{{B}}
\def\emC{{C}}
\def\emD{{D}}
\def\emE{{E}}
\def\emF{{F}}
\def\emG{{G}}
\def\emH{{H}}
\def\emI{{I}}
\def\emJ{{J}}
\def\emK{{K}}
\def\emL{{L}}
\def\emM{{M}}
\def\emN{{N}}
\def\emO{{O}}
\def\emP{{P}}
\def\emQ{{Q}}
\def\emR{{R}}
\def\emS{{S}}
\def\emT{{T}}
\def\emU{{U}}
\def\emV{{V}}
\def\emW{{W}}
\def\emX{{X}}
\def\emY{{Y}}
\def\emZ{{Z}}
\def\emSigma{{\Sigma}}

% entries of a tensor
% Same font as tensor, without \bm wrapper
\newcommand{\etens}[1]{\mathsfit{#1}}
\def\etLambda{{\etens{\Lambda}}}
\def\etA{{\etens{A}}}
\def\etB{{\etens{B}}}
\def\etC{{\etens{C}}}
\def\etD{{\etens{D}}}
\def\etE{{\etens{E}}}
\def\etF{{\etens{F}}}
\def\etG{{\etens{G}}}
\def\etH{{\etens{H}}}
\def\etI{{\etens{I}}}
\def\etJ{{\etens{J}}}
\def\etK{{\etens{K}}}
\def\etL{{\etens{L}}}
\def\etM{{\etens{M}}}
\def\etN{{\etens{N}}}
\def\etO{{\etens{O}}}
\def\etP{{\etens{P}}}
\def\etQ{{\etens{Q}}}
\def\etR{{\etens{R}}}
\def\etS{{\etens{S}}}
\def\etT{{\etens{T}}}
\def\etU{{\etens{U}}}
\def\etV{{\etens{V}}}
\def\etW{{\etens{W}}}
\def\etX{{\etens{X}}}
\def\etY{{\etens{Y}}}
\def\etZ{{\etens{Z}}}

% The true underlying data generating distribution
\newcommand{\pdata}{p_{\rm{data}}}
% The empirical distribution defined by the training set
\newcommand{\ptrain}{\widehat{p}_{\rm{data}}}
\newcommand{\Ptrain}{\widehat{P}_{\rm{data}}}
% The model distribution
\newcommand{\pmodel}{p_{\rm{model}}}
\newcommand{\Pmodel}{P_{\rm{model}}}
\newcommand{\ptildemodel}{\widetilde{p}_{\rm{model}}}
% Stochastic autoencoder distributions
\newcommand{\pencode}{p_{\rm{encoder}}}
\newcommand{\pdecode}{p_{\rm{decoder}}}
\newcommand{\precons}{p_{\rm{reconstruct}}}

\newcommand{\laplace}{\mathrm{Laplace}} % Laplace distribution

\newcommand{\E}{\mathbb{E}}
\newcommand{\Ls}{\mathcal{L}}
\newcommand{\R}{\mathbb{R}}
\newcommand{\emp}{\widetilde{p}}
\newcommand{\lr}{\alpha}
\newcommand{\reg}{\lambda}
\newcommand{\rect}{\mathrm{rectifier}}
\newcommand{\softmax}{\mathrm{softmax}}
\newcommand{\sigmoid}{\sigma}
\newcommand{\softplus}{\zeta}
\newcommand{\KL}{D_{\mathrm{KL}}}
\newcommand{\Var}{\mathrm{Var}}
\newcommand{\standarderror}{\mathrm{SE}}
\newcommand{\Cov}{\mathrm{Cov}}
% Wolfram Mathworld says $L^2$ is for function spaces and $\ell^2$ is for vectors
% But then they seem to use $L^2$ for vectors throughout the site, and so does
% wikipedia.
\newcommand{\normlzero}{L^0}
\newcommand{\normlone}{L^1}
\newcommand{\normltwo}{L^2}
\newcommand{\normlp}{L^p}
\newcommand{\normmax}{L^\infty}

\newcommand{\parents}{Pa} % See usage in notation.tex. Chosen to match Daphne's book.

\DeclareMathOperator*{\argmax}{argmax}
\DeclareMathOperator*{\argmin}{argmin}

\DeclareMathOperator{\sign}{sign}
\DeclareMathOperator{\sgn}{sgn}
\DeclareMathOperator{\Tr}{Tr}
\let\ab\allowbreak



% Operators the model notation needs beyond math_commands.tex.
\newcommand{\child}{\operatorname{ch}}
\newcommand{\parmap}{\pi}
\newcommand{\Dir}{\mathrm{Dir}}
\newcommand{\Mult}{\mathrm{Mult}}
\newcommand{\Top}{\mathrm{Top}}
\newcommand{\logit}{\mathrm{logit}}
\newcommand{\old}{\mathrm{old}}
\newcommand{\new}{\mathrm{new}}
\newcommand{\NB}{\mathrm{NB}}
\newcommand{\ZINB}{\mathrm{ZINB}}
\newcommand{\SHD}{\mathrm{SHD}}
\newcommand{\Prec}{\mathrm{Prec}}
\newcommand{\Rec}{\mathrm{Rec}}

\newtheorem{assumption}{Assumption}
\crefname{assumption}{assumption}{assumptions}
\Crefname{assumption}{Assumption}{Assumptions}
\crefrangelabelformat{assumption}{#3#1#4 to #5#2#6}
\let\origtheassumption\theassumption

\begin{document}

\title{
    Modeling Behavior Evolution Across Timescales:\\
    Nested Non-Stationary Dynamic Bayesian Networks
}

\author[1,$*$,$\dagger$]{First Author}
\author[1,2,$*$]{Second Author}
\author[2]{Third Author}

\affil[1]{Affiliation 1}
\affil[2]{Affiliation 2}
\affil[$*$]{Equal contribution}
\affil[$\dagger$]{Correspondence to \url{email@example.com}}

\maketitle

\begin{abstract}
One stream of behavior appears at several temporal resolutions at once, as a
sparse sequence of event starts at a fine resolution and as counts per hour or
per day above it, and the dependencies that govern it drift over time. A
dynamic Bayesian network is defined at one chosen resolution and does not force
a fine-scale prediction to aggregate to the coarse summary of the same
quantity; non-stationary variants keep the single resolution; temporal
abstraction supplies no probabilistic account of how a coarse count divides
into finer ones. We present nested non-stationary dynamic Bayesian networks,
which generate each finer scale by dividing its parent, so that counts and
durations reconcile exactly by construction rather than by a later correction,
and which carry two latent states per scale, a discrete regime for persistent
drift and a continuous factor for co-movement inside one epoch, so that
co-movement is not reported as a directed dependence. Structure and parameters
are estimated by a bidirectional variational structural expectation maximization
procedure whose two passes carry the two components of a scale-\(k\) predictor,
the block-constant one that only a coarser scale identifies and the within-block
one that only scale \(k\) does, and whose local score is conditioned on the
latent posterior. We give the
model, the estimator, the identifiability invariances that any evaluation of it
must respect, and a benchmark screened for identifiability before any metric is
read. In that benchmark the coarsest and the middle scale carry true lagged
edges and the finest carries none, so the finest scale scores a method for
staying silent. Against eight published structure learners the
latent-conditional score reaches an $F_1$ of $0.892$ at the coarsest scale,
where those learners reach $0.021$ to $0.346$, and of $0.690$ at the middle
scale, where they reach $0.109$ to $0.251$, and it is the only method whose
total structural distance improves on reporting nothing at all. Its recall of
true cross-label edges at the middle scale is $0.523$ against at most $0.245$
for any of the eight, which separates a dependence between two labels from the
autocorrelation of one. We also evaluate
against released Bayesian networks with exact conditional probability tables,
tiled into two-slice temporal networks, and we report a methodological finding
that governs how such a comparison must be run. A latent-variable model carries
free quantities that a latent-free learner does not, namely the number of
regimes, the factor dimension and the edge penalty, and declaring them rather
than estimating them makes the comparison a comparison of declarations. Choosing
all three on held-out log evidence, by the same rule and from the training data
alone, raises $F_1$ on five tiled networks from $0.532$ to $0.806$. Most of the
gap to hill climbing that the declared setting appeared to show was therefore a
property of the protocol. What remains is a recall deficit of about eight true
edges spread over three of the five networks, which we report as unexplained
rather than attribute to the model's block structure, an explanation we test
and rule out. On an observed wearable and self-report log with no ground truth,
$51$ participants over eight complete calendar weeks each, a model fitted to
the panel draws synthetic logs holding four of six behavior labels within $6$
percent of their real episode volume at exactly zero reconciliation error,
four times closer than an intercept and a participant level draws them, and an
enumerated grid of twenty-one non-stationarity forms accepts four of the six
rows that carry evidence net of a positive control. That grid does not separate
the fitted graph from the empty one on any row that carries evidence, which we
report as a bound on what fifty-one participants over eight weeks can resolve.
Agreement between changepoints found independently in the two logs is monotone
in the calendar scale for every label, which places the departure at the
within-day shape. On the same log, offering calendar nodes at the scale the
model says they are identified at is worth $2295.64$ nats over withholding them
and $33.90$ nats over offering them at every scale.
\end{abstract}

\section{Synopsis}
\label{sec:synopsis}

A nested non-stationary dynamic Bayesian network is built from a ladder of
temporal resolutions, two maps between adjacent rungs, one latent state and one
lagged graph per rung. The densities, the estimator and the proofs follow in
\cref{sec:model,sec:learning,sec:identifiability}.

\paragraph{The ladder.}
An event log is \(\mathcal D=\{(\ell_i,s_i,d_i)\}\): a label, a start, a
duration. Fix nested partitions of the study horizon indexed by
\(k=1,\dots,K\), finest to coarsest, and read off each epoch the number of
episodes that start in it and the duration they carry,
\begin{align}
 Z^{(k)}=\bigl(X^{(k)},Y^{(k)}\bigr),
 \qquad
 X_t^{(k,\ell)}=\#\{i:\ell_i=\ell,\ s_i\in I_t^{(k)}\},
 \qquad
 Y_t^{(k,\ell)}=\!\!\sum_{i:\,s_i\in I_t^{(k)}}\!\! d_i .
 \label{eq:syn_project}
\end{align}
One log, \(K\) views of it. Write \(\child_k(c)\) for the children of a
coarse epoch \(c\).

\paragraph{Folding and unfolding.}
Two maps move between adjacent views. The \emph{fold} is deterministic and
known; the \emph{unfold} is the model,
\begin{align}
 \underbrace{\bigl(\mathcal A_kz\bigr)_c=\!\!\sum_{t\in\child_k(c)}\!\! z_t}_{\text{fold, fine}\to\text{coarse}},
 \qquad
 \underbrace{Z^{(k)}\sim W_k\bigl(\cdot\mid Z^{(k+1)},U^{(k)},C^{(k)}\bigr)}_{\text{unfold, coarse}\to\text{fine}},
 \qquad
 \operatorname{supp}W_k=\mathcal A_k^{-1}\bigl(Z^{(k+1)}\bigr).
 \label{eq:syn_foldunfold}
\end{align}
The unfold is a distribution over the \emph{fiber}: over all the ways a coarse
total can be dealt out among its children. Generation starts at the top and
walks down,
\begin{align}
 Z^{(K)}\sim E_K\bigl(\cdot\mid U^{(K)},C^{(K)}\bigr),
 \qquad
 Z^{(K-1)}\sim W_{K-1},\ \dots,\ Z^{(1)}\sim W_1 ,
 \label{eq:syn_generate}
\end{align}
where \(U^{(k)}=(S^{(k)},H^{(k)})\) is a per-scale latent state, a discrete
regimen for persistent drift and a continuous factor for same-epoch
co-movement, and \(C^{(k)}\) holds calendar quantities as exogenous nodes.
Because every unfold lands inside the fiber it came from,
\begin{align}
 \mathcal A_k\circ W_k=\mathrm{id}
 \qquad\text{almost surely, hence}\qquad
 \mathcal A_{j-1}\cdots\mathcal A_kZ^{(k)}=Z^{(j)}\ \ \text{for all }k<j .
 \label{eq:syn_coherence}
\end{align}
Every draw, every forecast and every simulation reconciles exactly at every
resolution. Nothing is projected afterwards.

\paragraph{Shadow and detail.}
Each scale carries a lagged graph \(\mathcal G^{(k)}\) and a predictor
\(\eta^{(k)}\) built from lagged parents, latents and calendar. Split that
predictor along the blocks, with \(\Pi_k\) the block average,
\begin{align}
 \eta^{(k)}
 =\underbrace{\Pi_k\eta^{(k)}}_{\text{shadow}}
 \;\oplus\;
 \underbrace{(I-\Pi_k)\eta^{(k)}}_{\text{detail}} .
 \label{eq:syn_split}
\end{align}
The unfold is blind to its own shadow: adding a constant to every child of one
block leaves \(W_k\) unchanged, because the split it induces is unchanged.
Therefore
\begin{center}
\emph{the detail at scale \(k\) is identified at scale \(k\); the shadow at
scale \(k\) is identified only above it.}
\end{center}
This single sentence organizes the rest, and it has a spectral reading: folding
is a low-pass filter, so the resolutions partition frequency and each one owns
its own band. Periodic structure is then carried by a few Fourier harmonics,
each landing in the band that contains its frequency, and localized structure by
Haar coefficients on the block tree, where a level shift is simply a large one.
A changepoint therefore has a home scale exactly as a calendar quantity does.
Every parameter has a resolution that can see it. A latent regimen that does not switch inside a block is pure shadow
there, so below the coarsest scale a regimen must act on a feature that varies
within the block. A calendar quantity is pure shadow above the scale at which
its period fits, which is what gives each one a home: hour of day at the hour,
weekday at the day, a study trend at the week. And two adjacent scales share
exactly their shadow, which is the only thing a cross-resolution prior may tie.

\paragraph{The nested walk.}
Learning is one walk up the ladder and one walk down, repeated. The ascent
carries evidence, the descent carries context,
\begin{align}
 \text{ascent }(1\!\to\!K):\ \text{fold the details into sufficient statistics};
 \qquad
 \text{descent }(K\!\to\!1):\ \text{pass shadows down as priors} .
 \label{eq:syn_walk}
\end{align}
Both directions are necessary, and \cref{eq:syn_split} says why: neither scale
can identify what the other owns. One sweep updates the latent posterior
\(q(U)\), then the parameters, then proposes graph moves, each step
non-decreasing on one penalized bound
\begin{align}
 \mathcal J
 =\mathbb E_q\bigl[\log p(Z,U,\mathcal G,\Theta\mid C)\bigr]
 +\mathbb H[q]
 -\sum_k\Bigl[\lambda_k^{\mathrm{rel}}\!\!\sum_{r}z_{r,k}
 +\lambda_k^{\mathrm{rep}}\!\!\sum_{r}z_{r,k}\bigl(|\mathcal R_k(r)|-1\bigr)\Bigr] .
 \label{eq:syn_objective}
\end{align}
A scientific relation \(r\) is bought once at price
\(\lambda_k^{\mathrm{rel}}\); its representations, the choice of lag, quantity
and basis, are cheap refinements inside it. Pricing each representation
separately is what makes a simple relation lose to its own alternatives.

\paragraph{The cascade.}
The descent carries one more thing, and it is what makes the search
affordable. A relation whose lag is shorter than a coarse epoch cannot appear
there as a lag; it appears as same-epoch co-movement. Coarse co-movement is
therefore the shadow of a fine lagged edge, and it orders where to look,
\begin{align}
 \rho_{a\to b}^{(k)}
 =\max_{j>k}\ \bigl|\operatorname{cor}\bigl(Z^{(j,a)},Z^{(j,b)}\bigr)\bigr|
 \quad\longrightarrow\quad
 \text{visit candidates at scale }k\text{ in decreasing }\rho .
 \label{eq:syn_cascade}
\end{align}
This is an ordering and never a filter. Nothing is removed, so at full depth
the accepted set is exactly the unscreened one and the traversal depth is a
speed control with no effect on what is reachable. Coarse regimens order the
search the same way, by naming the epochs in which a relation is worth
examining at all.

\paragraph{What comes back.}
A fit returns, at every resolution, a graph, its conditional distributions, and
the latent regimen and factor paths. Running \cref{eq:syn_generate} forward
from a filtered state produces a forecast; running it without conditioning
produces a synthetic event log. Both are coherent by \cref{eq:syn_coherence},
so a claim made at one resolution is automatically a claim at all of them, and
a synthetic log can be re-projected through \cref{eq:syn_project} and compared
with the real one scale by scale. \Cref{sec:results:arm1} runs that comparison
on an observed log, and the pattern it finds is the one this synopsis predicts:
what a fit reproduces is coarse-scale drift, and where the two logs separate is
inside the day.

\section{Introduction}
\label{sec:introduction}

A behavioral event log is a sparse record of when episodes of labeled activities
start and how long they last. The same log admits several readings at once. Read
at a one-minute resolution it is a binary indicator per label per slot, marking
the slots at which an episode begins. Read at ten minutes it is a count of
episodes together with a total occupied duration. Read at an hour or a day it is
a coarse summary that carries the trend and almost none of the detail. The
readings are not three datasets. They are one process observed through three
windows, and every quantity at a coarse window is a sum of quantities at a fine
one.

Three properties of such logs hold simultaneously and each of them is awkward for
a model defined at a single resolution. The log is sparse, because most slots at
the finest resolution contain nothing. It is organized at several granularities
at the same time, because the questions asked of it range from which minute an
episode began to whether a month differs from the month before. And the
dependencies among labels drift, because a routine that holds in one season is
not the routine of the next, so a graph fitted to the whole horizon describes an
average of several regimes and none of them.

Dynamic Bayesian networks are the natural starting point. They represent a
temporal stochastic process in factored form, admit latent variables, and keep
the temporal dependence interpretable \citep{dean1989model, murphy2002dbn}.
Learning their structure from data has been studied since structure search was
carried over to the dynamic case, hidden variables and incomplete data included
\citep{friedman1998learning}, and the score-based, constraint-based and hybrid
families developed for static networks all transfer to the lagged case
\citep{heckerman1995learning, spirtes2000causation, chickering2002optimal,
tsamardinos2006max}. What the standard formulation does not supply is a
commitment about how the resolutions relate to each other. A network fitted at
ten minutes and a network fitted at an hour are two models, and nothing in either
of them forces the first to sum to the second.

\subsection{Four requirements on a single model}
\label{sec:intro_requirements}

This work takes the position that the three readings should be one model, and
that such a model has to meet four requirements at the same time. Meeting three
of them is comparatively easy; the difficulty is that each of the four constrains
how the others can be met.

\paragraph{Exact aggregation between a fine count and its coarse summary.}
Let the resolutions be indexed by \(k \in \{1,\dots,K\}\), with \(k=1\) the
finest, let \(\gT^{(k)}\) be the ordered set of epochs at scale \(k\) and
\(T_k\) their number, and let the scales be nested through a parent map
\(\parmap_k\) with child sets \(\child_k(c) = \{t : \parmap_k(t) = c\}\). Writing
\(X_t^{(k,\ell)}\) for the count of label \(\ell \in \gL\) in epoch \(t\) at
scale \(k\) and \(Y_t^{(k,\ell)}\) for the total duration of that label in that
epoch, measured in finest-scale slots, the requirement is
\begin{align}
\sum_{t \in \child_k(c)} X_t^{(k,\ell)} = X_c^{(k+1,\ell)},
\qquad
\sum_{t \in \child_k(c)} Y_t^{(k,\ell)} = Y_c^{(k+1,\ell)},
\label{eq:intro_reconciliation}
\end{align}
for every \(k<K\), every parent epoch \(c \in \gT^{(k+1)}\) and every label. The
identity is a property of the data, since a coarse count is literally the sum of
the fine counts inside it. A model that predicts the two scales independently
will violate it, and repairing the violation after the fact is a separate
estimation problem. The n\(^2\)DBN removes the problem instead of solving it. The
coarsest scale is generated freely, and every finer scale is generated by
allocating the parent total among its children through a distribution whose
support is the set of non-negative integer vectors summing to that total, so
\cref{eq:intro_reconciliation} holds almost surely rather than approximately. In
the reported runs the measured reconciliation error is exactly zero for counts
and for durations in every fit of every method, including the fit on an observed
log.

\paragraph{Learned rather than assumed lagged dependencies.}
Which labels predict which, and at which lag, is the scientific content of the
graph, and fixing it by hand decides in advance what the analysis can find. The
model therefore carries a lagged parent set \(\gP_{k,\ell} \subseteq \gL \times
\{1,\dots,p_k\}\) per label and per scale, and learns it by maximizing a local
score penalized by \(\lambda_k |\gP_{k,\ell}|\) inside a bounded candidate set.
Because every admissible parent is lagged, any graph over the unrolled variables
is acyclic, so the search needs no acyclicity constraint and the score decomposes
over targets.

\paragraph{Separate representation of persistent drift and same-epoch
co-movement.}
Two labels can move together for two quite different reasons. One is a persistent
change in the process itself, a regime that raises or lowers the rate of several
labels for as long as it lasts. The other is a contemporaneous coupling within a
single epoch that carries no direction. Neither is a lagged dependence, and a
model that offers only lagged edges will report both of them as edges. The
n\(^2\)DBN gives each its own latent variable. A regime \(S_t^{(k)} \in
\{1,\dots,M_k\}\) follows a context-dependent Markov chain and carries persistent
drift, and a factor \(H_t^{(k)} \in \sR^{r_k}\) carries same-epoch co-movement.
No observed variable is a parent of another observed variable in the same epoch.
Contemporaneous association is mediated entirely through \(S_t^{(k)}\) and
\(H_t^{(k)}\), which is what keeps co-movement out of the reported graph and, at
the same time, keeps the unrolled graph acyclic.

\paragraph{A mapping back to an event trace with wall-clock times.}
A multiscale summary that cannot be turned back into events answers a narrower
question than the one that was asked. At the finest scale the model constrains
counts to be binary, so \(X_t^{(1,\ell)} = 1\) means an episode of label
\(\ell\) begins in slot \(t\), and \(Y_t^{(1,\ell)}\) is that episode's length in
slots. With a slot width \(\Delta_1\) and a time origin \(t_0\), the episode
beginning at slot \(t\) occupies the wall-clock interval
\begin{align}
\left[\,
t_0 + (t-1)\,\Delta_1,\;
t_0 + \bigl(t - 1 + Y_t^{(1,\ell)}\bigr)\,\Delta_1
\,\right),
\label{eq:intro_episode}
\end{align}
and the trace of label \(\ell\) is the sequence of these intervals in increasing
order of start slot. The map is deterministic, so a coarse-scale prediction
propagates to a dated event list without a second estimation step.

\subsection{Contributions}
\label{sec:intro_contributions}

The paper makes five contributions.

First, it formulates the nested non-stationary dynamic Bayesian network, a single
probabilistic model of event counts and durations at nested resolutions in which
\cref{eq:intro_reconciliation} holds by construction, lagged structure is learned
per scale, persistent drift and same-epoch co-movement are carried by separate
latent variables, and \cref{eq:intro_episode} returns a dated event trace.

Second, it derives a bidirectional variational Structural EM procedure for the
model, in which the regime posterior is obtained by forward-backward recursion,
the factor posterior by a Gaussian approximation, the parameters by monotone
block updates, and the graph by a penalized local search whose score is an
expectation under the current posterior over the latent variables and includes
the block normalizer of the allocation. One sweep is two passes, fine to coarse
and then coarse to fine, and each pass identifies the component of a scale-$k$
predictor the other leaves unidentified.

Third, it measures what that latent-conditional score buys, by placing the
structure step behind a common interface with published structure learners and
running them on the same generated data, inside the same in-degree budget, with
the same parameters fitted afterwards, so that the only difference between two
rows of a results table is the graph. The benchmark is screened beforehand, so
that a low score is attributable to an estimator rather than to a scenario
containing nothing to find. The screen also fixes how many true parents each
scale carries. It places lagged edges at the coarsest and at the middle scale
and none at the finest, where a lagged edge on a binary start process can only
move which slot of a block begins an episode and does not pay for itself, so
the finest scale grades a method on staying silent.

Fourth, it checks the same code against external ground truth, using released
Bayesian networks tiled into two-slice temporal networks, and establishes what
such a check requires. A latent-variable model carries free quantities that a
latent-free learner does not, namely the number of regimes, the factor
dimension and the edge penalty. Fixing them by hand and then reporting a
comparison measures the choices, not the criteria. We therefore select the
regime count and the edge penalty on held-out log evidence from the training
split alone, holding the factor dimension at zero because the tiled networks
carry no same-slice factor, and we show that on those networks the selection
recovers the setting that maximizes $F_1$ without ever reading the true graph.
On the generated grid the same rule picks a different edge penalty from the one
that maximizes $F_1$, and rescoring the frontier at every scale that carries a
true edge shows this is a property of the selector rather than of the metric:
the middle scale, the coarsest scale and the mean over scales each disagree with
held-out evidence, and the two truth-bearing scales disagree with each other, so
no single penalty is best at both.

Fifth, it documents the failures. A residual leak of same-epoch structure into
finest-scale lagged self edges survives every mechanism excluded for it; the
mean-field evaluation of the allocation normalizer is contained by an outer
guard rather than removed; recall on external data remains below that of hill
climbing after selection; and parameter recovery does not meet the sample-size
criterion stated before the experiment was run, from a ground-truth start as
well as from a data-driven one.

\section{Related work and the benchmark landscape}
\label{sec:related}

\subsection{Dynamic Bayesian networks and structure learning}
\label{sec:related_dbn}

Dynamic Bayesian networks generalize hidden Markov models and linear state-space
models by writing a temporal process as a factored graphical model over a
repeated time slice \citep{dean1989model, murphy2002dbn}. The representation is
attractive precisely when the temporal dependence has to remain readable, since
each edge names a variable, a lag and a target. Structure learning for the
dynamic case inherits the machinery of the static case. Score-based search
maximizes a decomposable criterion, most often the Bayesian information criterion
or a Bayesian Dirichlet score \citep{heckerman1995learning}; constraint-based
search recovers a skeleton from conditional independence tests
\citep{spirtes2000causation}; greedy equivalence search operates on equivalence
classes \citep{chickering2002optimal}; and hybrid methods combine a local
neighborhood step with a scored search \citep{tsamardinos2006max}. The tree
special case has an exact maximizer \citep{chow1968approximating}, the
time-series literature contributes lag selection read as model selection
\citep{granger1969investigating}, and continuous-optimization formulations
replace the combinatorial search with a smooth acyclicity penalty and an
autoregressive design \citep{pamfil2020dynotears}. Open implementations of most
of these are available in a single library \citep{ankan2015pgmpy, ankan2024pgmpy}.

Two features of that literature matter here. Every one of these criteria is
computed from the observed variables alone, so none of them conditions on a
latent state while scoring a candidate parent. And the flagship survey of
Bayesian network structure learning covers 74 algorithms and states explicitly
that structure learning from time-series data is outside its scope, deferring to
a separate review \citep{kitson2023bnsurvey, nogueira2022cdtools}. The dynamic
case is a distinct literature with distinct tooling, and the present work sits
inside it.

\subsection{Non-stationary and time-varying variants}
\label{sec:related_nonstationary}

Classical dynamic Bayesian network learning assumes a stationary generating
process, which long behavioral horizons violate. The response splits into two
premises. In the first, the graph itself changes, and learning means recovering a
sequence of structures together with the times at which they change
\citep{robinson2010learning, lebre2010artiva}. In the second, the graph is shared
and only the parameters change across segments, which is the premise of the
allocation and coupling models of Grzegorczyk and Husmeier and of their
successors \citep{grzegorczyk2011nhdbncontinuous, dondelinger2013edison}. A
parallel line treats the parameters as varying smoothly rather than in segments,
estimating a time-varying network by kernel-weighted or penalized regression
\citep{song2009tvdbn, kolar2010estimating}. More recent work makes the regime an
explicit latent variable to be inferred jointly with per-regime structure
\citep{saggioro2020rpcmci, rahmani2025castor}.

The n\(^2\)DBN belongs to the second premise. The lagged graph at a scale is
shared across the horizon and the regime \(S_t^{(k)}\) modulates the emission and
allocation parameters, with the regime chain at a finer scale conditioned on its
parent-scale regime so that drift is coherent across resolutions. What separates
this model from the cited work is that the regime is one of two latent
mechanisms rather than the only one, and that the scales are coupled by an exact
aggregation identity rather than by a shared prior.

There is a specific hazard in this setting that the non-stationary literature has
already named. A persistent latent state raises the level of every variable it
drives for as long as it lasts, so each variable becomes positively
autocorrelated even with an empty parent set, and a criterion that conditions on
nothing attributes that autocorrelation to a lagged self edge
\citep{runge2018pitfalls}. This is not a defect of any particular algorithm. It
is a property of the hypothesis class those algorithms search, and it is the
mechanism the comparison in this paper measures.

\subsection{Temporal abstraction and hierarchical time series}
\label{sec:related_abstraction}

Reasoning at several temporal granularities rather than only over instantaneous
slices is an old concern in the Bayesian network literature, where temporal
abstraction supplies the vocabulary of intervals and episodes that a slice-based
model lacks \citep{burns2003temporal}. The forecasting literature supplies the
other half of the picture. A hierarchical time series carries an exact and known
summing map from the bottom level to the aggregates, and independent forecasts of
the levels are almost never coherent with it, so a reconciliation step projects
them back onto the space that satisfies the summing constraints
\citep{athanasopoulos2009tourism, wickramasuriya2019mint}. The summing map of
that literature is the direct analogue of \cref{eq:intro_reconciliation}.

The n\(^2\)DBN takes the constraint from the second literature and the episode
semantics from the first, and differs from both in where the constraint is
imposed. Reconciliation in hierarchical forecasting is a projection applied after
independent models have been fitted, and it is optimal under an assumption about
the base forecast errors. Here the constraint is a property of the generative
model, since a finer scale exists only as an allocation of its parent total, so
there is no incoherent intermediate object to project and no reconciliation
estimator to specify. The cost of that choice is that the fine scale is not free.
It competes for a fixed total, which turns out to have consequences the
limitations record.

\subsection{The benchmark landscape}
\label{sec:related_benchmarks}

A model whose central property is exact cross-scale aggregation invites an
obvious question. Why not score it on a public benchmark rather than on a
generator written for the purpose? We surveyed the datasets available for scoring
a dynamic Bayesian network learner and found a specific and citable gap.

Restricting attention to resources that carry both a released ground-truth graph
and a genuine time series, which is the set on which a structure learner can be
scored at all, leaves twelve. They are the DREAM4 and DREAM3 in-silico network
challenges \citep{marbach2010revealing, prill2010dream3} and the GeneNetWeaver
simulator behind them \citep{schaffter2011gnw}; the TimeGraph synthetic suite
\citep{ferdous2025timegraph}; CausalDynamics \citep{herdeanu2025causaldynamics};
CausalTime \citep{cheng2024causaltime}; the Causal Chambers
\citep{gamella2025causalchambers}; CausalRivers \citep{stein2025causalrivers};
the IRMA yeast circuit \citep{cantone2009irma}; the NetSim simulated fMRI
networks \citep{smith2011netsim}; FinanceCPT \citep{nauta2019tcdf}; and
CIPCaD-Bench \citep{menegozzo2022cipcad}. Two kinds of resource frequently
listed alongside them do not belong, since SynTReN generates steady-state data
rather than a time series, and the pseudotime benchmarks order cells along a
differentiation trajectory over a network that does not change, so what varies
across their samples is not time.

Not one of the twelve combines the three properties this model is about, namely
event counts, exact cross-scale aggregation and a latent regime. The gap is not
that any one of the three is unattainable. A regime change with something to
score against is present in five of the twelve, and count series that aggregate
exactly over a known hierarchy are routine in forecasting repositories. What no
resource does is release a ground-truth graph alongside all three at once.

DREAM4 is the most used quantitative benchmark in the field, with five networks
of ten genes and five of a hundred, replicated time series of twenty-one points
each, and a released directed gold standard; DREAM3 is the same construction at
ten, fifty and a hundred genes. The observations are continuous expression levels
rather than counts, the design is single resolution, and the perturbation applied
and then removed is a stimulus rather than a structural change, so the process is
stationary. GeneNetWeaver is the simulator that produced them and generates the
same kind of data at chosen sizes, with the same three absences.

TimeGraph is the closest synthetic match on the non-stationarity axis, since
non-stationarity is an explicit generator axis alongside trends, seasonality,
irregular sampling and missingness, and it ships the structural equations rather
than only the graph. Its series are four to eight continuous variables over 500
to 3000 steps at a single resolution, so it has a changing generator but no
counts and no aggregation.

The Causal Chambers are the strongest real-data item, and the only one giving a
real physical system, a known graph, known intervention targets and marked
changepoints at once. The regime-jump datasets run 320{,}000 observations in
thirty-two regimes of ten thousand, with an intervention drawn from an
eleven-variable pool at seventy percent of the regime boundaries. The observed
quantities are sensor readings rather than event counts, the regime is an applied
intervention rather than a latent state to be inferred, and there is no nesting,
so any hierarchy would be imposed by the analyst rather than intrinsic.

CausalRivers is the largest real-world time-series causal benchmark, with
hundreds of stations at fifteen-minute resolution over several years and a
released station graph, and its flood product contains a genuine distributional
shift. The shift is real, it is not labeled as a regime path, the measurements
are water levels rather than counts, and nothing aggregates exactly.

IRMA is the only classical real dataset with both a known network and an
experimentally imposed regime switch, since the growth medium is changed inside
the series. It is five genes, eight links and series of sixteen or twenty points,
which supports a demonstration and not a power analysis, and its readings are
again continuous and single-scale.

CausalDynamics is the largest synthetic resource, with more than fourteen
thousand graphs and more than fifty million samples drawn from coupled
differential equations and two idealized climate models, and its adjacency is
derived analytically from the governing equations. The trajectories are chaotic
or stochastic rather than regime-switching by design, and they are continuous and
single-scale.

The remaining four complete the picture without changing it. NetSim and
FinanceCPT are stationary simulations on fixed graphs, the second usefully
labeling its ground truth with the lag. CausalTime fits a deep model to real data
and emits synthetic series that keep real dynamics and carry a graph by
construction, which is the right instrument against the charge that fixed
synthetic benchmarks are too easy, but the graph is generated rather than
externally validated and none of the three properties is present. CIPCaD-Bench
supplies an industrial process with faults injected mid-run and a graph derived
from a plant flowsheet, which is the closest thing to a regime among the
industrial data, again without counts or nesting.

The tempting alternative is a static repository. The bnlearn Bayesian Network
Repository distributes fully parameterized networks with an exact directed graph
and exact conditional probability tables \citep{scutari2010bnlearn}, and a
package vendors them so that an experiment runs without network access
\citep{ankan2024pgmpy}. It has no temporal section of any kind. Turning one of
its networks into a two-slice temporal network means mapping each edge onto a
lag-1 edge and sampling, which is what the field already does when it needs a
dynamic benchmark, with published methodology and published precedent
\citep{trabelsi2013benchmarking, ouyang2023divideconquer}. The construction is
sound and this work uses it, but it manufactures the time series it needs. The
tiled network has no nesting, no counts that aggregate and no latent regime, so
three of the four requirements of \cref{sec:intro_requirements} are absent by
construction. There is also no dynamic counterpart to the repository itself. The
2025 expansion of the static collection adds a further large set of networks
drawn from the literature and no dynamic ones \citep{leonelli2025bnrep}, and the
reference reproducible
workflow for benchmarking structure learning has no temporal graph generator, no
time-series data module and no dynamic algorithm entries
\citep{rios2025benchpress}.

The conclusion is narrow and it is the justification for this work shipping its
own benchmark. Multiscale reconciliation cannot be scored on any released
benchmark, because no released benchmark contains it. What can be scored
externally is the structure learner, the local score and the table estimator on
flat discrete data, which is what the tiled networks are used for here, and the
resulting comparison is reported alongside the generated one rather than in place
of it.

\subsection{Released parameters, not only released graphs}
\label{sec:related_parameters}

A recovery claim about parameters needs more than a released graph, since the
generating conditional distributions have to be known exactly for the comparison
to mean anything. Most of the twelve release a graph and nothing else. Four
resources release true parameters. GeneNetWeaver attaches kinetic models of
transcription and translation to the subnetworks it extracts, so a generated
network carries its own parameters by construction. TimeGraph releases the
structural equations and the full generating metadata alongside each graph.
CausalDynamics derives its adjacency from governing equations that are themselves
released. And the tiled bnlearn networks carry exact conditional probability
tables, which is why the external check in this paper can score an estimated
table against a published one rather than only counting edges. The real systems
in the list, namely the Causal Chambers, CausalRivers and IRMA, have no true
parameters to release, and the DREAM challenges release edge lists rather than
the kinetics of the specific challenge networks.

Two further hazards of benchmark design deserve mention, since both would be
raised against any evaluation of this kind. A synthetic generator for a two-slice
network can induce an ordering of the variables by marginal variance that
continuous-optimization learners exploit for the wrong reason, so standardization
and a report of the residual sortability are both needed
\citep{reisach2021varsortability}. And structural Hamming distance is defined
over partially directed graphs in its original form \citep{tsamardinos2006max}
while several widely used libraries compute it on the directed adjacency, so a
paper has to state which convention it used before its distances can be compared
with anyone else's.

Finally, a falsification test that requires no ground truth exists and transfers
to real logs. The self-compatibility check compares graphs learned from different
variable subsets and asks whether they agree where they overlap
\citep{faller2024selfcompatibility}. It is necessary rather than sufficient, and
it is the only structural check available where the true graph is unknown.

\clearpage
\section{Compact formulation}
\label{sec:model}

This section summarizes N\textsuperscript{2}DBN from an observed event log to
the learned multiresolution graphs, conditional parameters, and generated
event logs. The following sections retain derivations and implementation
details. The construction combines DBN semantics
\citep{dean1989model,murphy2002dbn}, multiscale temporal representations
\citep{saria2007granularity,kang2022multiscale}, coherent probabilistic
allocation \citep{das_topdown2023}, and coupling between related networks
\citep{grzegorczyk2012globalcoupling,kamalabad2020edgewise}.

\paragraph{Event log and temporal granularity.}
\label{sec:model_scales}

An event log is
\(\mathcal D=\{(\ell_i,s_i,d_i)\}_{i=1}^N\), where \(\ell_i\in\mathcal L\)
is an event label, \(s_i\) is its start time, and \(d_i\in\mathbb N\) is its
duration in units of the finest epoch. Let \(k=1,\ldots,K\) index nested
temporal partitions from finest to coarsest, with epochs
\(\mathcal T^{(k)}=\{I_t^{(k)}\}\). Counts and total durations at each
granularity are
\begin{align}
 X_t^{(k,\ell)}
 &=\sum_{i=1}^N\mathbb{I}\{\ell_i=\ell,\ s_i\in I_t^{(k)}\},
 &Y_t^{(k,\ell)}
 &=\sum_{i=1}^Nd_i\mathbb{I}\{\ell_i=\ell,\ s_i\in I_t^{(k)}\}.
 \label{eq:model_log_projection}
\end{align}
Write \(Z^{(k)}=(X^{(k)},Y^{(k)})\). The map \(\pi_k\) assigns an epoch to
its containing epoch at scale \(k+1\). It induces the aggregation map, its
fiber, and event-log reconstruction at the finest granularity,
\begin{align}
 (\mathcal A_kZ^{(k)})_c
 &=\sum_{t:\pi_k(t)=c}Z_t^{(k)},
 &Z^{(k+1)}&=\mathcal A_kZ^{(k)},
 &\mathcal F_k(v)&=\{u:\mathcal A_ku=v\},
 \label{eq:model_aggregation}\\
 \mathcal E(Z^{(1)})
 &=\{(\ell,t,t+Y_t^{(1,\ell)}-1):X_t^{(1,\ell)}=1\}.
 \label{eq:model_episode}
\end{align}
Counts obey the capacity of their epochs, durations obey \(Y\geq X\), and
\(X_t^{(1,\ell)}\in\{0,1\}\).

Let \(\mathcal Q=\{\mathrm{cnt},\mathrm{dur}\}\). A learned edge is
\(e=((a,q_s),(b,q_t),d)\), with \(1\leq d\leq p_k\). For its indicator
\(w_{e,k}\), the lagged values available to target \((b,q)\) are
\begin{align}
 \mathcal P_{k,b,q}(t)
 &=\{Z_{t-d,q'}^{(k,a)}:w_{((a,q'),(b,q),d),k}=1\}.
 \label{eq:model_lag_values}
\end{align}
The learned event structure is
\(\mathcal G=\{\mathcal G^{(k)}\}_{k=1}^K\). Its edges have positive lag, so
every unrolled graph is acyclic. Aggregation and latent links are fixed parts
of the model construction and are excluded from \(\mathcal G\).

Let \(U_t^{(k)}=(S_t^{(k)},H_t^{(k)})\) combine a discrete regimen for
persistent change and a continuous factor for same-epoch co-movement. Context
is denoted by \(C_t^{(k)}\). Resolution-specific maps define the observation
parameters and latent dynamics,
\begin{align}
 \eta_{t,b,q}^{(k)}
 &=f_{k,b,q}(\mathcal P_{k,b,q}(t),U_t^{(k)},C_t^{(k)};\Theta^{(k)}),
 \qquad q\in\mathcal Q,
 \label{eq:model_count_score}\\
 U_t^{(k)}
 &\sim R_k(U_{t-1}^{(k)},U_{\pi_k(t)}^{(k+1)},C_t^{(k)};\Theta^{(k)}).
 \label{eq:model_regime_init}
\end{align}
Boundary arguments are omitted at \(t=1\) and \(k=K\). At the coarsest scale,
an emission map \(E_K\) specifies the count and duration-surplus families. At
finer scales, a normalized allocation map \(W_k\) defines a child kernel,
\begin{align}
 P_K(Z^{(K)}\mid U^{(K)},C^{(K)})
 &=E_K(\boldsymbol\eta^{(K)};\Theta^{(K)}),
 \label{eq:model_coarse_count_mean}\\
 K_k(Z^{(k)}\mid Z^{(k+1)},U^{(k)},C^{(k)})
 &=W_k(\boldsymbol\eta^{(k)},Z^{(k+1)};\Theta^{(k)}),
 \qquad \operatorname{supp}K_k=\mathcal F_k(Z^{(k+1)}).
 \label{eq:model_count_omega}
\end{align}
The implemented \(E_K\) uses zero-inflated negative binomial counts and a
negative binomial duration surplus. The map \(W_k\) uses constrained
Dirichlet-Multinomial allocations for counts and duration surplus. It includes
capacity, sequential depletion, the duration floor, the active-child set, and
the block normalizer.

Let \(\Psi_K\) contain the coarsest latent dynamics and emissions, and let
\(\Psi_k\), \(k<K\), contain the child dynamics and kernel. The joint model is
\begin{align}
 p(\mathcal G,\Theta,Z^{(1:K)},U^{(1:K)}\mid C^{(1:K)})
 &=p(\mathcal G)p(\Theta\mid\mathcal G)\,\Psi_K
   \prod_{k=1}^{K-1}\Psi_k.
 \label{eq:model_joint}
\end{align}
The fiber support gives exact reconciliation for every draw.

\begin{proposition}[Exact multiresolution coherence]
\label{prop:model_reconciliation}
For every parameter value and every draw of the model,
\begin{align}
 \sum_{t:\pi_k(t)=c}X_t^{(k,\ell)}&=X_c^{(k+1,\ell)},
 &\sum_{t:\pi_k(t)=c}Y_t^{(k,\ell)}&=Y_c^{(k+1,\ell)}
 \quad\text{almost surely},
 \label{eq:model_reconciliation}
\end{align}
and consequently \(\mathcal A_{j-1}\cdots\mathcal A_kZ^{(k)}=Z^{(j)}\) for
every \(k<j\).
\end{proposition}

\begin{proof}
\Cref{eq:model_count_omega} gives \(\operatorname{supp}K_k=\mathcal
F_k(Z^{(k+1)})\), and \(u\in\mathcal F_k(v)\) means \(\mathcal A_ku=v\) by
\cref{eq:model_aggregation}, which is the one-step claim. For \(j>k+1\),
suppose \(\mathcal A_{j-2}\cdots\mathcal A_kZ^{(k)}=Z^{(j-1)}\); applying
\(\mathcal A_{j-1}\) and using the one-step case at \(j-1\) gives the claim, so
the identity composes by induction on \(j-k\).
\end{proof}

\Cref{prop:model_reconciliation} is a statement about the arrays \((X,Y)\). The
reconstruction \(\mathcal E\) of \cref{eq:model_episode} additionally imposes a
one-slot floor, a horizon clip and, where labels are declared mutually
exclusive, an occupancy projection. \Cref{sec:identifiability} states each and
bounds the trace-level discrepancy they introduce; the arrays are exact.

\paragraph{What each resolution can identify.}
\label{sec:model_blockconstant}

The kernel \(W_k\) is normalized over a fiber, and that normalization removes a
subspace of its own predictor from the likelihood at scale \(k\). Fix a parent
epoch \(c\in\mathcal T^{(k+1)}\), write \(\child_k(c)=\{t:\pi_k(t)=c\}\), and
decompose the coordinates on that block orthogonally,
\begin{align}
 \sR^{\child_k(c)}
 &=\mathcal N_c^{(k)}\oplus\mathcal V_c^{(k)},
 &\mathcal N_c^{(k)}&=\operatorname{span}\{\mathbf 1_{\child_k(c)}\},
 &\mathcal V_c^{(k)}&=\Bigl\{v:\textstyle\sum_{t\in\child_k(c)}v_t=0\Bigr\},
 \label{eq:model_block_split}
\end{align}
with \(\Pi_k\) the block-averaging projector onto the block-constant
directions,
\((\Pi_kv)_t=|\child_k(\pi_k(t))|^{-1}\sum_{s\in\child_k(\pi_k(t))}v_s\).

\begin{proposition}[Block-constant invariance]
\label{prop:model_block_invariance}
For every parent epoch \(c\), every \(a\in\sR\) and every predictor
\(\boldsymbol\eta^{(k)}\),
\begin{align}
 W_k(\boldsymbol\eta^{(k)}+a\mathbf 1_{\child_k(c)},Z^{(k+1)})
 &=W_k(\boldsymbol\eta^{(k)},Z^{(k+1)}).
 \label{eq:model_block_invariance}
\end{align}
Hence \(\Pi_k\boldsymbol\eta^{(k)}\) is unidentified from \(\Psi_k\), and only
\((I-\Pi_k)\boldsymbol\eta^{(k)}\) is.
\end{proposition}

\begin{proof}
Within block \(c\) the allocation concentrations are
\(\alpha_t=\kappa_k\,\sigma_t(\boldsymbol\eta^{(k)})\) with \(\sigma\) the
softmax over \(\child_k(c)\). The softmax is invariant to a constant added to
every coordinate of the block, so \(\alpha\) is unchanged, and the constrained
Dirichlet-Multinomial density on \(\mathcal F_k(Z^{(k+1)})\) depends on
\(\boldsymbol\eta^{(k)}\) only through \(\alpha\).
\end{proof}

Three consequences organize the model.

First, \(\Pi_k\boldsymbol\eta^{(k)}\) is not lost. It acts on the data through
the parent total \(Z^{(k+1)}\), so it is identified from
\(\Psi_{k+1},\ldots,\Psi_K\). Every parameter has a resolution at which it is
identified, and the two directions of the learning pass carry exactly these two
components: the coarse-to-fine message supplies \(\Pi_k\boldsymbol\eta^{(k)}\)
and the fine-to-coarse message supplies the sufficient statistics for
\((I-\Pi_k)\boldsymbol\eta^{(k)}\). This is the sense in which the pass is
bidirectional.

Second, a term constant across a block contributes exactly zero at scale \(k\).
A latent regimen entering additively is block-constant on every block over
which its path does not switch, so at \(k<K\) it must act on a feature carrying
within-block variation,
\begin{align}
 \eta_{t,b,q}^{(k)}
 &\supseteq\beta_{S_t^{(k)},b,q}^{(k)}\,g_{t,b,q}^{(k)},
 \qquad (I-\Pi_k)g^{(k)}\neq0,
 \label{eq:model_regime_interaction}
\end{align}
with \(g\) prefix-measurable: the depletion feature, a lagged parent value, or
the recency column below. An additive regimen level is admitted at \(k=K\)
alone.

Third, the cross-resolution prior acts on the block-constant component, which
is the component two adjacent scales share.

\paragraph{A covariate that is never block-constant.}

The lag axis is the largest source of competing representations of one
relation, and one prefix-measurable summary of it removes most of that
competition. Let \(\delta_t^{(k,\ell)}=t-\max\{s<t:X_s^{(k,\ell)}>0\}\) be the
recency of label \(\ell\) and \(M_t^{(k,\ell)}\) its geometric-lag counterpart,
\begin{align}
 M_t^{(k,\ell)}(\rho)
 &=(1-\rho)\sum_{d\geq1}\rho^{\,d-1}Z_{t-d}^{(k,\ell)},
 \qquad\text{columns }
 \log\bigl(1+\delta_t^{(k,\ell)}\bigr),\;M_t^{(k,\ell)}.
 \label{eq:model_recency}
\end{align}
Both are measurable with respect to the strict past, so they preserve
acyclicity, and neither is block-constant at an interior scale, so both stay
identified under \cref{prop:model_block_invariance}.

\paragraph{The ladder is a band-pass decomposition.}
\label{sec:model_bands}

Summing over the children of a block is a low-pass filter, so
\cref{prop:model_block_invariance} has a spectral reading. Let \(n_k\) be the
number of finest slots in one scale-\(k\) epoch. A signal supported on
frequencies above \(1/n_{k+1}\) has zero average on every scale-\(k+1\) block
and is therefore pure detail at scale \(k\); a signal supported below
\(1/n_k\) is constant across each scale-\(k\) block and is pure shadow there.
The detail at scale \(k\) is the band
\begin{align}
 \mathcal B_k=\Bigl[\tfrac{1}{n_{k+1}},\ \tfrac{1}{n_k}\Bigr),
 \qquad
 \bigoplus_{k=1}^{K}\mathcal B_k=\text{the whole spectrum},
 \label{eq:model_bands}
\end{align}
so the resolutions partition frequency, and each one estimates the part of the
signal that lives in its own band. Two bases realize this, and the structure decides
which one applies.

\emph{Periodic structure takes a Fourier basis.} A calendar effect of period
\(P\) enters as a few harmonics,
\begin{align}
 \gamma_{P}(t)=\sum_{h=1}^{H}\Bigl[a_h\cos\tfrac{2\pi ht}{P}+b_h\sin\tfrac{2\pi ht}{P}\Bigr],
 \label{eq:model_harmonic}
\end{align}
which costs \(2H\) parameters where a free profile costs \(P-1\), and which is
smooth by construction. Harmonic \(h\) has frequency \(h/P\) and is therefore
identified at the unique scale whose band contains it. \Cref{eq:model_harmonic}
with \(H\) at its maximum recovers the free profile exactly, so the harmonic
order is a smoothness dial and not a modelling commitment. This matters most
where epochs are few: a week-scale effect estimated from nine epochs per
participant cannot afford a free profile and can afford two harmonics.

\emph{Localized structure takes a Haar basis.} A level shift is not periodic,
and a Fourier expansion represents it badly. In the block decomposition of
\cref{eq:model_block_split} the detail coefficients are Haar coefficients on
the block tree, and a level shift at scale \(k\) is exactly a large one,
\begin{align}
 \delta_t^{(k)}=\bigl\langle\eta^{(k)},\psi_t^{(k)}\bigr\rangle,
 \qquad
 \psi_t^{(k)}\in\mathcal V_{\pi_k(t)}^{(k)} .
 \label{eq:model_haar}
\end{align}
So a changepoint has a home scale in the same sense a calendar node does: a
shift that persists longer than one scale-\(k\) block appears as shadow there
and as detail above, and one that decays inside a block appears as detail
there and is invisible above. The latent regimen \(S^{(k)}\) is the
probabilistic version of the same object, with the switch times its
changepoints, and reading them off the posterior at each scale returns the
shifts the ladder can resolve.

\paragraph{Calendar quantities are nodes, and each has a home scale.}
\label{sec:model_calendar}

Hour of day, weekday, day index and week index are vertices of
\(\mathcal G\). Each is a known function of the clock
and has no parents, so an edge from one into a label may carry lag zero without
threatening acyclicity, and the conditional
\(P_{\mathrm{day}}(\text{walk}\mid\text{weekday}=0)\) becomes a table the model
reports. Whether such a node earns its place is decided by the same price as
any other candidate.

\Cref{prop:model_block_invariance} says where each can be earned. A calendar
node \(A\) is a deterministic function of the epoch, so it enters \(\Psi_k\)
only through \((I-\Pi_k)A\); it is therefore identified at scale \(k\) exactly
when it varies inside a scale-\(k\) parent block, and contributes nothing at any
coarser scale where it does not.

\begin{corollary}[Home scale of a calendar node]
\label{cor:model_calendar_home}
Let \(A\) have period \(P_A\) in finest slots and let \(n_k\) be the number of
finest slots in one scale-\(k\) epoch. Then \(A\) is identified at scale \(k\)
if and only if \(n_{k+1}>P_A\) is false and \(n_k<P_A\); equivalently, its home
scale is the finest \(k\) whose parent epoch spans at least one full period of
\(A\).
\end{corollary}

So hour of day is identified where a day is the parent, weekday where a week is
the parent, and a trend index wherever the horizon is the parent. Placing a
calendar node above its home scale buys nothing, and the price it would pay
there is wasted. This turns a modelling convention into a consequence of the
kernel.

For a typed relation \(r\), let \(z_{r,k}\) indicate that some lag is active.
Adjacent structures and standardized active effects evolve according to
\begin{align}
 \Pr(z_{r,k}=1\mid z_{r,k+1}=1)&=\pi_{r,k}^{\mathrm{keep}},
 &\Pr(z_{r,k}=1\mid z_{r,k+1}=0)&=\pi_{r,k}^{\mathrm{appear}},
 \label{eq:model_scale_coupling}\\
 \Pi_k\alpha_{r,k}\mid\alpha_{r,k+1}
 &\sim\mathcal N(\mathcal B_{k+1\to k}(\alpha_{r,k+1}),\tau_{r,k}^2I),
 &(I-\Pi_k)\alpha_{r,k}
 &\sim\mathcal N(0,\tau_{r,k,\perp}^2I).
 \label{eq:model_parameter_evolution}
\end{align}
The tie is on the block-constant component alone, because by
\cref{prop:model_block_invariance} that is the component the parent scale
identifies and the child scale does not. The within-block component has no
parent counterpart and carries its own prior.

The broadcast \(\mathcal B_{k+1\to k}\) is the adjoint of block averaging,
rescaled to the child's feature units,
\begin{align}
 \bigl(\mathcal B_{k+1\to k}\alpha\bigr)_{t}
 &=\frac{\operatorname{sd}(F_{r,k})}{\operatorname{sd}(F_{r,k+1})}\,
 \alpha_{\pi_k(t)}\,\frac{d_{k+1}}{d_k},
 \qquad
 \mathcal B_{k+1\to k}=\Pi_k\mathcal B_{k+1\to k},
 \label{eq:model_broadcast}
\end{align}
where \(d_k\) is the lag at which \(r\) is active at scale \(k\), so the ratio
\(d_{k+1}/d_k\) converts a parent-scale lag into the child's clock. Relations
and lag sets may change by resolution; when \(z_{r,k+1}=0\) the first line of
\cref{eq:model_parameter_evolution} is dropped and \(\alpha_{r,k}\) carries the
second prior alone. A Bernoulli gate on the tie, with a Beta-conjugate update
for \(\pi_{r,k}^{\mathrm{keep}}\) and \(\pi_{r,k}^{\mathrm{appear}}\), lets an
abrupt resolution change escape the coupling; \(\tau_{r,k}\) and the gate are
estimated, and \cref{sec:learning} gives their updates.

\paragraph{One price per relation.}

Scale, lag, quantity pairing and basis are alternative representations of one
scientific relation, and pricing each separately makes a simple relation
compete against its own alternatives. Let \(r=(a,b)\) be an ordered label pair
and let \(g\in\mathcal R(r)\) index a representation of it. The structural
penalty in \cref{eq:model_elbo} is nested,
\begin{align}
 \sum_k\lambda_k|\mathcal E^{(k)}|
 \;\longrightarrow\;
 \sum_k\Bigl[
 \lambda_k^{\mathrm{rel}}\sum_r z_{r,k}
 +\lambda_k^{\mathrm{rep}}\sum_r z_{r,k}\bigl(|\mathcal R_k(r)|-1\bigr)
 \Bigr],
 \qquad
 \lambda_k^{\mathrm{rel}}\gg\lambda_k^{\mathrm{rep}},
 \label{eq:model_group_price}
\end{align}
so a relation is bought once and its further representations are cheap
refinements. The search of \cref{sec:learning} accordingly proposes relations
and reads the lag, quantity and basis out inside the accepted group. Self-lags
\(r=(a,a)\) are a declared part of the construction, exempt from
\(\lambda^{\mathrm{rel}}\) and reported in their own stratum, so a label's own
autocorrelation cannot outbid a cross-label relation for a slot under the
in-degree bound.

Because \(\lambda_k^{\mathrm{rel}}\) prices one purchase out of \(C_k\)
candidate relations, it carries the search's own multiplicity. Under a
marginal-uniform prior on the active set the \(j\)-th accepted relation is
priced at
\begin{align}
 \lambda_k^{\mathrm{eff}}(j)
 &=\lambda_k^{\mathrm{rel}}+\log\frac{C_k-j+1}{j},
 \label{eq:model_multiplicity_price}
\end{align}
which is what makes a per-scale edge count comparable across resolutions whose
candidate fields differ in size.

\paragraph{Bidirectional structural EM.}

For a structured posterior \(q(U^{(1:K)})\), learning maximizes the penalized
evidence lower bound. An interior scale receives evidence from its own factor
and the factor through which it generates its children,
\begin{align}
 \mathcal J(q,\Theta,\mathcal G)
 &=\mathbb E_q[\log p(Z,U,\mathcal G,\Theta\mid C)]
   +\mathbb H[q]
   -\sum_{k=1}^K\lambda_k|\mathcal E^{(k)}|
   -\lambda_\Theta\Omega(\Theta),
 \label{eq:model_elbo}\\
 \log q_k^*(U^{(k)})
 &=\mathbb E_{q_{-k}}[\log\Psi_k+\log\Psi_{k-1}]+\mathrm{constant}.
 \label{eq:model_bidirectional_update}
\end{align}
Fine-to-coarse and coarse-to-fine messages therefore precede the three
structural EM coordinates \citep{friedman1998learning},
\begin{align}
 q^+&=\arg\max_q\mathcal J(q,\Theta,\mathcal G),
 \label{eq:model_em_steps}\\
 (\Theta^{(k)})^+
 &=\arg\max_{\Theta^{(k)}}\left\{
 \mathbb E_{q^+}[\log\Psi_k]+\Lambda_k(\Theta^{(k)})
 -\lambda_\Theta\Omega_k(\Theta^{(k)})\right\},
 \label{eq:model_parameter_step}\\
 \mathcal P_{k,b}^+
 &=\arg\max_{\substack{P\subseteq\mathcal C_{k,b}\\|P|\leq D_k}}
 \left\{\max_{\beta_P}\mathbb E_{q^+}[\log\Psi_{k,b}(P,\beta_P)]
 +\Gamma_k(P)-\lambda_k|P|\right\}.
 \label{eq:model_structure_step}
\end{align}
These are respectively the E-step, M-step, and S-step. Here
\(\Lambda_k(\Theta^{(k)})=\log p(\Theta^{(k)}\mid
\Theta^{(k-1)},\Theta^{(k+1)})\) couples adjacent parameter blocks, and
\(\Gamma_k(P)\) is the adjacent-graph log prior, with unavailable neighbors
omitted at the boundary scales. The S-step profiles edge
coefficients inside each normalized local likelihood. Search may be
greedy or exhaustive under the in-degree bound \(D_k\). Each accepted edge is
checked against shifted versions of its source. Regime count, factor dimension,
and penalty are selected by held-out evidence. A sweep is retained only when
its monitored penalized evidence improves; otherwise the best state is
restored with a smaller M-step learning rate. Cells within one allocation
block share a normalizer and active set, so their exact updates remain
dependent. Participants and conditionally independent blocks can run in
parallel.

Regime permutations, factor-basis changes, and block-constant score shifts
form observational equivalence classes. Parameter recovery is therefore
reported on invariant contrasts and aligned latent states. Missing values are
marginalized within the same fibers. Fitting returns a graph and conditional
parameters at every resolution. Simulation returns one event log whose
projected counts and durations reproduce all generated coarser views.

\paragraph{A population of logs.}
\label{sec:model_participants}

A study observes \(n\) participants, and days inside one participant are
repeated measures. Index participants by \(i\) and let the graph
\(\mathcal G\) and the fixed effects be shared while the level of each
predictor varies,
\begin{align}
 \eta_{i,t,b,q}^{(k)}
 &=f_{k,b,q}\bigl(\mathcal P_{k,b,q}(t),U_{i,t}^{(k)},C_t^{(k)};\Theta^{(k)}\bigr)
 +\bigl(B_i^{(k)}\bigr)_{b,q},
 &B_i^{(k)}&\sim\mathcal N\bigl(0,\Sigma_B^{(k)}\bigr).
 \label{eq:model_participant}
\end{align}
Latent paths are participant-specific; \(\Theta^{(k)}\) and
\(\mathcal G^{(k)}\) are not. Two things follow. A relation shared across
participants borrows strength, which is what makes a per-scale graph estimable
from horizons that are short in epochs. And a label present for some
participants and absent for others cannot enter another label's score through
participant identity, because \(B_i^{(k)}\) absorbs the level that identity
would otherwise supply. Lag histories never cross a participant boundary, and
an epoch whose lagged predictors fall outside the observed interval has its
emission factor masked while the epoch itself is retained, so the sums of
\cref{prop:model_reconciliation} stay intact.

\paragraph{Forecasting.}
\label{sec:model_forecast}

Let \(\mathcal H_T\) be the history through epoch \(T\) at every resolution.
The \(h\)-step predictive distribution is generated by the same factorization,
run forward at the coarsest scale and downward through the child kernels,
\begin{align}
 p\bigl(Z^{(1:K)}_{T+1:T+h}\mid\mathcal H_T,C\bigr)
 &=\int\Bigl[\prod_{s=1}^{h}
 E_K\bigl(\boldsymbol\eta^{(K)}_{T+s}\bigr)
 \prod_{k=K-1}^{1}
 W_k\bigl(\boldsymbol\eta^{(k)}_{T+s},Z^{(k+1)}_{T+s}\bigr)\Bigr]
 \,dq\bigl(U^{(1:K)}_{T+1:T+h}\mid\mathcal H_T\bigr).
 \label{eq:model_forecast}
\end{align}
Latent states are filtered to the origin and propagated by \(R_k\); no
observation after \(T\) is read. Ancestral sampling draws each path by
unrolling \(E_K\) and then descending the chain, and lagged predictors are
replaced by the path's own earlier draws.

\begin{corollary}[Coherent forecasts]
\label{cor:model_coherent_forecast}
Every draw from \cref{eq:model_forecast} satisfies
\(\mathcal A_kZ^{(k)}_{T+s}=Z^{(k+1)}_{T+s}\) for all \(k\) and
\(s\leq h\). The predictive distribution is therefore supported on the coherent
set, and its per-resolution marginals require no reconciliation.
\end{corollary}

\begin{proof}
Immediate from \cref{prop:model_reconciliation}, whose argument uses only the
support of \(K_k\) and holds for any conditioning set.
\end{proof}

\Cref{cor:model_coherent_forecast} places the model with the methods that
obtain coherence by construction \citep{pennings2017integrated,
rangapuram2021hiere2e,das_topdown2023}, in contrast with the reconciliation
projections that map incoherent base forecasts onto the coherent subspace
\citep{hyndman2011optimal,wickramasuriya2019mint}. Two properties separate it
inside that family. The support is the integer lattice intersected with the
coherent set, so counts and durations stay whole and the zero atom is
preserved; and the conditional law at each resolution is indexed by a graph
\(\mathcal G^{(k)}\) that is learned, so the forecast carries a per-resolution
structural account alongside its predictive distribution.

\section{Learning and inference}
\label{sec:learning}

The model is fitted by a bidirectional variational structural EM procedure in the sense of
\citet{friedman1998learning}, specialized to a multiscale latent-variable model with exact
reconciliation. Three steps alternate at every scale. The E-step infers the latent state, the
M-step fits the parameters, and the S-step selects the lagged parent sets. One outer sweep runs
fine to coarse and then coarse to fine, for the reason
\cref{prop:model_block_invariance} gives: the block-constant component of a
scale-\(k\) predictor is identified only above \(k\) and the within-block
component only at \(k\), so each direction carries the component the other
cannot. The loop monitors one scalar objective, and rejects any sweep that
lowers it.

This section states the objective the three steps share, gives each step in the form the
implementation uses, proves the three properties of the local structure score that make the
S-step well posed, and closes with the cost of one sweep. Two of the propositions correspond to
defects that made earlier versions of this method fail, so they are stated and proved rather
than asserted.

\subsection{The penalized objective}
\label{sec:learning:objective}

Write the multiscale observations as
\[
\gD^{(k)}
=
\Big\{
\big(\vx_t^{(k)}, \vy_t^{(k)}, \gK_t^{(k)}\big)
\Big\}_{t \in \gT^{(k)}},
\qquad
\vx_t^{(k)} = \big(X_t^{(k,\ell)}\big)_{\ell \in \gL},
\quad
\vy_t^{(k)} = \big(Y_t^{(k,\ell)}\big)_{\ell \in \gL},
\]
and collect them as $\gX = \{\vx_t^{(k)}\}_{k,t}$, $\gY = \{\vy_t^{(k)}\}_{k,t}$ and
$\gK = \{\gK_t^{(k)}\}_{k,t}$. All scales are aggregations of one event log, so the
reconciliation identities hold exactly in the observed arrays. The latent variables are the
regime trajectories $\gS = \{S_t^{(k)}\}_{k,t}$ and the factor trajectories
$\gH = \{H_t^{(k)}\}_{k,t}$. The structure is the collection of lagged parent sets
$\gG = \{\gP_{k,\ell}\}$, and $\Theta$ collects the parameters of the regime chains, the factor
chains, the coarsest-scale emissions and the finer-scale allocations.

The target is the penalized marginal log-likelihood
\begin{align}
(\hat{\Theta}, \hat{\gG})
=
\argmax_{\Theta,\;\gG \in \mathfrak{G}_{\mathrm{adm}}}
\;
\log p_{\Theta,\gG}(\gX,\gY \mid \gK)
\;-\;
\sum_{k=1}^{K} \lambda_k \sum_{\ell \in \gL} \big|\gP_{k,\ell}\big|
\;-\;
\lambda_\Theta \, \Omega_\Theta(\Theta),
\label{eq:learn_penalized_marginal}
\end{align}
where $\mathfrak{G}_{\mathrm{adm}}$ is the class of admissible multiscale structures (lagged
observed-to-observed edges only, fixed top-down inter-scale edges, and same-slice co-movement
carried only by $S_t^{(k)}$ and $H_t^{(k)}$), $\lambda_k$ is the edge penalty at scale $k$, and
$\Omega_\Theta$ is a parameter regularizer.

The marginal in \cref{eq:learn_penalized_marginal} is intractable because of $\gS$ and $\gH$, so
the method optimizes the penalized evidence lower bound
\begin{align}
\gJ(q,\Theta,\gG)
=
\E_{q(\gS,\gH)}\!\big[\log p_{\Theta,\gG}(\gX,\gY,\gS,\gH \mid \gK)\big]
+
\mathbb{H}[q]
-
\sum_{k=1}^{K} \lambda_k \sum_{\ell \in \gL} \big|\gP_{k,\ell}\big|
-
\lambda_\Theta \, \Omega_\Theta(\Theta),
\label{eq:learn_elbo}
\end{align}
with $\mathbb{H}[q]$ the entropy of the variational posterior. For a fixed structure and fixed
parameters the first two terms of \cref{eq:learn_elbo} lower bound
$\log p_{\Theta,\gG}(\gX,\gY \mid \gK)$, with the gap equal to
$\KL\!\big(q \,\|\, p(\gS,\gH \mid \gX,\gY,\gK)\big)$.

The variational family factorizes top-down across scales, and within a scale it keeps the regime
chain intact and conditions the factor on it,
\begin{align}
q(\gS,\gH)
=
q^{(K)}\!\big(\gS^{(K)},\gH^{(K)}\big)
\prod_{k=1}^{K-1}
q^{(k)}\!\big(\gS^{(k)},\gH^{(k)} \mid \bar{\gZ}^{(k+1)}\big),
\qquad
q^{(k)} = q_S^{(k)}\big(\gS^{(k)}\big)\, q_H^{(k)}\big(\gH^{(k)} \mid \gS^{(k)}\big).
\label{eq:learn_variational_family}
\end{align}
The conditioning set $\bar{\gZ}^{(k+1)}$ holds what a parent scale passes to its children, namely
the parent regime marginals mapped through the parent map,
$\big(\gamma^{(k+1)}_{\parmap_k(t)}(\cdot)\big)_{t \in \gT^{(k)}}$, and the observed parent totals
$\vx_{\parmap_k(t)}^{(k+1)}$ and $\vy_{\parmap_k(t)}^{(k+1)}$. The parent factor mean
$\bar{\vh}^{(k+1)}$ is available and is not used, because a coarser factor acts on its children
through the allocation totals it has already shaped.

Three coordinate ascent steps follow from \cref{eq:learn_elbo}. Each holds two of the three
arguments fixed,
\begin{align}
q^{\new} &= \argmax_q \gJ(q,\Theta^{\old},\gG^{\old}),
\qquad
\Theta^{\new} = \argmax_\Theta \gJ(q^{\new},\Theta,\gG^{\old}),
\nonumber\\
\gG^{\new} &= \argmax_{\gG \in \mathfrak{G}_{\mathrm{adm}}} \gJ(q^{\new},\Theta^{\new},\gG).
\label{eq:learn_three_steps}
\end{align}
The E-step of \cref{eq:learn_three_steps} is solved to the accuracy of the structured family, the
M-step block by block, and the S-step label by label. \cref{sec:learning:outer} states what
survives of the monotonicity that coordinate ascent would otherwise supply.

The parameter penalty deserves one remark. In the reported fits $\lambda_\Theta$ is zero and
$\Omega_\Theta$ acts through the feasible set rather than through the objective. The factor block
solves ridged linear systems with ridge $10^{-4}$, the negative binomial dispersions are floored
at $0.05$, and the allocation concentrations are searched over $[0.5,\,2000]$. These are
constraints on where a block update may land, not an additive term, so every objective quoted
below omits $\lambda_\Theta \Omega_\Theta$.

\subsection{E-step: latent-state inference}
\label{sec:learning:estep}

Fix a scale $k$ and suppress the superscript where it is unambiguous. The scale-local objective
is
\begin{align}
\gJ_k(q_S,q_H)
=
\E_{q_S q_H}\!\big[\log p_{\Theta_k,\gG_k}\big(\gD^{(k)},\gS^{(k)},\gH^{(k)} \mid \bar{\gZ}^{(k+1)}\big)\big]
+
\mathbb{H}[q_S]
+
\E_{q_S}\!\big[\mathbb{H}[q_H(\cdot \mid \gS)]\big],
\label{eq:learn_scale_objective}
\end{align}
and it is maximized by alternating a discrete update of $q_S$ and a Gaussian update of $q_H$.

\paragraph{Expected local evidence.}
The regime update needs, for every epoch $t$ and every regime $m \in \{1,\dots,M_k\}$, the
expected log evidence
\begin{align}
\psi_t^{(k)}(m)
=
\E_{q_H}\!\left[
\log p\!\left(
\vx_t^{(k)}, \vy_t^{(k)}
\;\middle|\;
S_t^{(k)}=m,\; H_t^{(k)},\;
\big\{\gP_{k,\ell}^{X}(t),\, \gP_{k,\ell}^{Y}(t)\big\}_{\ell \in \gL}
\right)
\right].
\label{eq:learn_psi}
\end{align}
The expectation in \cref{eq:learn_psi} is evaluated at the posterior mean
$\bar{\vh}_t^{(k)} = \E_q[H_t^{(k)}]$, a first-order approximation that costs one pass over a
$T_k \times M_k \times L$ array instead of a quadrature.

At the coarsest scale the emission is a zero-inflated negative binomial count paired with a
negative binomial duration restricted to the epochs that carry at least one event,
\begin{align}
\psi_t^{(K)}(m)
&=
\sum_{\ell \in \gL}
\log \mathrm{ZINB}\!\big(X_t^{(K,\ell)} \mid \varpi_t^{(K,\ell)}(m),\, \phi_\ell^{(K)},\, \mu_t^{(K,\ell)}(m)\big)
\nonumber\\
&\quad
+
\sum_{\ell \in \gL}
\mathbf{1}\!\big\{X_t^{(K,\ell)} > 0\big\}\,
\log \mathrm{NB}\!\big(Y_t^{(K,\ell)} \mid s_\ell^{(K)},\, \nu_t^{(K,\ell)}(m)\big),
\label{eq:learn_psi_coarse}
\end{align}
with the three outputs of the emission map in
\cref{eq:model_coarse_count_mean} evaluated at
the regime indicator of $m$ and at the factor mean $\bar{\vh}_t^{(K)}$. The count mean
$\log\mu_t^{(K,\ell)}(m)$ and the duration mean $\log\nu_t^{(K,\ell)}(m)$ carry the lagged features
$\log(1+X_{t-d}^{(K,\ell')})$ and $\log(1+Y_{t-d}^{(K,\ell')})$ selected by $\gP_{K,\ell}$, the
duration mean carries the own-epoch term $\log(1+X_t^{(K,\ell)})$ as well, and the zero-inflation
predictor $\logit \varpi_t^{(K,\ell)}(m)$ carries no lagged term, so a lagged edge reaches the
count channel through the mean alone.

At a finer scale $k < K$ the emission is the allocation of a parent block. For a child epoch $t$
with parent $c = \parmap_k(t)$, let $\eta_t^{(k,\ell)}(m)$ and $\zeta_t^{(k,\ell)}(m)$ be the
count and duration allocation scores of \cref{eq:model_count_score} evaluated at
the regime indicator of $m$ and at the factor mean. Both are affine in their coefficients. The
count score carries the parent totals $\log(1+X_c^{(k+1,\ell)})$ and $\log(1+Y_c^{(k+1,\ell)})$;
the duration score carries $\log(1+Y_c^{(k+1,\ell)})$ and the own-epoch count
$\log(1+X_t^{(k,\ell)})$ in place of the parent count total. Both carry the depletion term
$\log(1+W_t^{(k,\ell)})$ represented in \cref{eq:model_count_omega}, which records what the earlier children of
the block have already taken, together with the context, the factor mean, and the lagged features
selected by $\gP_{k,\ell}$. The allocation probabilities $\omega$ and $\rho$ are block softmaxes
of these scores, and the emission is the Dirichlet-Multinomial allocation log-likelihood
\begin{align}
\psi_t^{(k)}(m)
&=
\sum_{\ell \in \gL}
\Big[
\log\Gamma\!\big(X_t^{(k,\ell)} + \kappa\,\omega_t^{(k,\ell)}(m)\big)
-
\log\Gamma\!\big(\kappa\,\omega_t^{(k,\ell)}(m)\big)
\Big]
\nonumber\\
&\quad
+
\sum_{\ell \,:\, X_t^{(k,\ell)}>0}
\Big[
\log\Gamma\!\big(Y_t^{(k,\ell)} + \tau\,\rho_t^{(k,\ell)}(m)\big)
-
\log\Gamma\!\big(\tau\,\rho_t^{(k,\ell)}(m)\big)
\Big],
\label{eq:learn_psi_fine}
\end{align}
up to a term that depends on the data and the concentrations alone and therefore shifts every
regime by the same amount. The concentrations $\kappa$ and $\tau$ carry the overdispersion of the
observed split. Letting $\kappa,\tau \to \infty$ recovers the multinomial limit
$X_t^{(k,\ell)} \log \omega_t^{(k,\ell)}$, which overstates the evidence and lets spurious lagged
edges clear the penalty.

The block softmax couples the children of one parent, so a per-regime allocation probability is
not defined without a treatment of the block-mates. The implementation scores the child under its
own regime and the block-mates at their regime-averaged score,
\begin{align}
\omega_t^{(k,\ell)}(m)
=
\frac{\exp\!\big(\eta_t^{(k,\ell)}(m)\big)}
{\exp\!\big(\eta_t^{(k,\ell)}(m)\big) + \sum_{s \in \child_k(c) \setminus \{t\}} \exp\!\big(\bar{\eta}_s^{(k,\ell)}\big)},
\qquad
\bar{\eta}_s^{(k,\ell)} = \sum_{m'=1}^{M_k} \gamma_s(m')\, \eta_s^{(k,\ell)}(m'),
\label{eq:learn_loo_omega}
\end{align}
and likewise for $\rho$ with the score $\zeta$ and the active-child restriction. This mean-field
treatment of the block normalizer is the one place where the objective a step is accepted on
differs from the objective the next step reports, and \cref{sec:learning:outer} handles the
consequence.

\paragraph{Regime update.}
With $\psi^{(k)}$ in hand, $q_S^{(k)}$ is the chain distribution
\begin{align}
q_S\big(\gS^{(k)}\big)
\;\propto\;
p\big(S_1^{(k)} \mid \gK_1^{(k)}\big)
\prod_{t=2}^{T_k}
p\big(S_t^{(k)} \mid S_{t-1}^{(k)}, \gK_t^{(k)}, \bar{\gZ}^{(k+1)}\big)
\prod_{t=1}^{T_k}
\exp\!\big(\psi_t^{(k)}(S_t^{(k)})\big),
\label{eq:learn_qS}
\end{align}
whose transitions replace the parent regime of \cref{eq:model_regime_init} by its posterior
marginal,
\begin{align}
p\big(S_t^{(k)} = m \mid S_{t-1}^{(k)} = m',\, \bar{\gZ}^{(k+1)}\big)
=
\softmax_m\!\Big(
a_m^{(k)} + A_{m,m'}^{(k)}
+ \sum_{n=1}^{M_{k+1}} \gamma^{(k+1)}_{\parmap_k(t)}(n)\, C_{m,n}^{(k)}
+ (\vb_m^{(k)})^\top \gK_t^{(k)}
\Big).
\label{eq:learn_trans_logit}
\end{align}
The cross-scale term enters through the expected logit. That form is exact when the parent
posterior is degenerate, and it is the same first-order form the M-step differentiates, so the
two steps optimize one objective. Because \cref{eq:learn_qS} is a Markov chain with time-varying
transitions, a log-space forward-backward recursion returns the marginals $\gamma_t(m)$, the pair
marginals $\xi_t(m,m')$ and the log evidence $\log Z_k$ exactly, at cost $O(T_k M_k^2)$.

\paragraph{Factor update.}
Conditioned on the regime path the factor dynamics are linear Gaussian, but the emissions are
not, so $q_H$ is taken Gaussian and fitted by a regime-averaged Kalman filter followed by an RTS
smoother. The averaged system at epoch $t$ is
\begin{align}
\bar{\mF}_t &= \sum_{m} \gamma_t(m)\, \mF_m,
\qquad
\bar{\vg}_t = \sum_{m} \gamma_t(m)\, \mG_m \gK_t^{(k)},
\nonumber\\
\bar{\mQ}_t &= \sum_{m} \gamma_t(m)\Big[\mSigma_m + (\mF_m - \bar{\mF}_t)\, \mP_{t-1} (\mF_m - \bar{\mF}_t)^\top\Big],
\label{eq:learn_avg_lgssm}
\end{align}
where $\mP_{t-1}$ is the filtered covariance of the previous epoch. The second term of
$\bar{\mQ}_t$ is the Jensen gap that a single averaged transition matrix would otherwise discard.

At the coarsest scale the emission carries information about the factor, and that information is
injected as a Laplace pseudo-observation. Write $\psi_t^{(K)}(m;\vh)$ for
\cref{eq:learn_psi_coarse} with the factor held at an arbitrary $\vh$ rather than at
$\bar{\vh}_t^{(K)}$, and let
\[
\bar{\psi}_t(\vh)
=
\sum_{m} \gamma_t(m)\, \psi_t^{(K)}(m;\vh)
\]
be the regime-averaged emission log-likelihood at epoch $t$ as a function of the factor. Its
gradient and Hessian at the current mean are taken by central differences, the Hessian eigenvalues
are clipped above at $-10^{-6}$ to force negative definiteness, and the filter absorbs
\begin{align}
\mR_t = \big(-\nabla^2 \bar{\psi}_t(\bar{\vh}_t)\big)^{-1},
\qquad
\vz_t = \bar{\vh}_t + \mR_t\, \nabla \bar{\psi}_t(\bar{\vh}_t)
\label{eq:learn_laplace}
\end{align}
as an observation of $H_t^{(K)}$ under the identity observation model with noise covariance
$\mR_t$. This is one Newton step on the emission, converted into the moments a Kalman update
consumes. At a finer scale no pseudo-observation is formed and the smoother runs on the prior
alone, since the factor there enters only the allocation scores, whose curvature is already
represented through the block normalizer.

\paragraph{Alternation.}
The two updates are alternated. The regime update needs $\bar{\vh}$ to build $\psi$, and the
factor update needs $\gamma$ to build $\bar{\mF}_t$ and $\bar{\psi}_t$. The loop starts from
$\bar{\vh} = \vzero$ and a uniform $\gamma$, and stops when
$|\log Z_k - \log Z_k^{\mathrm{prev}}| < \varepsilon\,(1 + |\log Z_k|)$ or after a fixed
iteration budget. One case is exact in a single pass. At the coarsest scale with $r_K = 0$
nothing feeds back into the emissions, so the loop exits immediately. At a finer scale with
$r_k = 0$ the loop still runs, because the block normalizer of \cref{eq:learn_loo_omega} depends
on $\gamma$ through $\bar{\eta}$, so the two updates still have a fixed point to reach.

\subsection{M-step: parameter learning}
\label{sec:learning:mstep}

Given $q^{\new}$ and the current structure, the parameters solve
\begin{align}
\Theta_k^{\new}
=
\argmax_{\Theta_k}
\;
\E_q\!\big[\log p_{\Theta_k,\gG_k}\big(\gD^{(k)},\gS^{(k)},\gH^{(k)} \mid \bar{\gZ}^{(k+1)}\big)\big].
\label{eq:learn_mstep}
\end{align}
Given $q$, the expected complete-data log-likelihood separates into parameter blocks that share no
arguments, so the blocks can be updated in any order. Five blocks appear in total, and one scale
uses the regime block, the factor block, and then either the two coarsest-scale emission blocks or
the single allocation block.

\begin{enumerate}
\item \textbf{Regime transitions} $\{\va^{(k)}, A^{(k)}, C^{(k)}, \vb^{(k)}\}$ together with the
initial distribution. The objective is the expected log-probability of the regime path under
\cref{eq:learn_trans_logit}, a softmax conditional likelihood whose gradient is the difference
between expected and predicted transition mass,
\[
\frac{\partial}{\partial\, \logit_t(m \mid m')}
=
\xi_{t-1}(m',m) - \gamma_{t-1}(m')\, p_t(m \mid m').
\]
It is maximized by Adam. A tabular soft-Bayes estimate of the same block can be enabled as a
second proposal, in which case the proposal with the higher block objective is taken; the fits
reported here use the gradient proposal alone.

\item \textbf{Latent factor} $\{\mF_m, \mG_m, \mSigma_m, \vmu_0, \mSigma_0\}$. The expected
complete-data log-likelihood of a switching linear Gaussian chain is quadratic in $\mF_m$ and
$\mG_m$ given the smoothed moments, so the update is closed form, a regime-weighted least squares
solve per regime with ridge $10^{-4}$, followed by the weighted residual second moment for
$\mSigma_m$. One approximation is made here. The cross moment $\E_q[H_t H_{t-1}^\top]$ is replaced
by $\bar{\vh}_t \bar{\vh}_{t-1}^\top$, dropping the smoother cross-covariance, so the update is
exact in the means and optimistic in the covariances.

\item \textbf{Coarsest-scale counts.} The zero-inflated negative binomial coefficients and the
per-label dispersion are fitted by Adam on the $\gamma$-weighted log-likelihood of
\cref{eq:learn_psi_coarse}, with a zero observation contributing to the mean gradient through its
posterior responsibility for coming from the negative binomial rather than the inflation
component.

\item \textbf{Coarsest-scale durations.} The negative binomial duration coefficients and
dispersions are fitted by Adam on the same weighted objective, restricted to the epochs with a
positive count.

\item \textbf{Finer-scale allocations.} The count and duration score coefficients are fitted by
Adam on the expected Dirichlet-Multinomial allocation log-likelihood, whose gradient with respect
to a score is the block-centered digamma residual
\[
\frac{\partial}{\partial \eta_t^{(k,\ell)}}
=
\kappa\, \omega_t^{(k,\ell)} \Big[ v_t - \sum_{s \in \child_k(c)} \omega_s^{(k,\ell)} v_s \Big],
\qquad
v_t = \Psi\!\big(X_t^{(k,\ell)} + \kappa \omega_t^{(k,\ell)}\big) - \Psi\!\big(\kappa \omega_t^{(k,\ell)}\big),
\]
with $\Psi$ the digamma function. The concentrations $\kappa$ and $\tau$ are estimated afterwards,
one per channel, by golden-section search on the log concentration over $[0.5,\,2000]$, using the
complete Dirichlet-Multinomial
log-likelihood including the block-constant term that the per-epoch form of
\cref{eq:learn_psi_fine} drops, since that term is the only part of the likelihood that
identifies a concentration. Leaving $\kappa$ and $\tau$ at an assumed value makes every structure
score that uses this likelihood over- or under-confident, so they are fitted rather than fixed.
\end{enumerate}

A fixed number of Adam steps can overshoot, so each gradient-based update is accepted only if it
does not lower the expected log-likelihood it targets, and otherwise the previous value is kept.
The regime block is checked on its own; the two coarsest-scale emission blocks are proposed
together and accepted or rejected together, and so are the two allocation blocks. The factor block
is a closed-form solve rather than a gradient step and is applied without a check. Each checked
group is therefore monotone in its own objective. That does not make the sweep monotone in the
objective the outer loop monitors, for the reason given in \cref{sec:learning:outer}.

\subsection{S-step: structure learning}
\label{sec:learning:sstep}

For each scale $k$ and label $\ell$, let
$\gC_{k,\ell} \subseteq \gL \times \{1,\dots,p_k\}$ be the candidate set of lagged parents, where
$p_k$ is the maximum lag in \cref{eq:model_lag_values}. With no screening rule the candidate
set is the whole product, so $|\gC_{k,\ell}| = L\,p_k$. Write $D_k$ for the maximum in-degree
allowed at scale $k$, a quantity distinct from the context dimension $d_k$. The learned parent set
satisfies $\gP_{k,\ell} \subseteq \gC_{k,\ell}$ and $|\gP_{k,\ell}| \le D_k$.

\paragraph{The local score.}
Let $\vbeta_\gP$ collect the lag coefficients that a parent set $\gP$ contributes to the two linear
predictors of label $\ell$, namely $\beta_{\ell',d}^{(k,\ell)}$ in the count channel and
$\alpha_{\ell',d}^{(k,\ell)}$ in the duration channel for each $(\ell',d) \in \gP$, and let
$\Theta_{-\gP}$ be everything else, held at its current value. The local score is the expected
local log-likelihood \emph{maximized over} $\vbeta_\gP$,
\begin{align}
Q_{k,\ell}(\gP; q,\Theta)
=
\max_{\vbeta_\gP}
\;
\sum_{t \in \gT^{(k)}} \sum_{m=1}^{M_k}
\gamma_t(m)\,
\log p_{\vbeta_\gP,\Theta_{-\gP}}\!\left(
X_t^{(k,\ell)}, Y_t^{(k,\ell)}
\;\middle|\;
S_t^{(k)}=m,\; \bar{\vh}_t^{(k)},\; \gP^{X}(t),\, \gP^{Y}(t)
\right),
\label{eq:learn_local_score}
\end{align}
with $\gP^{X}(t)$ and $\gP^{Y}(t)$ the lagged parent values of \cref{eq:model_lag_values}, and the
structure update at $(k,\ell)$ is
\begin{align}
\hat{\gP}_{k,\ell}
=
\argmax_{\gP \subseteq \gC_{k,\ell},\; |\gP| \le D_k}
\;
Q_{k,\ell}(\gP; q,\Theta) - \lambda_k |\gP|.
\label{eq:learn_local_update}
\end{align}
At the coarsest scale the density in \cref{eq:learn_local_score} is the zero-inflated negative
binomial paired with the negative binomial duration, and the implementation evaluates the
$\gamma$-weighted sum exactly as written. At a finer scale it is the Dirichlet-Multinomial
allocation likelihood of \cref{eq:learn_psi_fine} with the block normalizer included, and the
implementation makes one further economy. Rather than form the leave-one-out allocation of
\cref{eq:learn_loo_omega} once per regime, it evaluates the allocation likelihood once at the
regime-averaged score $\bar{\eta}$, which is the $\gamma$-weighted sum of
\cref{eq:learn_local_score} to first order and coincides with it when the regime posterior is
degenerate. The gradient that drives the inner maximization is still contracted against $\gamma$.
Two features of \cref{eq:learn_local_score} are load-bearing, the block normalizer and the
maximization over $\vbeta_\gP$. Both were absent from earlier versions of the method, and the next
two propositions say exactly what their absence costs.

\begin{proposition}[An unnormalized local score has no maximizer]
\label{prop:learn_unnormalized}
Fix a scale $k < K$, a label $\ell$ and a parent set $\gP$. Let $\vz_t \in \sR^{J}$ collect the
features entering the count allocation score of epoch $t$, so that
$\eta_t^{(k,\ell)} = \vtheta^\top \vz_t$ for a coefficient vector $\vtheta \in \sR^{J}$, and
define the unnormalized local score
\begin{align}
\widetilde{Q}_{k,\ell}(\gP;\vtheta)
=
\sum_{t \in \gT^{(k)}} X_t^{(k,\ell)}\, \eta_t^{(k,\ell)}
=
\vtheta^\top \vw,
\qquad
\vw = \sum_{t \in \gT^{(k)}} X_t^{(k,\ell)}\, \vz_t .
\label{eq:learn_unnormalized_score}
\end{align}
If $\vw \neq \vzero$ then
$\sup_{\vtheta} \big[\widetilde{Q}_{k,\ell}(\gP;\vtheta) - \lambda_k |\gP|\big] = +\infty$ and the
supremum is not attained. Consequently, for any two parent sets whose feature sums are both
nonzero, the profiled unnormalized scores are both infinite, and their comparison carries no
information about the parent sets.
\end{proposition}

\begin{proof}
$\widetilde{Q}_{k,\ell}$ is linear in $\vtheta$, so for $c>0$ it satisfies
$\widetilde{Q}_{k,\ell}(\gP; c\vtheta) = c\,\vtheta^\top \vw$. Take $\vtheta = \vw$; then
$\widetilde{Q}_{k,\ell}(\gP; c\vw) = c\|\vw\|^2 \to \infty$ as $c \to \infty$, because
$\|\vw\|^2 > 0$. The penalty $-\lambda_k|\gP|$ is a finite constant and does not affect the
divergence, and no finite $\vtheta$ attains an infinite value. The same argument applies to any
other parent set with a nonzero feature sum, so \cref{eq:learn_local_update} would compare
$+\infty$ with $+\infty$.
\end{proof}

\begin{remark}
\label{rem:learn_normalizer_bounds}
The block-normalized score is bounded above and the difficulty disappears. Under the multinomial
allocation, $\sum_t X_t^{(k,\ell)} \log \omega_t^{(k,\ell)} \le 0$ because each
$\omega_t^{(k,\ell)} \in (0,1]$. Under the Dirichlet-Multinomial form of
\cref{eq:learn_psi_fine}, the per-epoch term satisfies, for integer $X \ge 0$ and
$a = \kappa\omega \in (0,\kappa]$,
\[
\log\Gamma(X+a) - \log\Gamma(a) = \sum_{j=0}^{X-1} \log(a+j) \;\le\; \sum_{j=0}^{X-1} \log(\kappa+j),
\]
since the middle expression is nondecreasing in $a$ and the right side is its value at
$\omega = 1$. The right side depends on the data and on $\kappa$ alone, so summing over epochs and
labels, and repeating the argument on the duration channel with $\tau$, bounds the whole score by
a constant free of the allocation coefficients. The supremum of
\cref{eq:learn_local_score} is therefore finite and two parent sets are compared as two finite
numbers. The block normalizer is what supplies the bound, because it is the only term that
charges for raising one child's score without lowering its block-mates.
\end{remark}

\begin{proposition}[Scoring at the incumbent coefficient can only delete edges]
\label{prop:learn_profiling}
Let $\widehat{Q}_{k,\ell}(\gP;\Theta)$ denote the local expected log-likelihood of
\cref{eq:learn_local_score} evaluated at the current $\Theta$ rather than maximized over
$\vbeta_\gP$. Fix $\gP \subseteq \gC_{k,\ell}$ and suppose the parameterization holds the count and
duration lag coefficients of every candidate outside $\gP$ at zero, that is
$\beta_e^{(k,\ell)} = \alpha_e^{(k,\ell)} = 0$ for every $e \in \gC_{k,\ell} \setminus \gP$, which
is what a parameterization indexed by the current parent set supplies when $\gP$ is that set. Let
$\lambda_k > 0$.
\begin{enumerate}
\item For every $e \in \gC_{k,\ell} \setminus \gP$,
\begin{align}
\Big[\widehat{Q}_{k,\ell}(\gP \cup \{e\};\Theta) - \lambda_k|\gP \cup \{e\}|\Big]
-
\Big[\widehat{Q}_{k,\ell}(\gP;\Theta) - \lambda_k|\gP|\Big]
=
-\lambda_k \;<\; 0 .
\label{eq:learn_addition_is_minus_lambda}
\end{align}
A search that accepts a move only when the penalized score strictly increases therefore never
adds an edge, and the parent set it returns is a subset of the one it started from.
\item The profiled score $Q_{k,\ell}$ of \cref{eq:learn_local_score} is monotone under enlargement
of the parent set,
\begin{align}
Q_{k,\ell}(\gP \cup \{e\}) \;\ge\; Q_{k,\ell}(\gP),
\qquad\text{so}\qquad
Q_{k,\ell}(\gP \cup \{e\}) - Q_{k,\ell}(\gP) - \lambda_k \;\ge\; -\lambda_k ,
\label{eq:learn_profiled_gain}
\end{align}
with equality exactly when the extra coefficient buys no improvement in fit, and strict positivity
exactly when the best achievable improvement exceeds $\lambda_k$.
\end{enumerate}
\end{proposition}

\begin{proof}
For part 1, the parent set enters the local log-likelihood only through the lagged term
$\sum_{(\ell',d) \in \gP} \beta_{\ell',d}^{(k,\ell)} \log\!\big(1 + X_{t-d}^{(k,\ell')}\big)$ of
the count predictor and its counterpart
$\sum_{(\ell',d) \in \gP} \alpha_{\ell',d}^{(k,\ell)} \log\!\big(1 + Y_{t-d}^{(k,\ell')}\big)$ in
the duration predictor. Adding $e = (\ell',d)$ with $\beta_e^{(k,\ell)} = \alpha_e^{(k,\ell)} = 0$
adds the terms $0 \cdot \log(1+X_{t-d}^{(k,\ell')}) = 0$ and $0 \cdot \log(1+Y_{t-d}^{(k,\ell')})
= 0$ to the predictors of every epoch. The predictors are therefore unchanged, hence so is every
emission factor, hence so is $\widehat{Q}_{k,\ell}$. The penalty grows by exactly $\lambda_k$,
which gives \cref{eq:learn_addition_is_minus_lambda}. A hill climb that keeps the parameterization
indexed by its current parent set meets the hypothesis at every state it reaches, so no addition
ever increases the penalized score, the only accepted moves are removals, and the returned set is
contained in the initial set.

For part 2, $\vbeta_\gP$ ranges over $\sR^{2|\gP|}$, two coefficients per edge. The map
$\vbeta_\gP \mapsto (\vbeta_\gP, \vzero)$ embeds $\sR^{2|\gP|}$ into $\sR^{2|\gP|+2}$, the
coefficient space of $\gP \cup \{e\}$, and leaves the objective value unchanged, by the same
argument as part 1. A maximum over a superset is at least the maximum over the subset, which gives
\cref{eq:learn_profiled_gain}. Equality holds exactly when the maximizer over the larger space
attains no more than the maximizer over the smaller one, that is, when the additional coefficients
buy nothing.
\end{proof}

\begin{remark}
\label{rem:learn_profile_practice}
Profiling the candidate's coefficients before comparing is therefore what makes an addition
possible at all, and it is what gives $\lambda_k$ its intended meaning as the improvement in fit
that a new coefficient must buy. The implementation approximates the inner maximum of
\cref{eq:learn_local_score} by normalized-gradient ascent from the incumbent coefficients, with a
backtracking line search that accepts only improving steps, so the returned value is never below
the value at the starting point. The monotonicity of \cref{eq:learn_profiled_gain} therefore
holds exactly for the true maximizer and up to the accuracy of the ascent in practice. The fitted
coefficients of the winning parent set are written back into $\Theta$, so a newly accepted edge
starts from a fitted value rather than from zero.
\end{remark}

\begin{proposition}[Per-label decomposition]
\label{prop:learn_decomposition}
Fix $q$ and $\Theta$.
\begin{enumerate}
\item For every assignment $(\gP_{k,\ell}) \in \prod_{k,\ell} 2^{\gC_{k,\ell}}$, the unrolled
multiscale graph is acyclic.
\item Write $\gJ^{\ast}(q,\Theta,\gG)$ for the value of $\gJ$ after maximization over the lag
coefficients $\{\vbeta_{\gP_{k,\ell}}\}_{k,\ell}$ alone, with the rest of $\Theta$ held fixed.
Then the dependence of $\gJ^{\ast}$ on $\gG$ is a sum of per-label terms,
\begin{align}
\gJ^{\ast}(q,\Theta,\gG)
=
R(q,\Theta)
+
\sum_{k=1}^{K} \sum_{\ell \in \gL}
\Big[ Q_{k,\ell}(\gP_{k,\ell}; q,\Theta) - \lambda_k |\gP_{k,\ell}| \Big],
\label{eq:learn_separable}
\end{align}
where $R$ does not depend on $\gG$.
\item Consequently the joint structure problem splits into $K L$ independent subset-selection
problems, and the collection of per-label maximizers of \cref{eq:learn_local_update} is a
maximizer of $\gJ^{\ast}$ over $\mathfrak{G}_{\mathrm{adm}}$.
\end{enumerate}
\end{proposition}

\begin{proof}
For part 1, give each node of the unrolled graph the key $\big(-k,\; t,\; \sigma\big)$, compared
lexicographically, where $k$ is the scale, $t$ is the epoch index within that scale, and
$\sigma \in \{0,1,2,3\}$ takes the value $0$ on the regime node $S_t^{(k)}$, $1$ on the factor node
$H_t^{(k)}$, $2$ on the count nodes $X_t^{(k,\ell)}$, and $3$ on the duration nodes
$Y_t^{(k,\ell)}$. Every admissible edge strictly increases this key. An inter-scale edge runs from
scale $k+1$ to scale $k$ and increases the first component, and no edge runs the other way. A
within-scale edge between two epochs runs from epoch $t-d$ to epoch $t$ with $d \ge 1$, leaving the
first component fixed and increasing the second, which covers the lagged observed edges, the two
latent chains, and the depletion feature, since the latter reads only children of the same block
that precede $t$. Within one epoch, $S_t^{(k)} \to H_t^{(k)}$ raises $\sigma$ from $0$ to $1$, a
latent-to-observed edge raises $\sigma$ from $0$ or $1$ to $2$ or $3$, and the count-to-duration
edge $X_t^{(k,\ell)} \to Y_t^{(k,\ell)}$ raises $\sigma$ from $2$ to $3$. A directed cycle would
have to return to its starting key, which a strictly increasing sequence cannot do. Since the
candidate sets contain no lag-zero pair, the argument covers every assignment, and no acyclicity
constraint has to be imposed on the search.

For part 2, the expected complete-data log-likelihood at scale $k$ splits into a regime-chain
term, a factor-chain term, and the per-epoch emission terms of \cref{eq:learn_psi_coarse} or
\cref{eq:learn_psi_fine}. The first two do not involve $\gG$. The emission terms are sums over
labels, and $\gP_{k,\ell}$ enters only the linear predictor of label $\ell$. At a finer scale the
block normalizer of \cref{eq:learn_loo_omega} normalizes over the children of one parent block
for one fixed label, so it couples epochs within a label but never two labels. The entropy
$\mathbb{H}[q]$ and the parameter penalty do not involve $\gG$. The coefficients
$\vbeta_{\gP_{k,\ell}}$ appear in the label-$\ell$ terms of scale $k$ and in no other term, so the
maximization defining $\gJ^{\ast}$ distributes over the labels and scales, and collecting the
label-$\ell$ terms and adding the edge penalty gives \cref{eq:learn_separable} with the profiled
local score of \cref{eq:learn_local_score}.

For part 3, the feasible set
$\prod_{k,\ell} \{\gP \subseteq \gC_{k,\ell} : |\gP| \le D_k\}$ is a product and, by part 1,
carries no joint constraint. A sum of terms over a product set is maximized coordinatewise, so
maximizing each bracket of \cref{eq:learn_separable} separately maximizes the whole.
\end{proof}

\begin{corollary}[Exact per-target search is feasible]
\label{cor:learn_exact_search}
Under the in-degree bound $D_k$, an exhaustive search of one subproblem in
\cref{eq:learn_local_update} evaluates
$\sum_{j=0}^{D_k} \binom{|\gC_{k,\ell}|}{j} = O\!\big(|\gC_{k,\ell}|^{D_k}\big)$ parent sets, which
is polynomial in the candidate-set size for fixed $D_k$. By \cref{prop:learn_decomposition},
solving the $L$ subproblems at a scale exactly solves the structure problem at that scale exactly,
with no acyclicity projection and no search over variable orderings.
\end{corollary}

The decomposability in \cref{eq:learn_separable} is the property that score-based Bayesian network
learning relies on generally \citep{heckerman1995learning,chickering2002optimal}; what is added
here is that the lagged restriction removes the acyclicity constraint that normally couples the
per-target problems back together \citep{murphy2002dbn,pamfil2020dynotears}.

\paragraph{Search.}
The default search is greedy forward-backward hill climbing, in the manner of the local search
stage of max-min hill climbing \citep{tsamardinos2006max}. Each pass evaluates every single-edge
addition that respects the in-degree bound and every single-edge removal, accepts the moves that
improve the penalized score, and stops when a pass finds no improvement. Scores are memoized on
the parent set, since the passes revisit the same sets.

Greedy hill climbing accepts one edge at a time, so a pair of edges that pays only jointly is
unreachable from the incumbent set. Three alternatives are implemented against that. Beam search
carries a fixed-width frontier of parent sets of increasing size. Exhaustive enumeration is exact,
and by \cref{cor:learn_exact_search} it is affordable whenever the subset count within the
in-degree bound is small, which holds for every fit reported here. A genetic search treats a
parent set as an individual and recombines and mutates them, for candidate sets too large to
enumerate. \Cref{sec:results:search} reports that all three return the same parent sets as the
greedy search on both benchmarks, so on these data the greedy answer is already the certified
per-target optimum.

The edge penalty is $\lambda_k = \lambda \log T_k$, the Bayesian information criterion price of
the two coefficients (one count, one duration) that an edge adds at a scale with $T_k$ epochs, so
that a single $\lambda$ behaves consistently across scales and sample sizes. The value of
$\lambda$ is not a constant of the method. It is selected together with the latent orders on
held-out log evidence, as described next.

\paragraph{Selecting the latent orders and the penalty.}
The regime count $M_k$, the factor dimension $r_k$ and the penalty multiplier $\lambda$ are free
quantities of the model rather than quantities a single fit determines. Each candidate triple is
fitted on a training split by the identical procedure and scored by
$-\sum_k \log Z_k$ on a held-out split, and the triple with the lowest score is kept. Nothing else
is selected this way. The graph budget $(\text{max lag}, D_k)$ and the block lengths of the
hierarchy state which lags and which timescales are of interest, and a held-out score cannot
supply a question, so both remain declarations that the protocol discloses.

\subsection{The outer loop and its monotonicity guard}
\label{sec:learning:outer}

One sweep is two passes. The fine-to-coarse pass runs the three steps at every
scale from $k=1$ up to $k=K$; the coarse-to-fine pass runs them from $k=K$ back
down to $k=1$, passing $\bar{\gZ}^{(k+1)}$ to each child scale.
\Cref{prop:model_block_invariance} is what makes both necessary. A scale-$k$
predictor splits into a block-constant part, which the child kernel is blind to
and which scale $k+1$ identifies, and a within-block part, which scale $k$
identifies. Each pass carries the part the other leaves unidentified. The
monitored objective is
\begin{align}
\gJ^{\mathrm{sweep}}
=
\sum_{k=1}^{K} \log Z_k
-
\sum_{k=1}^{K} \lambda_k \sum_{\ell \in \gL} \big|\gP_{k,\ell}\big|,
\label{eq:learn_monitored}
\end{align}
the summed per-scale log evidence returned by the forward-backward recursion, less the edge
penalty the structure step trades against it, under the same penalty schedule
$\lambda_k = \lambda \log T_k$ the structure step uses. The log evidence $\log Z_k$ is the
scale-local bound on the regime chain at the current expected local evidence, so it already
carries the entropy of $q_S^{(k)}$. Tracking the log evidence alone would leave out the
term the structure step trades against, and a sweep that correctly drops an edge lowers
$\sum_k \log Z_k$ while raising $\gJ^{\mathrm{sweep}}$, so the log evidence on its own would read
as a step backwards. The convergence
test is relative. The fit stops once the last two changes in \cref{eq:learn_monitored} each
satisfy $|\Delta \gJ^{\mathrm{sweep}}| < \varepsilon\,(1 + |\gJ^{\mathrm{sweep}}|)$. An absolute
test cannot fire when the objective is in the hundreds or thousands, and a fit under an absolute
test runs its full sweep budget however long it has been converged.

Coordinate ascent on \cref{eq:learn_elbo} would make \cref{eq:learn_monitored} nondecreasing, and
it does not, for one identifiable reason. The allocation normalizer of a finer scale is evaluated
at the regime-averaged score, as in \cref{eq:learn_loo_omega}. The M-step accepts an allocation
block when the allocation log-likelihood computed at
$\bar{\eta} = \sum_m \gamma(m)\,\eta(m)$ does not decrease. The next E-step reports $\log Z_k$
from a likelihood in which the child under evaluation carries its own regime's score while its
block-mates carry $\bar{\eta}$. The two quantities agree when the regime posterior is degenerate
and differ otherwise, so the quantity a parameter block is accepted on is not the quantity the
next inference step reports, and a sweep in which every block passed its own check can still lower
\cref{eq:learn_monitored}.

The outer loop turns that residual into a guard rather than a decreasing trace. A sweep whose
objective falls below the best sweep so far is rejected. The state (parameters, graph and
posterior summaries) reverts to the best sweep and the sweep is retried with the M-step learning
rate halved. Two consecutive rejections are retried this way and a third stops the fit at its best
state; an accepted sweep resets the count. A rejection whose
shortfall is itself below the convergence tolerance is read as convergence rather than as failure
and stops the fit immediately. The fit returns its best sweep rather than its last one, and the
rejected objectives are recorded separately, so the rejections stay measurable while the recorded
trace stays monotone. The five simulation experiments reported in this monograph record 1176 runs,
of which 1056 drive the outer loop; the remaining 120 are the oracle-parameter reference, which
runs one E-step and no sweep. Of those 1056, 860 triggered the guard at least once, and 1893
sweeps were rejected in total. The residual the guard absorbs is therefore a routine feature of
the fit rather than an edge case.

\subsection{Shift confirmation of accepted edges}
\label{sec:learning:confirm}

The penalty $\lambda_k$ is a fixed price, and at a scale where a persistent latent regime or a
shared parent-block total autocorrelates every label, a lagged feature that carries no information
about the target can still clear it. Each edge the search accepts is therefore put through a
second test against a null built from its own feature.

Fix an accepted edge $e \in \hat{\gP}_{k,\ell}$ and let
$\gP_{-e} = \hat{\gP}_{k,\ell} \setminus \{e\}$. The realized gain of the edge is
$G_e = Q_{k,\ell}(\hat{\gP}_{k,\ell}) - Q_{k,\ell}(\gP_{-e})$, penalties excluded. For
$b = 1,\dots,B$, a surrogate $\tilde{e}^{(b)}$ is built by displacing the rows of the design
column of $e$, and its gain over the same base set,
$G^{(b)} = Q_{k,\ell}\big(\gP_{-e} \cup \{\tilde{e}^{(b)}\}\big) - Q_{k,\ell}(\gP_{-e})$, is
recorded. The edge is kept only if $G_e$ exceeds an upper quantile of $\{G^{(b)}\}_{b=1}^{B}$ at a
level fixed in advance. The penalty cancels from both sides of that comparison, since the real
edge and every surrogate enlarge the base set by one. The experiments reported here use $B = 32$
copies and the $95\%$ quantile.

The displacement is chosen so that the null retains everything about the feature except its
alignment with the target. At the coarsest scale it is a whole-series circular shift by a random
offset drawn away from the endpoints, which preserves the marginal distribution of the feature and
its autocorrelation. A plain row shuffle would destroy the autocorrelation as well and would
understate the floor, because a lagged feature is attractive to the score precisely for being
autocorrelated, which is also what the latent regime explains. At a finer scale the column is
rolled independently inside each parent block, so block membership survives into the null. Gains
that come from the shared block total, from a persistent latent state, or from the realized
allocation split of a block therefore appear in the null too, and only epoch-level alignment can
separate a real edge from it. The test runs to a fixed point, since removing one edge changes the
base set of the survivors and hence their gains, so the survivors are rechecked after each
removal.

The scope of the test is worth stating plainly. It removes edges whose apparent gain is matched by
a feature with the same marginal distribution, the same autocorrelation and, at a finer scale, the
same block membership. It does not remove an edge whose source is genuinely aligned with the
target at the epoch level, whether or not that alignment is causal, so a common driver that is
itself epoch-aligned survives. It applies only to the edges the search proposed, so it can lower
recall but can never raise it, and it makes no correction for the number of candidates the search
examined before proposing them. It costs $B$ additional score evaluations per accepted edge per
pass rather than per candidate.

\subsection{Computational and memory complexity}
\label{sec:learning:complexity}

Write $T_k$ for the number of epochs at scale $k$, $L = |\gL|$ for the number of labels, $M_k$ for
the number of regimes, $r_k$ for the factor dimension, $d_k$ for the context dimension, $D_k$ for
the in-degree bound, and $\bar{C}_k = L^{-1}\sum_{\ell \in \gL} |\gC_{k,\ell}|$ for the average
candidate-set size, which equals $L p_k$ when nothing is screened. Let $I_E$ be the number of
E-step alternations, $n_\Theta$ the number of M-step gradient steps per block, $n_Q$ the number of
gradient steps used to profile one candidate parent set, and $B$ the number of shifted copies per
accepted edge.

\paragraph{E-step.}
The lagged design contributions are built once per E-step, at cost $O(T_k L \bar{C}_k)$. Each
alternation then evaluates the emission arrays, which have shape $T_k \times M_k \times L$, runs
the forward-backward recursion, and runs the smoother,
\begin{align}
O\!\Big(
T_k L \bar{C}_k
\;+\;
I_E \big[\, T_k M_k L + T_k M_k^2 + T_k r_k^3 \,\big]
\Big).
\label{eq:learn_cost_estep}
\end{align}
At the coarsest scale with $r_K > 0$ the Laplace step of \cref{eq:learn_laplace} adds a finite
difference gradient and Hessian, which take $2 r_K$ and $2 r_K (r_K + 1)$ evaluations of
$\bar{\psi}$ respectively, plus one eigendecomposition per epoch,
\begin{align}
O\!\Big(
I_E \big[\, T_K r_K^2 M_K L + T_K r_K^3 \,\big]
\Big).
\label{eq:learn_cost_laplace}
\end{align}

\paragraph{M-step.}
The gradient-based blocks cost one pass over the same arrays per step, and the factor block is a
closed-form solve per regime,
\begin{align}
O\!\Big(
n_\Theta \big[\, T_k M_k L + T_k L \bar{C}_k + T_k M_k^2 \,\big]
\;+\;
M_k r_k^3 + T_k M_k r_k^2
\;+\;
T_k L
\Big),
\label{eq:learn_cost_mstep}
\end{align}
where the last term is the golden-section search for the concentrations, whose iteration count is
fixed.

\paragraph{S-step.}
Evaluating \cref{eq:learn_local_score} for one parent set of size at most $D_k$ profiles at most
$D_k$ coefficients per channel, and each inner step touches a $T_k \times M_k$ array and a
$T_k \times D_k$ design matrix, so one score evaluation costs $O\!\big(n_Q T_k (M_k + D_k)\big)$.
Greedy hill climbing evaluates $O(D_k |\gC_{k,\ell}|)$ parent sets per label, since the in-degree
bound caps the number of productive passes. Exhaustive enumeration under the same bound evaluates
$O\!\big(|\gC_{k,\ell}|^{D_k}\big)$ of them. Shift confirmation adds $O(B D_k^2)$ evaluations per
label, since each of at most $D_k$ accepted edges draws $B$ surrogates and every removal restarts
the check. The S-step at scale $k$ therefore costs
\begin{align}
O\!\Big(
n_Q T_k (M_k + D_k)\; L \big[\, D_k \bar{C}_k + B D_k^2 \,\big]
\Big)
\quad\text{greedy},
\qquad
O\!\Big(
n_Q T_k (M_k + D_k)\; L \, \bar{C}_k^{\,D_k}
\Big)
\quad\text{exhaustive}.
\label{eq:learn_cost_sstep}
\end{align}

\paragraph{One sweep.}
Summing \cref{eq:learn_cost_estep,eq:learn_cost_mstep,eq:learn_cost_sstep} over the scales gives
the cost of one outer iteration under the greedy search,
\begin{align}
O\!\left(
\sum_{k=1}^{K}
T_k
\Big[
I_E \big( M_k L + M_k^2 + r_k^3 \big)
+
n_\Theta \big( M_k L + L \bar{C}_k + M_k^2 \big)
+
n_Q (M_k + D_k)\, L D_k \big( \bar{C}_k + B D_k \big)
\Big]
\right).
\label{eq:learn_cost_sweep}
\end{align}
The structure term dominates whenever the candidate set is not screened, which is the setting the
model is usually run in, so candidate screening is the first lever on wall-clock cost. The regime
term $M_k^2$ and the factor term $r_k^3$ matter only at large latent orders.

\paragraph{Memory.}
The dominant arrays are the observed data, the emission tables of shape
$T_k \times M_k \times L$, the pair marginals $\xi_t(m,m')$ of shape
$(T_k - 1) \times M_k \times M_k$, and the smoothed factor covariances of shape
$T_k \times r_k \times r_k$,
\begin{align}
O\!\left(
\sum_{k=1}^{K}
T_k
\big[\, L + M_k L + M_k^2 + r_k^2 \,\big]
\right).
\label{eq:learn_cost_memory}
\end{align}
Beyond the posterior summaries $\bar{\gZ}^{(k+1)}$, which occupy
$O\!\big(T_{k+1}(M_{k+1} + r_{k+1} + L)\big)$, no step needs more than one scale's arrays
resident at a time, so either pass of a sweep can stream the scales.

\section{Identifiability, event reconstruction and modeling assumptions}
\label{sec:identifiability}

The model is expressive, and expressiveness buys ambiguity. A persistent shift
in the rate of a label can be written as a regime effect, as a slow drift of
the latent factor, or as a lagged self edge; a block of the hierarchy can be
filled in the same way by many different allocation scores. Part of that
ambiguity is statistical and shrinks as the record grows. Part of it is exact
and survives any sample size, because the parameterization contains directions
along which the observed distribution does not move at all. The two kinds have
to be handled differently. The first is a question of power and is settled by
screening a scenario before a metric is read from it; the second is a question
of what a metric is entitled to measure, and no amount of data settles it.

This section fixes the finest-scale conventions that turn the bottom of the
hierarchy into an event log, states the exact invariances as a proposition,
derives from them the quantities a parameter-recovery metric may compare, and
collects the assumptions, failure modes and safeguards under which the learned
structure and parameters carry an interpretation.

\subsection{Durations in finest-scale slots}
\label{sec:slot_durations}

Let $\Delta_1 > 0$ be the wall-clock length of one epoch at the finest scale,
called a slot. Every duration in the hierarchy is expressed as a whole number
of slots,
\begin{align}
Y_t^{(k,\ell)}
:=
\frac{\text{total wall-clock length of the episodes of label } \ell
      \text{ that start in epoch } t \text{ at scale } k}{\Delta_1}
\;\in\; \sN_0 ,
\label{eq:slot_normalized_duration}
\end{align}
so that the wall-clock duration is recovered as $\Delta_1 Y_t^{(k,\ell)}$. An
episode is booked in full to the epoch in which it starts, which is what makes
\cref{eq:slot_normalized_duration} additive over the children of a block even
when an episode runs past the end of the epoch it started in. The convention is
not cosmetic, and three consequences of it are used throughout.

First, one unit of $Y$ denotes the same amount of time at every scale, so the
duration reconciliation constraint
$\sum_{t \in \child_k(c)} Y_t^{(k,\ell)} = Y_c^{(k+1,\ell)}$
is a statement about integers. A multinomial split of an integer total
satisfies it exactly and almost surely, where a continuous duration model would
have to be corrected afterwards by a projection, as hierarchical forecast
reconciliation does \citep{wickramasuriya2019mint}.

Second, the entire model becomes discrete. Durations are then allocated by the
same Dirichlet-multinomial mechanism as counts, over the same child blocks and
with the same softmax scores, which is why the invariances of
\cref{prop:invariances} apply to the count and the duration branch in identical
form.

Third, at the finest scale a duration is an extent rather than a summary. An
episode that starts in slot $t$ and lasts $y$ slots occupies slots $t$ through
$t + y - 1$, so the reconstruction of \cref{sec:event_reconstruction} needs no
rounding step and introduces no error of its own.

Write $n_k$ for the number of finest-scale slots inside one epoch at scale $k$,
as in \cref{sec:model_scales}, so that $n_1 = 1$ and $n_K$ is the largest block
in the hierarchy. Once episodes are atomic in the sense of
\cref{assm:atomic_resolution}, an epoch at scale $k$ can start at most one
episode of a label per slot, which gives the standing bound
$X_t^{(k,\ell)} \le n_k$ on counts. Durations carry no matching per-epoch
bound. A duration is booked in full to the epoch in which its episode starts,
so $Y_t^{(k,\ell)}$ measures an extent that may run past the end of epoch $t$
and is bounded by the parent total it was allocated from,
$Y_t^{(k,\ell)} \le Y_{\parmap_k(t)}^{(k+1,\ell)}$. Iterating that bound to the
top and using the coarsest-scale ceiling of \cref{sec:atomic_resolution} gives
$Y_t^{(k,\ell)} \le n_K$ at every scale, which is a statement about the record
and not about one epoch. In particular $Y_t^{(1,\ell)}$ is the length of the
episode starting in slot $t$ and is routinely larger than one.

\subsection{Atomic resolution and the capacity cap}
\label{sec:atomic_resolution}

\begin{assumption}[Atomic resolution at the finest scale]
\label{assm:atomic_resolution}
The finest scale is chosen fine enough that a slot starts at most one episode
of any given label,
\[
X_t^{(1,\ell)} \in \{0,1\}
\qquad
\text{for all } t \in \gT^{(1)} \text{ and all } \ell \in \gL .
\]
\end{assumption}

\Cref{assm:atomic_resolution} is a statement about the choice of $\Delta_1$,
not about the mechanism that generates behavior. Under it, $X_t^{(1,\ell)}$
stops being a count and becomes an indicator that a new episode of label $\ell$
begins in slot $t$, which is what makes the finest scale reconstructible as an
event log. If the assumption fails for some label, the chosen resolution is too
coarse for that label, and the remedy is either a finer $\Delta_1$ or a
label-specific micro-allocation step below the finest scale. The assumption is
not free for the generative model either, since a Dirichlet-multinomial split
of a parent total places no bound on how much of that total lands in one child.
Two devices enforce it.

\paragraph{Capacity cap with redistribution.}
Fix a scale $k < K$, a parent epoch $c \in \gT^{(k+1)}$ and a label $\ell$, and
let $(\omega_t^{(k,\ell)})_{t \in \child_k(c)}$ be the softmax allocation
weights of \cref{eq:model_count_omega}. Every child carries the capacity $n_k$.
Write $(\tilde{X}_t^{(k,\ell)})_{t \in \child_k(c)}$ for the drawn allocation
and
\begin{align}
E_c^{(k,\ell)}
=
\sum_{t \in \child_k(c)}
\max\!\left( \tilde{X}_t^{(k,\ell)} - n_k, \; 0 \right)
\label{eq:capacity_excess}
\end{align}
for the part of it that does not fit. The draw is clipped,
\begin{align}
X_t^{(k,\ell)} = \min\!\left( \tilde{X}_t^{(k,\ell)}, \; n_k \right),
\label{eq:capacity_clip}
\end{align}
and the $E_c^{(k,\ell)}$ removed units are then reassigned one at a time, each
unit going to a child drawn from those that still have room,
\begin{align}
\Pr\!\left[\text{unit goes to } t\right]
=
\frac{\omega_t^{(k,\ell)} \,
      \mathbb{I}\{X_t^{(k,\ell)} < n_k\}}
     {\sum_{s \in \child_k(c)} \omega_s^{(k,\ell)} \,
      \mathbb{I}\{X_s^{(k,\ell)} < n_k\}} ,
\label{eq:capacity_redistribution}
\end{align}
with $X^{(k,\ell)}$ updated before the next unit is placed. The operation moves
units between siblings of one block and creates none, so the block total of
\cref{eq:model_reconciliation} holds after the repair exactly as it held
before it, and reconciliation across the whole hierarchy is untouched. The
receiver is drawn along the same allocation weights that produced the draw
being repaired, so the repair expresses no preference the model does not
already carry. It is well posed whenever the parent total fits in the block,
$X_c^{(k+1,\ell)} \le |\child_k(c)| \, n_k = n_{k+1}$, which is the count bound
of \cref{sec:slot_durations}. At the coarsest scale there is no parent total to
preserve, so the cap is applied directly to the emission draw, which truncates
the count and the duration emission at $n_K$ rather than moving units between
epochs.

\paragraph{Duration floor.}
An episode occupies at least the slot in which it starts, so before the
duration allocation distributes anything, every active child
$t \in \gA_c^{(k,\ell)}$ of \cref{eq:model_count_omega} receives a floor of
$X_t^{(k,\ell)}$ slot units, and the allocation splits only the remainder
$Y_c^{(k+1,\ell)} - \sum_{t \in \gA_c^{(k,\ell)}} X_t^{(k,\ell)}$.
The floor keeps $Y_t^{(k,\ell)} \ge X_t^{(k,\ell)}$ at every scale, which is
exactly the condition under which the next finer scale can apply the same
floor, and at the finest scale it is the condition under which every start
carries a nonempty episode. The coarsest-scale draw does not guarantee
$Y_t^{(K,\ell)} \ge X_t^{(K,\ell)}$ on its own, since the duration mean is
floored at one slot but the negative binomial draw around it is not floored at
the count. Where a parent block therefore holds fewer duration units than its
own count, the floor cannot be met, and the total is instead split across the
active children in proportion to their counts, which again preserves the block
total.

\begin{remark}
Both devices act inside one block and preserve its total. Reconciliation in
this model is a property of the generative mechanism rather than a constraint
imposed on its output, and neither the capacity cap nor the duration floor
weakens that property.
\end{remark}

\subsection{Event reconstruction as a discrete marked process}
\label{sec:event_reconstruction}

Under \cref{assm:atomic_resolution} the finest-scale arrays determine an event
log. Identify the finest epochs with the half-open wall-clock intervals
$I_t = [(t-1)\Delta_1,\, t\Delta_1)$, take the start set $\gB^{(\ell)}$ of
\cref{eq:model_episode}, and define the episode length in slots as
\begin{align}
L_t^{(\ell)} := \max\!\left( Y_t^{(1,\ell)}, \, 1 \right),
\qquad t \in \gB^{(\ell)} .
\label{eq:start_set_and_length}
\end{align}
The floor of one slot in \cref{eq:start_set_and_length} is the reconstruction
counterpart of the duration floor of \cref{sec:atomic_resolution}. A start with
zero booked duration units can arise when a parent epoch carries fewer units
than it has active children, and an episode of zero extent is not an event.

Each start induces an episode, written as a label with a start slot and an end
slot,
\begin{align}
e_t^{(\ell)}
=
\left(
\ell, \;
t, \;
f_t^{(\ell)}
\right),
\qquad
f_t^{(\ell)} = \min\!\left( t + L_t^{(\ell)} - 1, \; T_1 \right),
\qquad t \in \gB^{(\ell)},
\label{eq:episode}
\end{align}
where the end slot is inclusive and is clipped at the end of the record.
\Cref{eq:episode} is \cref{eq:model_episode} made exact at its two boundary
cases, a start that received no duration unit and an episode that would outrun
the record. The reconstructed trace of label $\ell$ is the sequence
$\gE^{(\ell)} = \langle e_t^{(\ell)} : t \in \gB^{(\ell)} \rangle$ in
increasing order of start slot, and the reconstruction map
\begin{align}
\gR :
\left\{ X_t^{(1,\ell)}, Y_t^{(1,\ell)} \right\}_{t,\ell}
\;\longmapsto\;
\left\{ \gE^{(\ell)} \right\}_{\ell \in \gL}
\label{eq:reconstruction_map}
\end{align}
is deterministic. Because its arguments are produced by exact top-down
allocation, the number of reconstructed episodes of a label under any coarser
epoch equals that epoch's count, and the reconstructed lengths sum to that
epoch's duration except where \cref{eq:start_set_and_length} lifts a zero to
one slot or \cref{eq:episode} clips at the horizon.

The traces induce the finest-scale occupancy process
\begin{align}
O_s^{(\ell)}
:=
\sum_{t \in \gB^{(\ell)}}
\mathbb{I}\!\left\{ t \le s \le f_t^{(\ell)} \right\},
\qquad s \in \gT^{(1)},
\label{eq:occupancy}
\end{align}
which counts how many episodes of label $\ell$ are active in slot $s$. Some
labels admit at most one active episode at a time, and for those the
reconstruction may return $O_s^{(\ell)} > 1$, since nothing in the allocation
prevents a long episode from running past the next start of the same label.
Collect such labels in $\gL_{\mathrm{ex}} \subseteq \gL$ and require
$O_s^{(\ell)} \le 1$ for every slot $s$ and every $\ell \in \gL_{\mathrm{ex}}$.

The requirement is enforced by a deterministic projection that truncates an
episode at the start of the next one. Write the start slots of label $\ell$ as
$b_1 < b_2 < \dots < b_{N_\ell}$, so that $N_\ell = |\gB^{(\ell)}|$. The
projected end slots are
\begin{align}
\hat{f}_{b_j}^{(\ell)}
&=
\max\!\left(
b_j, \;
\min\!\left( f_{b_j}^{(\ell)}, \; b_{j+1} - 1 \right)
\right),
\qquad j = 1,\dots,N_\ell-1,
\label{eq:exclusivity_projection}
\\
\hat{f}_{b_{N_\ell}}^{(\ell)}
&=
f_{b_{N_\ell}}^{(\ell)},
\label{eq:exclusivity_projection_last}
\end{align}
with the outer maximum keeping every projected episode at least one slot long
and the horizon clip inherited from \cref{eq:episode}.

\begin{remark}
\label{rem:projection_is_a_safeguard}
The projection changes no start slot and removes no episode, so the count of
episodes per label is preserved; it only shortens durations. It is the third
place, alongside the one-slot floor and the horizon clip, where the
reconstruction can depart from exact duration reconciliation, because the slots
it removes are not returned to a sibling. It is therefore a
safeguard rather than a mechanism, and the burden of avoiding overlap belongs
to the model, through the lagged structure, the context features and the
label-specific constraints. Evaluation reflects both of the things the
projection touches, scoring start slots by detection $F_1$ against the true
start set and extents by the intersection over union of the occupied slot sets
$\{ s \in \gT^{(1)} : O_s^{(\ell)} > 0 \}$ of \cref{eq:occupancy}.
\end{remark}

\subsection{Fundamental invariances}
\label{sec:invariances}

The invariances below are exact. Along each of them the joint distribution of
the observed counts and durations given the context is literally unchanged, so
no estimator, no sample size and no regularizer can distinguish the two
parameter values, and any evaluation that appears to distinguish them is
measuring its own coordinate system.

\begin{proposition}[Exact invariances of the parameterization]
\label{prop:invariances}
Fix a scale $k$ and let $\Theta'$ be obtained from $\Theta$ by any one of the
following transformations. Then $\Theta$ and $\Theta'$ induce the same
distribution of $(X,Y)$ given the context, that is, they are observationally
equivalent.
\begin{enumerate}
\item \textbf{Regime relabeling.} For a permutation $\sigma$ of
$\{1,\dots,M_k\}$, replace $S_t^{(k)}$ by $\sigma(S_t^{(k)})$ and permute every
block indexed by a scale-$k$ regime consistently, so that $a_m^{(k)} \mapsto
a_{\sigma^{-1}(m)}^{(k)}$, $A_{m,m'}^{(k)} \mapsto
A_{\sigma^{-1}(m),\sigma^{-1}(m')}^{(k)}$, and likewise for the context weights
$\vb_m^{(k)}$, the initial logits $\iota_m^{(k)}$, the factor blocks
$(\mF_m^{(k)}, \mG_m^{(k)}, \mSigma_m^{(k)})$ and the regime effects
$\vbeta_{S,\ell}^{(k)}$, $\valpha_{S,\ell}^{(k)}$,
$\boldsymbol{\delta}_{S,\ell}^{(k)}$. The cross-scale matrix $C^{(k)}$ within
the transition map of \cref{eq:model_regime_init} is indexed by a scale-$k$ regime in its first slot
and by a scale-$(k+1)$ regime in its second, so $\sigma$ permutes the first
index of $C^{(k)}$ and the second index of $C^{(k-1)}$, and the invariance
holds one scale at a time and, applied at every scale, jointly.

\item \textbf{Factor rotation and sign.} For any invertible
$\mR \in \sR^{r_k \times r_k}$, replace $H_t^{(k)}$ by $\mR H_t^{(k)}$, the
factor dynamics by $(\mR \mF_m^{(k)} \mR^{-1}, \, \mR \mG_m^{(k)}, \,
\mR \mSigma_m^{(k)} \mR^\top)$ for every $m$, the factor initialization by
$(\mR \vmu_0^{(k)}, \mR \mSigma_0^{(k)} \mR^\top)$, and every loading vector by
$\vbeta_{H,\ell}^{(k)} \mapsto \mR^{-\top} \vbeta_{H,\ell}^{(k)}$, and
similarly for $\valpha_{H,\ell}^{(k)}$ and $\boldsymbol{\delta}_{H,\ell}^{(k)}$.
Writing $\mB_H^{(k)} \in \sR^{r_k \times L}$ for the matrix whose $\ell$-th
column is $\vbeta_{H,\ell}^{(k)}$, only the row space of $\mB_H^{(k)}$ is
pinned down, and in particular a sign flip of a factor coordinate is invisible.

\item \textbf{Block-constant shifts of an allocation score, $k < K$.} For
constants $(v_\ell)_{\ell \in \gL}$, any one of the shifts
\begin{align*}
\beta_{0,\ell}^{(k)} \mapsto \beta_{0,\ell}^{(k)} + v_\ell,
\qquad
&\beta_{P,\ell}^{(k)} \mapsto \beta_{P,\ell}^{(k)} + v_\ell,
\qquad
\beta_{D,\ell}^{(k)} \mapsto \beta_{D,\ell}^{(k)} + v_\ell,
\\
&\beta_{S,m,\ell}^{(k)} \mapsto \beta_{S,m,\ell}^{(k)} + v_\ell
\;\;\text{for every } m,
\end{align*}
leaves every count allocation probability of \cref{eq:model_count_omega}
unchanged, and the corresponding shifts of $\alpha_{0,\ell}^{(k)}$,
$\alpha_{P,\ell}^{(k)}$ and $\alpha_{S,m,\ell}^{(k)}$ leave the duration
allocation produced by \cref{eq:model_count_omega} unchanged.

\item \textbf{Intercept trade at the coarsest scale.} For constants
$(v_\ell)_{\ell \in \gL}$, the simultaneous shift
\[
\beta_{0,\ell}^{(K)} \mapsto \beta_{0,\ell}^{(K)} + v_\ell,
\qquad
\beta_{S,m,\ell}^{(K)} \mapsto \beta_{S,m,\ell}^{(K)} - v_\ell
\;\;\text{for every } m,
\]
leaves the count mean unchanged, and the same trade in
$(\delta_{0,\ell}^{(K)}, \boldsymbol{\delta}_{S,\ell}^{(K)})$ and in
$(\alpha_{0,\ell}^{(K)}, \valpha_{S,\ell}^{(K)})$ leaves the zero-inflation
probability and the duration mean unchanged. Only the sums
$\beta_{0,\ell}^{(K)} + \beta_{S,m,\ell}^{(K)}$ are pinned down.

\item \textbf{Destination-free shifts of the regime logits.} For any constants
$v_0$, $(v_{m'})_{m'}$, $(v_n)_n$ and any vector $\vv \in \sR^{d_k}$, the
shifts
$a_m^{(k)} \mapsto a_m^{(k)} + v_0$,
$A_{m,m'}^{(k)} \mapsto A_{m,m'}^{(k)} + v_{m'}$,
$C_{m,n}^{(k)} \mapsto C_{m,n}^{(k)} + v_n$
and
$\vb_m^{(k)} \mapsto \vb_m^{(k)} + \vv$,
applied for every destination $m$, leave every regime transition probability
unchanged. Only the row-stochastic transition table at a given context is
pinned down.
\end{enumerate}
\end{proposition}

\begin{proof}
Item 1 is a relabeling of a summation index. The latent regime is summed out of
the observed likelihood, and permuting the label of the summand together with
every block it indexes leaves each term of the sum and therefore the sum
itself.

Item 2 holds because the factor enters every emission and every allocation
score only through the inner products $(\vbeta_{H,\ell}^{(k)})^\top H_t^{(k)}$
and their $\valpha$ and $\boldsymbol{\delta}$ analogues, and
$(\mR^{-\top}\vbeta_{H,\ell}^{(k)})^\top (\mR H_t^{(k)})
= (\vbeta_{H,\ell}^{(k)})^\top H_t^{(k)}$. The transformed dynamics are those
of the process $\mR H_t^{(k)}$, and a linear map of a Gaussian is Gaussian with
the stated covariance, so the joint law of the inner products is unchanged.
Rows of $\mR^{-\top}\mB_H^{(k)}$ are linear combinations of the rows of
$\mB_H^{(k)}$ and $\mR^{-\top}$ is invertible, so the row space is the
invariant object.

Item 3 holds because the allocation probability of a child is a softmax over
the children of one block, \cref{eq:model_count_omega}, and a softmax is
invariant to adding to every score of the block a quantity that does not depend
on the child. The intercept is such a quantity trivially. The parent terms
$v_\ell \log(1 + X_c^{(k+1,\ell)})$ and $v_\ell \log(1 + Y_c^{(k+1,\ell)})$
depend on the parent epoch $c$ but on no child of it, so they are constant
inside the block. A per-label constant added to the whole regime column
contributes $v_\ell$ for every child, whatever regime that child occupies,
because exactly one regime is active per epoch. The same argument applies
verbatim to the duration score over the active children $\gA_c^{(k,\ell)}$,
with $\alpha_{P,\ell}^{(k)}$ in place of the two parent terms. No other feature
of \cref{eq:model_count_score} is block-constant, since the
factor, the context, the depletion feature $W_t^{(k,\ell)}$, the own-epoch
count and the lagged values all differ between children of the same block.

Item 4 holds because the coarsest-scale linear predictors see the regime only
through $\beta_{0,\ell}^{(K)} + \beta_{S,m,\ell}^{(K)}$ evaluated at the active
regime $m = S_t^{(K)}$, and the trade leaves that sum, hence
$\mu_t^{(K,\ell)}$, hence the emission law; the argument for
$\varpi_t^{(K,\ell)}$ and $\nu_t^{(K,\ell)}$ is the same predictor written with
$\boldsymbol{\delta}$ and $\valpha$.

Item 5 holds because the transition probability of
\cref{eq:model_regime_init} is a softmax over the
destination regime $m$, and each listed shift adds to the logit vector a
quantity that is constant in $m$ for every fixed conditioning value of the
previous regime, the parent regime and the context. The shift of
$\vb_m^{(k)}$ is likewise constant in $m$ in the initial logits of
\cref{eq:model_regime_init}, so the initial distribution is unchanged as well.
\end{proof}

\begin{remark}
\label{rem:not_invariant}
The complement is as important. A lagged coefficient
$\beta_{\ell',d}^{(k,\ell)}$ is identified, because its regressor
$\log(1 + X_{t-d}^{(k,\ell')})$ varies across the children of a block; so is
the depletion coefficient $\beta_{\mathrm{dep},\ell}^{(k)}$, whose regressor
$\log(1 + W_t^{(k,\ell)})$ is built from what earlier children of the same
block have already taken and therefore varies inside the block by
construction; so is the count coupling $\alpha_{X,\ell}^{(k)}$, and so are the
dispersions $\phi_\ell^{(K)}$, $s_\ell^{(K)}$ and the concentrations
$\kappa_k$, $\tau_k$. The contrast of the regime effect across regimes is
identified even where its per-label level is not. These are the directions in
which a learner can be right or wrong.
\end{remark}

\subsection{Identified quantities and what a recovery metric may compare}
\label{sec:identified_quantities}

\Cref{prop:invariances} has a direct operational consequence. A parameter is
not a measurable quantity; a function of parameters that is constant on the
orbits of the invariances is. Reading off the proposition, the identified
objects per scale are the following.

\begin{itemize}
\item The row-stochastic regime transition table at a reference context, which
by item 5 is what the transition logits determine.
\item The row space of the factor loading block, by item 2.
\item At the coarsest scale, the array of sums
$\beta_{0,\ell}^{(K)} + \beta_{S,m,\ell}^{(K)}$ over regimes $m$ and labels
$\ell$, by item 4, together with the count dispersion $\phi_\ell^{(K)}$ and the
duration dispersion $s_\ell^{(K)}$.
\item At a finer scale, the regime effect after its per-label mean over regimes
is removed, by item 3. Equivalently, and more directly interpretable, the
shares that a block would give its children if those children sat in different
regimes, which is the softmax of the centered regime effect down the regime
axis.
\item The lagged coefficients on the true support, and the total absolute
lagged mass a learner places off that support.
\end{itemize}

The parameter-recovery metric is built from this list and reads nothing outside
it. It first aligns the regimes of the learned model with those of the truth by
the permutation that minimizes the squared distance between the identified
regime levels, so that a model differing only by a relabeling is aligned before
anything is scored. It then reports, per scale, the mean total variation
distance between the transition tables, the mean absolute error of the
identified emission blocks and that error divided by the spread of the true
block, one minus the mean best absolute cosine between true and learned loading
directions, the mean absolute difference of the implied allocation shares, and
the two lagged-coefficient summaries. The emission blocks are the identified
regime level at every scale and, at the coarsest scale, the count dispersion
alongside it. The absolute value in the cosine is the sign invariance of item 2
made explicit. The single headline number averages the scale-free entries
across scales.

\begin{remark}
The loading comparison uses the mean best absolute cosine rather than a
subspace distance. The two agree when $r_k = 1$ and when the transformation
relating the two loading blocks is a signed permutation; under a general
rotation the cosine statistic is a conservative summary of the row-space
agreement of item 2 rather than an exact one.
\end{remark}

The reason for this care is negative rather than positive. A metric that
compares raw parameter arrays entry by entry charges a learner for every step
it takes along a direction of \cref{prop:invariances}, which is a direction the
data cannot see. Two models that induce the same distribution of the observed
record then receive different scores, and the ranking they induce is partly a
ranking of coordinate conventions. This is the same class of failure that scale
artifacts introduce into structure benchmarks, where a criterion picks up a
property of the simulation rather than the property under study
\citep{reisach2021varsortability,reisach2023r2sortability}, and the same reason
structural comparison is done modulo an equivalence class rather than on raw
adjacency \citep{peters2015sid,runge2018pitfalls}.

The construction has an empirical anchor. In the parameter-recovery experiment
an oracle method returns the generating parameters without fitting anything,
and its headline error is zero to machine precision on all $27$ scenario and
replicate rows of \texttt{results/exp7/exp7\_results.csv}, the largest value
being $2.8 \times 10^{-17}$. The floor is exact by construction, which is what
makes the errors reported by the fitted methods above it interpretable as
error rather than as bookkeeping.

\subsection{Separation of explanatory roles}
\label{sec:role_separation}

The three mechanisms that carry temporal and cross-label dependence are the
latent regime, the latent factor and the lagged edges. They are not
interchangeable, and the model needs all three because each covers a pattern
the other two either cannot express or would express by distorting something
else.

The regime $S_t^{(k)}$ carries persistent drift. It is a discrete state with a
sticky transition table, so it changes rarely and, while it holds, shifts the
level of counts and durations for many labels at once. A record in which the
base rate of a label is high for a stretch of epochs and low afterwards is what
a regime is for.

The factor $H_t^{(k)}$ carries same-epoch co-movement. It is continuous,
low-dimensional and shared across labels within one epoch, so it explains why
two labels move together at the same time without asserting that either causes
the other. This is the mechanism that makes the model a dynamic Bayesian
network in the strict sense, because it lets all contemporaneous association
among observed variables be mediated by a latent quantity rather than by a
same-epoch directed edge.

The lagged parent sets $\gP_{k,\ell}$ carry directed dependence. An edge
$(\ell',d) \to \ell$ at scale $k$ states that the value of label $\ell'$ at lag
$d$ shifts the conditional law of label $\ell$ now, given everything else
admissible.

Dropping any one of the three does not simplify the model; it redirects the
variation into the remaining two. Without a regime, a persistent level shift is
absorbed by lagged self edges, since a persistent state is autocorrelated and a
lagged self edge is the cheapest way for a criterion that conditions on nothing
to express autocorrelation. This is the over-selection mechanism that the
results section quantifies. Without a factor, same-epoch co-movement has no
admissible carrier and is expressed as directed dependence at whatever lag the
sampling makes it look like. Without lagged edges, genuine directed dependence
is pushed onto a latent state, which then splits into more regimes than the
phenomenon has.

\begin{assumption}[Role separation]
\label{assm:role_separation}
Within a scale, parent-scale summaries explain the transfer of a coarse budget
to finer epochs, the lagged parent sets explain within-scale directed
dependence, the regime explains persistent drift in local dynamics, and the
factor explains same-epoch co-movement not attributable to lagged structure.
\end{assumption}

\Cref{assm:role_separation} is partly enforced by construction. Prohibiting
same-epoch directed edges among observed variables prevents the graph from
absorbing contemporaneous association; exact reconciliation prevents finer
scales from redefining coarser totals. The remainder is a modeling commitment,
and the ablation that removes each component in turn is what tests it.

\subsection{Modeling assumptions}
\label{sec:assumptions}

The following assumptions are what the interpretation of a fitted model rests
on. They are stated so that a reader can check them against a record before
reading a result off it.

\begin{assumption}[Correct multiscale aggregation]
\label{assm:aggregation}
The observed arrays at all scales are deterministic aggregations of one
finest-scale event log, so the reconciliation constraints hold in the data and
not only in the model.
\end{assumption}

\begin{assumption}[Consistent label semantics across scales]
\label{assm:label_semantics}
A label denotes the same phenomenon at every scale. If aggregating over a
coarser window changes what a label means, the hierarchy relates quantities
that are not comparable and the top-down construction has no referent.
\end{assumption}

\begin{assumption}[Admissible lag structure]
\label{assm:lag_structure}
All directed dependencies among observed variables are lagged and within a
scale. No directed edge joins two observed variables in the same epoch, and all
contemporaneous association is mediated by $H_t^{(k)}$ and $S_t^{(k)}$.
\end{assumption}

\begin{assumption}[Finite candidate parents and sparse truth]
\label{assm:candidates}
For each scale $k$ and label $\ell$ the candidate set
$\gC_{k,\ell} \subseteq \gL \times \{1,\dots,p_k\}$ is finite and contains the
true parent set, and the true parent set has bounded cardinality.
\end{assumption}

\begin{assumption}[Distinct regime signatures]
\label{assm:distinct_regimes}
At each scale, distinct regimes induce distinct conditional behavior in at
least one of the transition table, the factor dynamics, or the emissions of at
least one label. No two regimes are observationally identical beyond the
relabeling of item 1 of \cref{prop:invariances}.
\end{assumption}

\begin{assumption}[Factor normalization]
\label{assm:factor_normalization}
The innovation covariance is fixed at a reference regime, which cuts the
freedom of item 2 of \cref{prop:invariances} from the general linear group to
the orthogonal group. No location normalization is imposed, since
the continuous component of \cref{eq:model_regime_init} carries no intercept and the mean of the
factor is therefore already determined by $\vmu_0^{(k)}$, $\mF^{(k)}$,
$\mG^{(k)}$ and the context rather than free to be set. Recovery is claimed for
the loading row space and not for individual loadings.
\end{assumption}

\begin{assumption}[Allocation feature non-redundancy]
\label{assm:nonredundancy}
Within a child block, the features entering an allocation score are not
perfectly collinear. Beyond the block-constant shifts of item 3 of
\cref{prop:invariances}, distinct allocation parameters then induce distinct
allocation probabilities.
\end{assumption}

\begin{assumption}[Stationarity within a regime]
\label{assm:within_regime_stationarity}
Conditional on the regime, the context and the parent-scale summaries, local
dynamics are approximately stable. The model is non-stationary across regimes,
not within one.
\end{assumption}

\begin{assumption}[Regularity]
\label{assm:regularity}
Parameters lie in the interior of the parameter space and every conditional
distribution used is well defined and non-degenerate over the support of the
data.
\end{assumption}

Two remarks bound what these assumptions buy. First, exact global
identifiability of the graph is not available and is not claimed. Under
\cref{assm:candidates} the target is restricted identifiability inside the
admissible class, and an edge is recoverable when its inclusion changes the
conditional law of $(X_t^{(k,\ell)}, Y_t^{(k,\ell)})$ given all other
admissible predictors by a non-negligible amount over the support of the data.
Two nearly collinear candidate predictors are not separable by any criterion at
finite sample size, so a recovered graph is a statement about statistically
distinguishable dependencies relative to the candidate set and the penalty
$\lambda_k$, not about a unique mechanism. Second,
\cref{assm:lag_structure} makes the factor the carrier of contemporaneous
association, but a low-dimensional factor does not adjust for every unobserved
influence. Learned lagged edges are predictive temporal relations unless
stronger assumptions are added, which is the standing caveat for structure
learned from observational time series with possible latent confounding
\citep{gerhardus2020lpcmci}. Regime label switching and the resulting need to
align before comparing are shared with the non-stationary dynamic Bayesian
network literature more broadly
\citep{robinson2010nsdbn,grzegorczyk2011nhdbncontinuous}.

\subsection{Failure modes and safeguards}
\label{sec:failure_modes}

Each assumption above fails in a characteristic way, and each failure has a
signature that can be checked.

A finest scale that is too coarse violates \cref{assm:atomic_resolution}. Its
signature is finest-scale counts above one before the capacity cap acts, and
its consequence is that reconstruction stops being exact, since several
episodes of a label share a slot and their boundaries are no longer recoverable
from $(X^{(1)}, Y^{(1)})$.

Redundant candidate predictors violate \cref{assm:nonredundancy} and
\cref{assm:candidates}. Their signature is a structure score that is nearly
flat over a set of competing edges, and their consequence is that recovered
parent sets vary across replicates while predictive performance does not.

An over-flexible factor absorbs genuine lagged dependence. Raising $r_k$ gives
the factor enough freedom to reproduce, as same-epoch co-movement, a pattern
that a lagged edge would explain. The signature is a fall in the recall of the
true edges with little movement in the held-out likelihood.

Too many regimes split one behavior across several states. The signature is a
transition table with several rows that are close to each other and regimes
whose identified emission levels differ by little, which also makes the regime
alignment of \cref{sec:identified_quantities} unstable.

Misspecified allocation features confound top-down transfer with within-scale
dynamics. If the competition among the children of a block for a fixed parent
total is not represented by a feature that varies inside the block, the only
instrument the structure search has for it is a lagged self edge, and
within-block competition is then reported as structure.

The safeguards follow from the list. Keep the regime cardinalities $M_k$ and
factor dimensions $r_k$ small, and prefer the smallest values that a held-out
criterion supports. Bound the candidate lag sets and the in-degree. Penalize
graph and parameter complexity together. Screen a scenario for identifiability
before reading a metric from it, by checking that the likelihood gain of a true
edge exceeds both the gain of the best non-edge and the penalty an edge must
beat, and by comparing latent recovery against the oracle attainable with the
true parameters; \cref{sec:benchmark_identifiability} gives the screen used
here and the reading it returns. Ablate the components against each other, so
that parent-scale summaries, factor and lagged edges are each shown to carry
predictive value the others do not.

\begin{remark}
The model is a structured, partially identifiable multiscale latent-variable
model. Its parameters and its graph are meaningful relative to the admissible
class, the chosen resolutions, and the normalization and sparsity assumptions
above, and any claim read from a fitted model inherits those qualifications.
\end{remark}

\section{Benchmark design and experimental protocol}
\label{sec:benchmark}

The claims this model makes are recovery claims. It claims to find the lagged
parent sets \(\gP_{k,\ell}\) that generated a multiscale log, to track the
latent regime \(S_t^{(k)}\) and the latent factor \(H_t^{(k)}\) that drive it,
and to do both while the counts and durations reconcile exactly across scales.
None of these can be checked against a quantity the data supply, so the
evaluation has to supply the quantity itself. This section states what the
generated benchmark contains, how a scenario is configured, and, before any of
that is used, how a scenario is screened for whether it contains anything to
recover at all.

\subsection{Why recovery is not measurable on an observed log}
\label{sec:benchmark_why}

On an observed event log the true parent sets, the true regime path and the
true parameters \(\Theta^\star\) are all unavailable, and no amount of care in
the fitting procedure produces them. A precision against \(\gP_{k,\ell}^\star\)
is not a hard quantity to compute on such a log; it is an undefined one. What
an observed log does support is prediction on held-out time, exact
reconciliation, the stability of a selected edge across subsamples, and
falsification tests that need no ground truth, such as checking whether graphs
learned from different variable subsets agree. Those are necessary conditions
and none of them is a recovery measurement.

A generator removes the obstruction by construction. It yields the graph, the
regime path, the factor trajectory and every conditional distribution exactly,
so precision, recall, structural Hamming distance, regime accuracy and
parameter error all become computable. It also yields something an observed log
cannot give at any price, which is control over the difficulty of the problem.
The generator can be interrogated about a scenario before an estimator is
pointed at it, and \cref{sec:benchmark_identifiability} makes that
interrogation part of the protocol rather than an afterthought.

The cost is equally clear. A recovery number measured on data drawn from the
model being evaluated establishes self-consistency and nothing wider. Two
design decisions bound that cost. The first is a family of scenarios generated
under deliberate violations of the modeling assumptions, so that graceful
degradation rather than exact recovery is what is measured there. The second is
running the identical evaluation code against networks released by others,
where the ground truth was fixed by someone with no stake in this model.

\subsection{The four generating families}
\label{sec:benchmark_families}

Four families span the axes the model is built around, and each one isolates a
different component by removing the others.

\paragraph{Family I, stationary and sparse.}
Every scale carries a single regime and no factor, \(M_k = 1\) and \(r_k = 0\),
with a fixed sparse lagged graph and parameters constant in time. What remains
is structure recovery and the top-down allocation, with no latent dynamics to
confound either.

\paragraph{Family II, regime switching.}
At least two regimes act at the coarsest scale and \(r_k = 0\) everywhere, with
regime-specific count levels, duration levels and allocation scores, and a
transition matrix built so that the regimes persist and occupy their states
about equally. Which further scales carry a regime is not a free choice; it is
whatever the screen of \cref{sec:benchmark_identifiability} admits. This family
tests recovery of non-stationary behavior and is the one in which a persistent
latent state is available to confound a latent-free structure criterion.

\paragraph{Family III, shared latent factor.}
A single regime and a factor of dimension at least one, \(M_k = 1\) and
\(r_k \ge 1\), with a stable first-order factor dynamic per scale and loadings
that put correlated same-slice fluctuations on every label. The lagged graph
stays sparse, so the question is whether the factor absorbs contemporaneous
co-movement without distorting recovery of the lagged structure.

\paragraph{Family IV, deliberate misspecification.}
The same structure with one generating assumption broken on purpose, in two
variants. The heavier-tail variant fixes the coarsest-scale negative binomial
dispersion at a large constant, which drives the generated counts toward the
Poisson limit and leaves the fitted dispersion chasing a value the estimator
cannot reach. The omitted-context variant enlarges the coarsest-scale context
coefficients so that the context genuinely drives the emission level, and the
fitted model has to absorb its effect elsewhere. A third setting names the
reference scenario, the same structure with nothing misspecified, which is what
makes the degradation measurable rather than merely observed.

\paragraph{A family overrides the scenario it is given.}
The family is not a label attached to a configuration; it is a projection of
that configuration onto its own subfamily, applied before any parameter is
drawn. Family I forces \(M_k = 1\) and \(r_k = 0\) at every scale. Family II
forces \(r_k = 0\), raises the coarsest scale to \(\max(2, M_K)\), and leaves
every other scale as the configuration states, so a scale asked to carry no
regime carries none. Family III forces
\(M_k = 1\) and \(r_k \gets \max(1, r_k)\). Asking family I for four regimes
therefore returns one, and asking family II for a factor returns none. The
consequence for the grid is that the regime count is a live axis only inside
family II and the factor dimension only inside family III. A grid that swept
\(M_k\) and \(r_k\) across all families would have generated many scenarios
that are byte-identical to one another and reported them as independent
evidence.

\subsection{Scenario configuration and the two grids}
\label{sec:benchmark_grid}

A scenario fixes the tuple
\begin{align}
\bigl(
K,\;
T_1,\;
(b_k)_{k=1}^{K-1},\;
\gL,\;
(M_k),\;
(r_k),\;
(p_k),\;
(D_k^\star),\;
\varpi_0,\;
\text{family},\;
\text{seed}
\bigr),
\label{eq:bench_scenario}
\end{align}
where \(b_k\) is the branching factor from scale \(k\) to scale \(k+1\),
\(p_k\) is the maximum lag searched at scale \(k\) in
\cref{eq:model_lag_values}, \(D_k^\star\) is the maximum
in-degree of the true graph there, and \(\varpi_0\) is the structural zero
probability that fixes the zero-inflation intercept, set to \(0.1\), \(0.35\)
or \(0.65\) for the low, medium and high zero-inflation levels. The epoch
counts follow from the branching,
\begin{align}
T_{k+1} = \left\lfloor T_k / b_k \right\rfloor,
\qquad k = 1,\dots,K-1,
\label{eq:bench_epochs}
\end{align}
and the wall-clock length of an epoch follows from the length of one finest
slot together with the same branching.

Two hierarchies are used, and both are aligned to a clock rather than chosen
for arithmetic convenience. The minute grid sets \(K = 3\), \(T_1 = 20{,}160\)
and branching \((10, 6)\) at a one-minute finest slot, giving \(20{,}160\),
\(2{,}016\) and \(336\) epochs at one minute, ten minutes and one hour over
fourteen days. The day grid sets \(K = 3\), \(T_1 = 8{,}640\) and branching
\((12, 24)\) at a five-minute finest slot, giving \(8{,}640\), \(720\) and
\(30\) epochs at five minutes, one hour and one day over thirty days. The
second exists so that the study repeats when the coarsest scale is a calendar
day, which changes the number of coarsest epochs available for estimation by
an order of magnitude.

Both grids fix the same generating conditions. The finest scale is constrained
to binary event starts, \(X_t^{(1,\ell)} \in \{0,1\}\), which is a modeling
choice about the labels rather than a requirement of the event-log reading. What
makes a reconstructed trace an event log is that every episode carries a start
and an end **at the resolution of the finest scale itself**, and nothing finer:
if that scale is an hour, an episode runs from one hour to another and may begin
and end in the same one. Two episodes of a label beginning in the same finest
epoch are therefore admissible, since the log does not claim to order events
inside an epoch, and no claim below the declared finest scale is supported by
anything the generator produced. The
coarsest log-mean is set from a target rate of \(0.15\) events per finest epoch
conditional on the label not being structurally absent, since fixing the
coarsest mean directly empties the finest scale once the counts have been
spread over sixty slots. The lagged coefficients and the regime contrasts carry
multipliers of \(10\) and \(3\) on their base draws, which is what places them
in the regime \cref{sec:benchmark_identifiability} certifies. The searched lag
bound is \((p_k) = (3,2,1)\) from finest to coarsest and the true in-degree
bound is \((D_k^\star) = (0,2,1)\), so the true edges live at the ten-minute
and the hourly scales and the finest scale is left empty. The latent states are
placed by the same rule and land differently: the regime count is
\((M_k) = (1,1,M)\) on the minute grid and \((1,M,M)\) on the day grid, since
the screen of \cref{sec:benchmark_identifiability} admits a regime at the hour
scale of the day grid and at no scale below the coarsest of the minute grid,
while the factor dimension \(r_k\) is uniform across scales because the factor
clears everywhere.

The minute grid holds twenty scenarios. Family I contributes four, from
\(L \in \{5,10\}\) crossed with low and high zero inflation. Family II
contributes eight, from \(L\) crossed with \(M \in \{2,4\}\) at the coarsest
scale and one regime below it, and the two zero inflation levels. Family III contributes eight, from
\(L\) crossed with \(r_k \in \{1,2\}\) and the two levels. The day grid adds
two at medium zero inflation, one from family II with two regimes at the hour
and day scales and one from family III with \(r_k = 1\), for twenty-two
scenarios in total, each generated
at three independent replicates.

Family IV sits outside this grid and carries its own. Degradation under
misspecification is a paired comparison against an otherwise identical
well-specified scenario, so the misspecification type has to be the swept axis
rather than one setting among twenty-two. That grid holds twelve
configurations, from the three misspecification settings crossed with
\(L \in \{5,10\}\) and \(M \in \{2,4\}\) at the coarsest scale and medium zero
inflation, on the minute clock and with the same lag bound, in-degree bound,
regime placement, rate, edge-strength and regime-strength calibration as above,
again at three replicates.

One asymmetry in the search space is deliberate. The in-degree bound the
learner searches under is \(D_k = \max(D_k^\star, 1)\) rather than
\(D_k^\star\) itself, so the finest scale, which carries no true edge, is
searched with a budget of one parent per label instead of zero. A budget of
zero would make that scale trivially correct for every method and would destroy
the precision test it exists to carry.

Each replicate is split chronologically at coarsest-scale boundaries into
training, validation and test partitions holding \(70\%\), \(15\%\) and
\(15\%\) of the coarsest epochs. Splitting on coarsest-scale boundaries is what
keeps the reconciliation identity
\(\sum_{t \in \child_k(c)} X_t^{(k,\ell)} = X_c^{(k+1,\ell)}\) true inside each
partition; a split at an arbitrary finest-scale index would cut a block in half
and break it. The split is chronological rather than random because the model
is sequential and the question is generalization forward in time. Every method
receives the same training partition, the same initialization and the same
search space, namely the scenario's own \((p_k, D_k)\), since a search space
narrower than the generating one caps recall before a single score is computed.

\subsection{The identifiability screen}
\label{sec:benchmark_identifiability}

A recovery metric carries information only when the generated data contain what
is being recovered. A scenario in which a true edge leaves less evidence than
the best non-edge produces a low \(F_1\) for every consistent estimator, and
reporting that number as a property of an estimator is a category error. The
screen described here is applied to every scenario before it enters a study,
and it decided the shape of the grid in \cref{sec:benchmark_grid}.

The screen fits nothing. It generates one replicate, runs latent-state
inference at the true parameters \(\Theta^\star\) and the true graph
\(\gG^\star\), and asks what a perfect estimator would see. Write \(q\) for the
posterior that inference returns under \(\Theta^\star\) and \(\gG^\star\), and
define the unpenalized gain of a single candidate lagged parent
\(e = (\ell', d)\) for label \(\ell\) at scale \(k\) as
\begin{align}
\Delta_{k,\ell}(e)
&=
\max_{\beta,\alpha}\;
\E_q\!\left[
\log p\bigl(X^{(k,\ell)}_{1:T_k}, Y^{(k,\ell)}_{1:T_k}
\bigm| S^{(k)}, H^{(k)}, \gK^{(k)}, X^{(k+1,\ell)}, Y^{(k+1,\ell)},
\{e\}, \beta, \alpha, \Theta_k^\star\bigr)
\right]
\nonumber \\
&\qquad -
\E_q\!\left[
\log p\bigl(X^{(k,\ell)}_{1:T_k}, Y^{(k,\ell)}_{1:T_k}
\bigm| S^{(k)}, H^{(k)}, \gK^{(k)}, X^{(k+1,\ell)}, Y^{(k+1,\ell)},
\emptyset, \Theta_k^\star\bigr)
\right],
\label{eq:probe_gain}
\end{align}
with the coefficients of the candidate profiled out, exactly as the structure
step scores it. The parent totals \(X^{(k+1,\ell)}\) and \(Y^{(k+1,\ell)}\) are
in the conditioning set at a finer scale, where the likelihood is the
Dirichlet-multinomial allocation of those totals over \(\child_k(c)\), and are
absent at \(k = K\), where it is the zero-inflated negative binomial count and
negative binomial duration pair. The expectation is likewise exact at \(k = K\)
and approximate below it, since a finer-scale allocation score is evaluated at
its regime average \(\bar{\eta}_t^{(k,\ell)}\) before the likelihood is formed,
which is the substitution the structure step itself makes.

Three quantities are then read off per scale, the median of
\(\Delta_{k,\ell}\) over the true edges, the \(95\)th percentile of
\(\Delta_{k,\ell}\) over the candidates that are not true edges, and the price
\(\log T_k\) that a plain Bayesian information criterion attaches to one added
parent. That price is a full log-epoch rather than half of one because a
parent adds two coefficients, one on the count channel and one on the duration
channel, at \(\tfrac{1}{2}\log T_k\) each. The screen also reports the number
of epochs of that scale in the training split, the fraction of label-epoch
cells there carrying at least one event, and the best-permutation regime
accuracy that the true parameters themselves attain.

The scenario is worth running a structure study on only when the three gains
order correctly,
\begin{align}
\operatorname{med}_e \Delta_{k,\ell}(e)
\;\gg\;
Q_{0.95}\bigl(\Delta_{k,\ell} \mid e \notin \gP_{k,\ell}^\star\bigr)
\;>\;
\log T_k ,
\label{eq:probe_ordering}
\end{align}
and the latent metrics are interpretable only relative to the oracle regime
accuracy. \Cref{tab:probe} gives the reading for the benchmark scenario, a
family II configuration on the minute grid with five labels, two regimes and
low zero inflation.

\begin{table}[t]
\centering
\caption{Identifiability screen on the benchmark scenario, one replicate,
latent-state inference at the true parameters and the true graph. The minute
scale carries no true edge, so its true-edge gain is undefined and the entry is
left blank. Gains are in nats.}
\label{tab:probe}
\begin{tabular}{llrrrrrr}
\toprule
\(k\) & epoch & \(T_k\) (train) & density & true edges & true gain & non-edge \(Q_{0.95}\) & \(\log T_k\) \\
\midrule
1 & one minute  & 14{,}100 & 0.255 & 0 &        &  32.82 & 9.55 \\
2 & ten minutes &  1{,}410 & 0.530 & 5 &  31.50 &  11.32 & 7.25 \\
3 & one hour    &      235 & 0.856 & 4 & 182.26 & 141.91 & 5.46 \\
\bottomrule
\end{tabular}
\end{table}

The reading of \cref{tab:probe} is the design of the grid. At the ten-minute
scale a true edge is worth \(31.50\) nats against a non-edge level of \(11.32\)
and a penalty of \(7.25\), and at the hourly scale it is worth \(182.26\)
against \(141.91\) and \(5.46\), so \eqref{eq:probe_ordering} holds at both and
the edges there are recoverable in principle. The whole gain scale at the hour
sits higher, non-edges included, because an edge there is scored against the
count and duration likelihood itself rather than against an allocation of a
parent total over children. At the minute scale the ordering fails before any
estimator is involved, since a non-edge already reaches \(32.82\) nats against a
penalty of \(9.55\), so a criterion charging the plain price would accept
non-edges on evidence alone. Placing true edges there would produce recovery
numbers that measure the scenario rather than the method.

That reading is one scenario, and the placement of the edges follows from the
same screen applied to all twenty-two. Each was generated once more at a nonzero
in-degree on every scale, and the quantity recorded was the median true-edge
gain less the floor it has to clear, \(\max(Q_{0.95}, \log T_k)\). The
ten-minute scale returns \(-4.42\), \(+7.68\) and \(+15.92\) nats for families I,
II and III with \(17\) of \(22\) scenarios clearing, and the hourly scale
\(-9.89\), \(+6.29\) and \(+11.15\) with \(14\) of \(22\). The finest scale fails
in every family, at \(-6.88\), \(-16.50\) and \(-8.64\), with only \(2\) of
\(22\) clearing. Family I clears at neither of the two coarser scales on this
draw, and its four scenarios split by zero inflation rather than by anything the
screen controls: at the hourly scale the two low-inflation scenarios return
\(-37.09\) and \(-79.24\) while the two high-inflation ones return \(+17.31\) and
\(+22.21\). A denser count series gives a non-edge more to exploit there, and the signal is
not what separates those scenarios: their realized effect sizes are $0.3244$ and
$0.3407$, so the failing arm is if anything marginally the stronger one. The
coarsest coefficients are pinned to one value, $0.3$, by the bound that keeps the
count autoregression stable, and that makes the edge-strength multiplier inert
at that scale; what
it does not make constant is the difficulty, since the evidence an edge leaves is
that coefficient times the spread of the column it multiplies, which runs
$0.2655$ to $0.7907$ across the grid. The grid therefore carries true edges at the
ten-minute and the hourly scales and leaves the finest one empty.

That failure is a property of the generator rather than an accident of one draw.
The finest scale is a binary event-start process, so a lagged edge there can
only shift which slot of a block begins an episode, which is worth less than the
price of the edge and is small relative to the freedom the allocation already
has. Leaving the scale empty converts it from a failed recovery test into a
working precision test. Any edge a method reports there is structure the data do
not contain, which makes the count of such edges a direct, interpretable
false-positive measurement with no threshold to choose. The converse has to be
stated with it, so that silence there is not read as recovery. No measurement on
this grid can credit a method for finding a finest-scale edge, because none can
be placed there.

An earlier version of this grid left the coarsest scale empty as well. The
screen that justified it had been read on scenarios generated with no
coarsest-scale edge at all, where every gain it could observe there was a
non-edge gain and the whole scale of them sat below the price. Regenerating at a
nonzero in-degree and rescreening reverses that reading, and the same screen is
what keeps the finest scale out.

The latent states need the same screen, and they need it stated with a floor
rather than a ceiling alone. The ceiling is the best-permutation regime accuracy
the true parameters themselves attain. An earlier version of this section
reported that ceiling as \(0.652\) at the minute scale, \(0.661\) at ten minutes
and \(1.000\) at the hour, and concluded that a fit near either of the first two
was at the limit the data allow. That conclusion does not follow, because a
ceiling bounds nothing on its own. The generated chain spent \(0.71\) of its
epochs in one state at the minute scale, so naming that state scored \(0.71\)
and the ceiling sat below the constant predictor. A fit at that ceiling was not
a successful inference reported as a failure; it was no inference at all, and
every parameter block driven by that path was scored against noise.

The screen therefore reports both, and refuses a latent state whose ceiling does
not clear its floor. A regime is scored against the share of its modal state,
which is what a constant predictor reaches. A factor has no constant predictor,
so it is scored against the correlation a circularly shifted copy of its own
trajectory reaches, which is the null the edge confirmation already uses. Write
\(\pi^\star_{\max}\) for the modal share of the generating path at scale \(k\)
and \(a_k^\star\) for the accuracy attained under \(\Theta^\star\) and
\(\gG^\star\); the state is admissible at that scale only when
\begin{align}
a_k^\star - \pi^\star_{\max} \;\ge\; \delta ,
\label{eq:latent_ordering}
\end{align}
with \(\delta = 0.05\), so that a path beating the constant by a thousandth is
not counted as recovered.

Applied to the nine family II scenarios, \eqref{eq:latent_ordering} holds at the
coarsest scale of both grids in \(9\) of \(9\), by \(+0.468\) to \(+0.600\). It
fails at the minute scale in \(0\) of \(8\), with every margin inside \(0.021\)
of the floor, and holds at the ten-minute scale in \(1\) of \(8\), which does not
survive raising the regime contrast and is therefore a draw rather than a signal.
The hour scale of the daily grid is the exception and keeps its regime, clearing
four times out of four at \(+0.157\) and above; its blocks hold \(24\) children
against the minute grid's \(6\), and that is the channel a regime has to act
through once it no longer reaches the count emission. The contrast is not the
lever, which had to be measured before the scale could be blamed: raising it
from \(3\) to \(8\) to \(20\) moves the two finer scales from \(0\) and \(2\) of
\(9\) clearing to \(1\) and \(2\) of \(9\), while the coarsest clears \(9\) of
\(9\) throughout. The grid therefore carries \(M_k = [1, 1, M]\) on the minute
hierarchy and \([1, M, M]\) on the daily one, decided exactly as the in-degree
is.

The generating chain is built so that this comparison is worth making. Its
transition logits used to be drawn freely, so the stationary distribution was
whatever the draw implied and the modal share at the coarsest scale was
\(0.672\) at two regimes. Each row is now \(q e_m + (1 - q)\pi\), a self
transition mixed with a draw from \(\pi\), which leaves \(\pi\) stationary
exactly for any persistence \(q\), with \(\pi\) drawn within ten percent of an
equal share. The modal share fell to \(0.511\). That rebalance changed no verdict
in \eqref{eq:latent_ordering}, since the finer scales fail on the ceiling sitting
at the floor rather than on the floor being high, but it is what makes a reported
accuracy at the coarsest scale readable.

The factor clears at every scale of every scenario, by \(+0.34\) to \(+0.78\) in
correlation, including at the minute scale where no lagged edge can be placed and
no regime can be recovered. Same-slice co-movement survives the event-start
reading and persistent drift does not, which is why \(r_k\) is uniform across
scales and \(M_k\) is not. The same argument applies to every parameter block
whose identified form is coarser than its stored form.

Stated generally, the screen turns a claim about an estimator into a claim
about an estimator on a scenario that has been certified to contain the signal.
It costs one replicate and one pass of latent-state inference per scenario,
which is negligible against the cost of the studies it protects, and it is the
step that makes a low recovery number attributable.

\subsection{Methods under comparison}
\label{sec:benchmark_methods}

Seven entries share one interface, so that the protocol fits them on identical
data and the difference between two rows is the single design decision that
separates them.

The \emph{full model} runs all three learning steps at every scale. The
\emph{single-scale variant} switches off the cross-scale coupling, so no
parent-scale posterior reaches latent-state inference or parameter learning and
the regime chain at each scale sees only its own data and context. The
allocation still ties a finer scale to its parent totals, since that is what
makes the counts reconcile, so this entry isolates the value of passing the
parent posterior downward rather than the value of the hierarchy itself. The
\emph{no-factor} ablation forces \(r_k = 0\) at every scale, which sends all
same-slice co-movement into the regimes, the lagged edges or the noise. The
\emph{no-regime} ablation collapses the chain to \(M_k = 1\), which removes
non-stationarity while leaving the hierarchy and the graph intact. The
\emph{empty graph} fixes \(\gP_{k,\ell} = \emptyset\) at every scale with the
structure step switched off, leaving top-down transfer and the latent model as
the only mechanisms; it is the reference against which a structural Hamming
distance becomes meaningful, because a method whose distance exceeds it is
worse than reporting nothing. The \emph{oracle graph} installs
\(\gG^\star\) and switches the structure step off, which separates parameter
and latent-state recovery from search error.

The seventh entry, \emph{oracle parameters}, fits nothing at all. It runs
latent-state inference at \(\Theta^\star\) and \(\gG^\star\) and reports the
resulting metrics. It is a ceiling and not a competitor, and comparing a fitted
model against it answers a question no other row can, which is whether a
disappointing latent metric reflects a learner that failed or a scenario that
does not identify the quantity. The screen of
\cref{sec:benchmark_identifiability} already fixes that ceiling numerically for
the regime path, and the oracle row extends it to every metric the protocol
computes.

When a published structure learner supplies the graph, it enters through the
same interface, with the learned edges replacing \(\gG\) and nothing else
changed. One implementation detail decides whether that comparison is valid at
all. The parameter step updates the coefficient of a lagged parent only where a
coefficient for that parent already exists, so a graph installed from outside
the structure step has no effect on any likelihood until a zero coefficient is
created for each of its parents. Without that step every externally supplied
graph is scored exactly as if it had reported the empty graph, and a table of
apparently distinct methods measures nothing.

\subsection{Metrics}
\label{sec:benchmark_metrics}

\paragraph{Structure recovery, per scale.}
Write
\(E_k = \{(\ell', d, \ell) : \ell \in \gL,\; (\ell', d) \in \gP_{k,\ell}\}\)
for the lagged edge set that a collection of parent sets at scale \(k\)
defines, and
\(\hat{E}_k\) and \(E_k^\star\) for its learned and true values. The two are
compared as sets of triples, giving true positives, false positives and false
negatives and hence
\begin{align}
\mathrm{Prec}_k = \frac{|\hat{E}_k \cap E_k^\star|}{|\hat{E}_k|},
\qquad
\mathrm{Rec}_k = \frac{|\hat{E}_k \cap E_k^\star|}{|E_k^\star|},
\qquad
\mathrm{SHD}_k = \bigl|\hat{E}_k \,\triangle\, E_k^\star\bigr| ,
\label{eq:bench_structure}
\end{align}
with \(F_{1,k}\) the harmonic mean of the first two. Reporting per scale rather
than pooled is forced by the design, since one of the three scales carries no
true edges and pooling would let recall at another scale hide false positives
there. The empty-versus-empty case is scored by convention rather than by
formula. When \(\hat{E}_k\) and \(E_k^\star\) are both empty the scale is
perfectly recovered, so precision, recall and \(F_1\) are all set to \(1\);
treating them as \(0/0\) and recording zero drags the macro average down by a
whole scale in every scenario, since the finest scale is empty in all of them.
The convention has to be declared wherever the macro average over scales is
quoted, because it hands a free \(1\) at the finest scale to any method that
proposes nothing there, and a method that proposes nothing anywhere collects
that term as well.

\paragraph{Latent-state recovery.}
The regime label is identifiable only up to permutation, so the estimate
\(\hat{S}_t^{(k)} = \argmax_m \gamma_t(m)\) is scored under the relabeling that
suits it best,
\begin{align}
\mathrm{Acc}_S^{(k)}
=
\max_{\sigma \in \mathfrak{S}_{M_k}}
\frac{1}{T_k}\sum_{t=1}^{T_k}
\mathbb{I}\bigl\{\sigma(\hat{S}_t^{(k)}) = S_t^{(k)\star}\bigr\},
\label{eq:bench_accs}
\end{align}
enumerated exactly for \(M_k \le 8\) and matched greedily on the contingency
table above that. Accuracy alone rewards a degenerate estimate on an unbalanced
regime path, so the adjusted Rand index and the normalized mutual information
are reported alongside it; both are invariant to relabeling by construction and
the adjusted Rand index is corrected for chance agreement.

Two readings of \eqref{eq:bench_accs} that a table would otherwise get wrong. A
scale whose true path never switches has no regime to recover, and
\eqref{eq:bench_accs} returns exactly \(1\) there for any estimate whatsoever, so
all three figures are reported as not applicable rather than as a perfect score.
Since the grid places a regime only where the screen of
\cref{sec:benchmark_identifiability} admits one, that is the usual case at the
two finer scales of the minute hierarchy. A scale whose true path switches while
the fitted model emits a constant is scored, and the score is worth reading: the
only available relabeling maps that constant onto one true regime, so
\eqref{eq:bench_accs} returns that regime's frequency, which is what a constant
predictor earns and is a floor rather than a recovery. Regime accuracy from a
model with \(M_k = 1\) is therefore not comparable with regime accuracy from one
with \(M_k > 1\). The factor is
identified only up to a rotation, so pointwise recovery of \(H_t^{(k)}\) is not
the quantity of interest. The metric averages, over the inferred components at
a scale, the largest absolute correlation that component attains against any
true component, which is invariant to the sign and to the ordering of the
components but not to a general rotation of them.

\paragraph{Aggregation consistency.}
The model claims exact reconciliation, so the claim is measured rather than
assumed,
\begin{align}
\mathrm{RecErr}_X
=
\sum_{k=1}^{K-1}\sum_{\ell \in \gL}\sum_{c \in \gT^{(k+1)}}
\Bigl| X_c^{(k+1,\ell)} - \sum_{t \in \child_k(c)} X_t^{(k,\ell)} \Bigr|,
\label{eq:bench_recx}
\end{align}
and identically for \(\mathrm{RecErr}_Y\) with \(Y\) in place of \(X\). The
quantity is a property of whatever multiscale array it is handed. In the
protocol that array is the held-out partition of the generated data, which
satisfies the identity by construction because the finer scales were produced
by allocating the parent total, so every entry scores exactly zero and the
column certifies the generator and the split rather than separating methods.
It becomes a comparison only when it is applied to an array a method
reconstructed itself, where a scale predicted without reference to its parent
leaves a residual whose magnitude is hierarchical inconsistency in the units of
the data.

\paragraph{Held-out log evidence, and how it must be read.}
The predictive quantity is the negative log evidence on the validation
partition, \(-\sum_k \log Z_k\), where \(\log Z_k\) is what latent-state
inference returns at scale \(k\); the recovery metrics are computed on the test
partition instead. It sums a zero-inflated negative binomial count likelihood
and a negative binomial duration likelihood at the coarsest scale with a
Dirichlet-multinomial allocation likelihood at each finer scale, the allocation
terms up to an additive constant that depends only on the data, since the
block-level normalizer and the multinomial coefficient shift every candidate by
the same amount and are dropped. That constant is the reason the quantity
compares methods on one dataset and does not compare datasets. Its
magnitude also grows with \(T_k\) and with the count level, so an average of
raw values across scenarios is dominated by whichever scenario happens to be
largest. The defensible summary is the paired per-fit difference between two
methods on the same replicate, reported as the fraction of fits in which one
beats the other together with the mean difference.

\paragraph{Parameter recovery.}
Structure recovery reports which edges returned and is silent on whether the
conditional distributions attached to them are right. The parameter metric
scores the learned blocks directly against the generating ones, and it scores
only the identified form of each block. Two alignments are needed. The regime
permutation is resolved once per scale by matching identified emission levels,
and every regime-indexed block is then compared under that one relabeling. The
factor loading is identified only up to a rotation, since rotating the factors
rotates the loadings back, so the loadings are compared through their row
space, by one minus the mean best absolute cosine between the normalized true
and learned loading vectors.

Which form is identified differs by scale, and comparing raw arrays instead
charges for directions the data cannot see. At the coarsest scale the count
mean is
\(\exp\bigl(\beta_{0,\ell}^{(K)} + \beta_{S,\ell}^{(K)}[m] + \cdots\bigr)\), so
only the sum \(\beta_{0,\ell}^{(K)} + \vbeta_{S,\ell}^{(K)}\) is pinned down and
a constant may move between the two freely. At a finer scale the allocation is
a softmax over the children of one block, so every term constant inside that
block cancels exactly and only the contrast of the regime effect across regimes
survives,
\begin{align}
\text{identified level of label } \ell \text{ at scale } k < K
\;=\;
\vbeta_{S,\ell}^{(k)}
- \frac{1}{M_k}\sum_{m=1}^{M_k} \beta_{S,\ell}^{(k)}[m].
\label{eq:bench_identified}
\end{align}
The reported blocks are then the total variation distance between transition
tables, the mean absolute lag coefficient error on the true support, the
absolute learned lag mass placed off that support, the emission level error
divided by the spread of the true block, the loading subspace error and the
mean absolute difference of the child shares a regime implies for the
allocation. The scale-free entries are averaged into a single figure, and the
headline is the sample-size ladder rather than any single value, since
parameter recovery is a consistency claim and an error that does not fall with
\(T_1\) is a bias rather than an estimate.

\paragraph{Edge AUPRC.}
The candidate edge sets here are sparse, so the negative class dominates and a
single thresholded \(F_1\) confounds the quality of a ranking with the choice
of operating point. Where a method returns scored candidates rather than a
hard set, the candidates are ranked by score descending and the precision at
each recall gain is summed and divided by the number of true edges. Equally
scored candidates form one indistinguishable block and are all credited the
precision at the end of that block, so the value does not depend on the order
in which the scores happen to be enumerated, and a true edge that was never
scored contributes nothing while still counting in the denominator. The
quantity is undefined when a scale has no true edges, which is the honest
outcome rather than a value of zero.

\subsection{Protocol items that decide whether the evaluation survives review}
\label{sec:benchmark_review}

Four reporting decisions separate a structure-learning evaluation that can be
compared against published work from one that cannot. They are listed here
because each has a documented failure mode and because each is cheap to satisfy
in advance and impossible to repair afterwards.

\paragraph{Standardize, then report both sortability statistics.}
A synthetic generator that draws coefficients on a fixed positive range induces
a marginal variance ordering that agrees with the causal order, and a
continuous-optimization learner then scores well by sorting on variance rather
than by discovering anything \citep{reisach2021varsortability}. The design here
standardizes the response and the lagged design matrix before the
continuous-optimization entry is solved, which removes the raw-scale advantage.
Standardizing is necessary and not sufficient, since an \(R^2\)-based
sortability survives it \citep{reisach2023r2sortability}, so both varsortability
and \(R^2\)-sortability belong in the reported description of any generated
grid. Both are computed here. Over the $44$ scale fits of the grid that carry
true edges, two per scenario, varsortability is $0.563 \pm 0.233$ and
\(R^2\)-sortability is $0.452 \pm 0.221$, both near the chance level of $0.5$.
The generator does not induce a variance ordering that agrees with the causal
order, so no entry in the comparison tables can be scoring for that reason.

The tiled bnlearn networks used as external ground truth are the more sortable
of the two constructions. Their varsortability runs from $0.462$ on sachs to
$0.773$ on child, and their \(R^2\)-sortability from $0.621$ on child to
$1.000$ on earthquake, so a learner has ordering information there that the
generated grid does not supply. That is a property of the tiling rather than of
any method scored against it, and it is stated because the construction is
widely used and, as far as we are aware, this statistic has not been reported
for it.

\paragraph{State which structural Hamming distance is meant.}
The original definition is over partially directed acyclic graphs
\citep{tsamardinos2006mmhc}, where a reversed edge and an unoriented edge are
distinct kinds of error. Most released code computes the distance on the
directed adjacency instead, and some implementations additionally offer a
switch for whether a reversed edge costs one or two. The two definitions are
different quantities and a distance reported without saying which is not
commensurable across papers. This work computes the distance on the directed
adjacency of the lagged edge set, as in \eqref{eq:bench_structure}, counting
each missing and each spurious edge once. Within the hypothesis class searched
here every candidate edge points forward in time and its reversal is not
admissible, so the orientation ambiguity that separates the two definitions in
the static case does not arise inside the search space.

\paragraph{Report AUPRC, not AUROC alone.}
Under the class imbalance of sparse edge scoring a large number of false
positives moves the area under the receiver operating characteristic very
little while moving precision a great deal, so the two curves rank methods
differently and only the precision-recall view reflects the imbalance
\citep{davis2006prroc}. The area under the precision-recall curve is therefore
the ranking metric reported wherever a method exposes scored candidates, and
the thresholded triple of \eqref{eq:bench_structure} is reported alongside it
rather than in place of it.

\paragraph{Report regime recovery separately from edge recovery.}
A persistent latent state raises the level of every label it drives for as long
as it lasts, so each label is autocorrelated even with an empty parent set, and
a criterion that conditions on nothing attributes that autocorrelation to the
only mechanism in its hypothesis space, a lagged self edge. This is the
published mechanism by which autocorrelation inflates false positives in
time-series structure learning \citep{runge2018pitfalls}, and it is exactly the
confound the empty finest scale of \cref{sec:benchmark_grid} is built to expose.
Because a method can track the regime path well and still over-select edges, or
the reverse, the regime metrics of \eqref{eq:bench_accs} are reported per scale
and are never folded into the structural columns. Pooling the two into a single
score would hide precisely the trade-off the benchmark exists to measure.

\section{Results}
\label{sec:results}

Every number below comes from a stored run under \texttt{results/} and was read
back from the file named beside it. The synthetic grid holds $22$ scenarios
drawn from the three generating families, each fitted on three replicates, with
the replicate seed fixed by the scenario so that the same scenario and
replicate index reproduce the same dataset. Across the five grid experiments
this is $1176$ scored fits; the parameter ladder adds $81$ and the external
benchmark adds $120$.

One property of the generator governs how the structural tables read. Over the
$66$ (scenario, replicate) pairs of the full grid the generator places $731$
lagged edges, $482$ of them at scale $k=2$ and $249$ at scale $k=3$, so the
mean number of true edges per fit is $11.076$ and two of the three scales carry
a recovery test. Scale $1$ carries no edge in any replicate, so $\gP_{1,\ell}$
is empty for every label $\ell$. That emptiness is a measured property of the
generator rather than a convenience. The finest scale is a binary event-start
process, so a lagged edge there can only shift which minute of a block begins
an episode, which is worth less than the price of the edge, and the screen of
\cref{sec:benchmark_identifiability} rejects it in every family. A scale with
no true edge and no proposed edge is scored as perfectly recovered, with
precision, recall and $F_1$ all equal to $1$; a scale with no true edge and at
least one proposed edge scores $0$ on all three. Reported precision at scale
$1$ is therefore the fraction of fits that proposed nothing there, not a ratio
of counts, and any mean of $F_1$ across the three scales still carries that
free term. One of the $66$ pairs happens to carry no true edge at scale $3$
either, which is why the empty-graph row scores $0.015$ there rather than
exactly $0$.

\subsection{Structure recovery on the synthetic grid}
\label{sec:results:structure}

Table~\ref{tab:results:exp1} reports the ideal-recovery experiment
(\texttt{results/exp1/exp1\_results.csv}, $198$ fits), where the generator, the
number of regimes and the number of factors are all given to the learner and
only $\gP_{k,\ell}$ has to be found. Every fit on the grid
starts at the generating parameters, and the structure search starts at the
generating parent sets, so the structure reported for the full model is where
the penalized score moves from the truth rather than where a cold search
arrives; the empty-graph row clears those parent sets and runs no search.

\begin{table}[t]
\centering
\caption{Structure recovery on the synthetic grid, means over $22$ scenarios
and three replicates. Scale $1$ carries no true edge, so its precision, recall
and $F_1$ record the fraction of fits that proposed nothing there. In the
across-scales row the $F_1$ column is the mean over the three scales and the
SHD column is their sum. Source
\texttt{results/exp1/exp1\_results.csv}.}
\label{tab:results:exp1}
\begin{tabular}{llrrrr}
\toprule
Method & Scale & Precision & Recall & $F_1$ & SHD \\
\midrule
n$^2$DBN      & $k=1$ & 0.273 & 0.273 & 0.273 & 2.318 \\
              & $k=2$ & \textbf{0.864} & \textbf{0.598} & \textbf{0.680} & \textbf{3.379} \\
              & $k=3$ & \textbf{0.985} & \textbf{0.870} & \textbf{0.908} & \textbf{0.515} \\
              & across &      &       & 0.620 & 6.212 \\
\midrule
Empty graph   & $k=1$ & 1.000 & 1.000 & 1.000 & 0.000 \\
              & $k=2$ & 0.000 & 0.000 & 0.000 & 7.303 \\
              & $k=3$ & 0.015 & 0.015 & 0.015 & 3.773 \\
              & across &      &       & 0.338 & 11.076 \\
\midrule
Oracle parameters & all & 1.000 & 1.000 & 1.000 & 0.000 \\
\bottomrule
\end{tabular}
\end{table}

The total structural Hamming distance favors the full model. It averages
$6.212$ against $11.076$ for the empty graph, so proposing a structure saves
$4.864$ edge errors per fit on average against proposing nothing. Paired within
(scenario, replicate), the full model is better on $61$ pairs, tied on $5$ and
worse on none of $66$.

The decomposition locates the gain, and it falls on two scales rather than
one. At scale $2$ the paired difference is $-3.924$ edge errors per fit, better
on $58$ pairs, tied on $6$ and worse on $2$. At scale $3$ it is $-3.258$,
better on $64$ pairs, tied on $2$ and worse on none. At scale $1$, where the
truth is empty, the full model proposes on average $2.318$ edges that cannot be
correct, losing every pair in which it proposes anything at all ($0$ wins, $18$
ties, $48$ losses); it proposes nothing there on $18$ of the $66$ fits and
$3.188$ edges on each of the rest. Summed, a gain of about seven edge errors
across the two scales that carry the truth is offset by about two false
positives at the scale that carries none, so the total keeps most of the gain,
and reporting it alone would still hide all three effects.

The silence at the finest scale is selective rather than a general reluctance
to propose. The same fits propose at least one edge at scale $2$ on $63$ of
$66$ and at scale $3$ on $65$ of $66$.

Precision at scale $2$ runs $0.851$ on family I, $0.775$ on family II and
$0.959$ on family III, against recall $0.655$, $0.502$ and $0.668$. At scale
$3$ precision runs $0.917$, $1.000$ and $1.000$ against recall $0.833$, $0.836$
and $0.921$. The selection is conservative in the same way in all three
families, and the coarsest scale is the easier of the two scales that carry
true edges, returning more of the truth at a higher precision than the middle
scale in every family.

\subsection{Prediction against the empty graph}
\label{sec:results:prediction}

Held-out fit is scored by the validation loss $-\sum_k \log Z_k$ evaluated on
the validation split with the learned $\Theta$ and the learned graph, where
$\log Z_k$ is the normalizer the scale-$k$ E-step returns, the forward-backward
partition function of the regime chain evaluated with the factor posterior
folded into the emission terms. Lower is better. Because the comparison is
between two methods fitted on the same split, the paired difference is the
quantity of interest and the level is not.

Over the $66$ pairs of \texttt{results/exp1/exp1\_results.csv} the full model
reaches a mean of $-36296.45$ against $-34951.83$ for the empty graph. The mean
paired difference is $-1344.63$ nats and the median is $-632.20$; the full model
is better on $64$ of $66$ pairs. Expressed relative to the
empty-graph value of each pair, the mean reduction is $4.0$ percent. The
selected edges therefore pay for themselves in held-out evidence on the same
fits where they also reduce the total edge distance.

The oracle-parameter row records a validation loss of $-3485.14$, far above
both. That row installs the generating parameters and fits nothing, so it is
the ceiling for structure and for latent state, not for prediction; a
parameter set tuned on the training split will fit a held-out split of the same
series better than the parameters that generated it.

\subsection{Latent-state recovery against the oracle ceiling}
\label{sec:results:latent}

Table~\ref{tab:results:latent} scores the recovered regime path $S_t^{(k)}$
against the generating path, with the oracle-parameter row as the ceiling.

\begin{table}[t]
\centering
\caption{Latent regime recovery, means over the full grid. The oracle row
installs the generating $\Theta$ and runs inference only, so it bounds what any
estimator on this grid can reach. Source
\texttt{results/exp1/exp1\_results.csv}.}
\label{tab:results:latent}
\begin{tabular}{lrrrrrr}
\toprule
& \multicolumn{3}{c}{Adjusted Rand index} & \multicolumn{3}{c}{Best-permutation accuracy} \\
\cmidrule(lr){2-4}\cmidrule(lr){5-7}
Method & $k=1$ & $k=2$ & $k=3$ & $k=1$ & $k=2$ & $k=3$ \\
\midrule
n$^2$DBN          & --- & 0.295 & 0.907 & --- & 0.767 & 0.961 \\
Empty graph       & --- & 0.237 & 0.908 & --- & 0.736 & 0.967 \\
Oracle parameters & --- & 0.219 & 0.967 & --- & 0.742 & 0.987 \\
\bottomrule
\end{tabular}
\end{table}

Most of this table is empty on purpose, and that is the result rather than a
gap in it. A scale whose generating path never switches has no regime to
recover, and the screen of \cref{sec:benchmark_identifiability} places a regime
only where its own oracle clears the constant predictor, so on the minute grid
scales $1$ and $2$ carry one state and are reported as not applicable rather
than as a perfect score. The scale-$2$ column holds the three day-grid fits,
whose hour scale does carry a regime, and the scale-$3$ column holds $25$ of the
$27$ family II fits, the two omissions being replicates whose coarsest test
partition happened not to switch.

Where a regime is placed, it is recovered. At the coarsest scale the full model
reaches an adjusted Rand index of $0.907$ and an accuracy of $0.961$ against an
oracle at $0.967$ and $0.987$, and it recovers the path exactly on $17$ of its
$25$ scored fits, the same count as the oracle. The gap to the ceiling is
$0.060$ of index and $0.026$ of accuracy. At the hour scale of the day grid the
three rows are $0.295$, $0.237$ and $0.219$ of index, where the fitted model is
ahead of the oracle on three fits, which is inside the noise of three fits and
should not be read as beating a ceiling.

An earlier version of this benchmark placed a regime at all three scales and
reported near-identical rows for the full model, the empty graph and the oracle,
with scale-$2$ accuracies of $0.529$, $0.521$ and $0.516$ against a floor of one
half. That was not three methods tying; it was three methods scoring a quantity
the data did not identify. The screen removed those scales from the grid, and
what is left is the one comparison that carries information.

The latent factor $H_t^{(k)}$ tells the same story with a smaller gap. On the
family III scenarios the mean absolute correlation between the recovered and
generating factor trajectories is $0.457$, $0.470$ and $0.694$ at the three
scales for the full model, against $0.540$, $0.555$ and $0.786$ for the oracle.
The estimator gives up between $0.083$ and $0.092$ of correlation at each
scale, and the ceiling itself is well below one.

\subsection{The contribution of the cross-scale coupling}
\label{sec:results:coupling}

The coupling experiment (\texttt{results/exp2/exp2\_results.csv}, $132$ fits)
removes the downward pass, so that the regime chain of scale $k$ is driven only
by that scale's own data and context $\gK_t^{(k)}$ and never by the posterior
of its parent epoch $\parmap_k(t)$. Everything else, including the allocation
that ties $X_t^{(k)}$ and $Y_t^{(k)}$ to their parent totals, is unchanged.

On the minute grid the coupling now buys exactly nothing, and the reason is
structural rather than empirical. Its sixty fits return identical graphs and
identical validation losses to the last recorded digit, on every column. With
$M_k = [1,1,M]$ the parent posterior passed downward at scales $1$ and $2$ is a
distribution over one state, which is the same object whether or not it is
passed, so the coupled and uncoupled arms are the same computation. The
experiment does not test the coupling there.

The six day-grid fits are the ones that test it, since the hour scale of that
hierarchy does carry a regime with a two-state parent above it. There the
uncoupled arm is ahead on held-out evidence by a mean of $362.5$ nats, with the
coupled model better on $2$ of the $6$ pairs; the structural columns are
identical. Six fits with that spread support no claim in either direction, so
the honest statement is that this benchmark no longer measures the contribution
of the cross-scale coupling.

This is a direct cost of the screen, and it is worth stating as one. An earlier
version of this grid put a regime at all three scales and recorded the coupling
as worth $61.89$ nats and $0.061$ edge errors, all of it on family II. What the
coupling was passing downward there was a finer-scale regime whose own oracle
could not beat a constant predictor, so that figure measured the benchmark
rather than the mechanism. Removing the unrecoverable latent removed the
measured benefit with it. Testing the coupling properly needs a hierarchy in
which two adjacent scales both carry a latent the screen admits, which on the
present grid is one hierarchy out of two and three scenarios out of twenty-two.

\subsection{Component ablation}
\label{sec:results:ablation}

The ablation (\texttt{results/exp3/exp3\_results.csv}, $270$ fits over the $18$
family II and family III scenarios) removes one component at a time. Dropping
the factor sets $r_k = 0$; dropping the regime sets $M_k = 1$.

\begin{table}[t]
\centering
\caption{Component ablation on the family II and family III scenarios, means
over $18$ scenarios and three replicates. Source
\texttt{results/exp3/exp3\_results.csv}.}
\label{tab:results:ablation}
\small
\begin{tabular}{lrrrrrrr}
\toprule
Method & $F_1$ $k{=}2$ & $F_1$ $k{=}3$ & SHD $k{=}1$ & SHD $k{=}3$ & SHD all & Acc.\ $k{=}3$ & Val.\ loss \\
\midrule
n$^2$DBN          & \textbf{0.677} & \textbf{0.918} & 2.362 & \textbf{0.500} & \textbf{6.259} & \textbf{0.961} & \textbf{-37814.78} \\
No latent factor  & 0.677 & 0.918 & 2.352 & 0.537 & 6.259 & 0.961 & -37767.39 \\
No regime chain   & 0.676 & 0.784 & 2.389 & 1.796 & 7.574 & 0.376 & -37692.81 \\
Empty graph       & 0.000 & 0.000 & 0.000 & 3.815 & 11.130 & --- & -36454.00 \\
Oracle parameters & 1.000 & 1.000 & 0.000 & 0.000 & 0.000 & 0.987 & -4033.88 \\
\bottomrule
\end{tabular}
\end{table}

The two components contribute unequally, and one of them contributes nothing
structural. Removing the factor leaves the total distance equal to three
decimals, equal on $45$ of the $54$ pairs and better on $4$ against worse on
$5$, and leaves scale-$2$ and scale-$3$ $F_1$ unchanged to three decimals. Its
whole contribution is $47.39$ nats of held-out evidence, identical on $27$
pairs, better on $24$ and worse on $3$. Removing the regime chain costs
$121.97$ nats, better on all $27$ pairs that differ and worse on none, and it
also costs structure at the coarsest scale. The scale-$3$ distance rises from
$0.500$ to $1.796$, scale-$3$ $F_1$ falls from $0.918$ to $0.784$, the total
distance rises from $6.259$ to $7.574$ and the mean $F_1$ across scales falls
from $0.606$ to $0.561$.

Scale-$2$ $F_1$ alone does not show this. It is $0.676$ without the regime
chain against $0.677$ with it, so the regime pays for itself at the coarsest
scale and not at the middle one. That is where persistent drift was always
claimed to live, and it is now the only scale at which the screen of
\cref{sec:benchmark_identifiability} admits a regime on this grid, which is why
the regime's held-out contribution is $122$ nats here where an earlier version
of this benchmark recorded $1250$. That earlier figure credited the model for
tracking a finer-scale regime whose own oracle could not beat a constant, so it
measured the benchmark rather than the method. The regime accuracy column makes
the same point from the other side: with the chain removed the fitted path is
constant, and the $0.376$ it scores is the frequency of the modal true regime,
which is the floor rather than a recovery. The regime chain is what makes the
coarsest scale recoverable, and the factor is a small likelihood term that the
structure search does not see. That is an honest negative for the factor as a
structural device.

\subsection{Behavior under misspecification}
\label{sec:results:misspec}

The misspecification experiment
(\texttt{results/exp4/exp4\_results.csv}, $36$ fits) generates from a family
that violates a modeling assumption and fits the same model. The
\emph{heavier tail} arm draws counts from a distribution with a heavier tail
than the fitted emission; the \emph{omitted context} arm generates with a
context that is then withheld from $\gK_t^{(k)}$ at fit time. No fit failed and
no run raised an error.

\begin{table}[t]
\centering
\caption{Behavior under misspecification, means over four scenario shapes and
three replicates per arm. Structure metrics only, since the three arms generate
from different data and their held-out losses are not on a common scale. Source
\texttt{results/exp4/exp4\_results.csv}.}
\label{tab:results:misspec}
\small
\begin{tabular}{lrrrrrrr}
\toprule
Arm & Prec.\ $k{=}2$ & Rec.\ $k{=}2$ & $F_1$ $k{=}2$ & $F_1$ $k{=}3$ & SHD $k{=}1$ & SHD $k{=}2$ & Sweeps \\
\midrule
Well specified   & 0.965 & 0.578 & 0.707 & 1.000 & 3.000 & 3.333 & 6.33 \\
Heavier tail     & 0.875 & 0.611 & 0.710 & 0.976 & 3.583 & 3.917 & 6.75 \\
Omitted context  & 0.846 & 0.536 & 0.650 & 0.889 & 3.667 & 4.500 & 6.17 \\
\bottomrule
\end{tabular}
\end{table}

The two arms fail in different directions, and they reach different scales. A
heavier tail than the emission assumes trades precision for recall at the middle
scale, precision falling from $0.965$ to $0.875$ while recall rises from $0.578$
to $0.611$, so its $F_1$ there is level at $0.710$ against $0.707$; what it costs
is the coarsest scale, whose $F_1$ falls from an exact $1.000$ to $0.976$, and
the empty finest scale, where false positives rise from $3.000$ to $3.583$.
Withholding the context costs both halves of the middle scale, precision
$0.846$ and recall $0.536$, and it is the arm that reaches the coarsest scale
hardest, where precision falls to $0.917$ and recall to $0.875$. Neither arm
breaks the model. Both keep the aggregation exact, at a reconciliation error of
exactly $0$ for counts and durations on all $36$ fits, and neither takes
materially longer to converge than the well-specified arm, at $6.75$ and $6.17$
outer sweeps against $6.33$.

The mean of $F_1$ across scales is the wrong summary for this table and is left
out of it. That mean is $0.590$ under the heavier tail against $0.569$ under the
well-specified arm, reversing the ordering of the coarsest scale, and the
reversal comes from the empty finest scale, where the heavier-tail arm proposes
nothing on one of its twelve fits and the other two arms on none of theirs. The
heavier tail is not the easier setting.

\subsection{Against eight published structure learners}
\label{sec:results:comparison}

The comparison experiment (\texttt{results/exp5/exp5\_results.csv}, $540$ fits)
runs eight published structure learners on the same counts $X_t^{(k)}$ that
the local score sees, then fits this model's parameters on whatever graph each
returns, so that every row is scored by the same metric code on the same data.

\begin{table}[t]
\centering
\caption{Eight published structure learners against the latent-conditional
local score and against reporting nothing, means over $18$ scenarios and three
replicates. Sorted by total structural Hamming distance. Source
\texttt{results/exp5/exp5\_results.csv}.}
\label{tab:results:comparison}
\small
\begin{tabular}{lrrrrrrr}
\toprule
Method & Prec.\ $k{=}2$ & Rec.\ $k{=}2$ & $F_1$ $k{=}2$ & $F_1$ $k{=}3$ & SHD all & $\times$ empty & Val.\ loss \\
\midrule
n$^2$DBN local score & \textbf{0.867} & \textbf{0.585} & \textbf{0.677} & \textbf{0.918} & \textbf{6.259} & \textbf{0.56} & \textbf{-37814.78} \\
Empty graph          & 0.000 & 0.000 & 0.000 & 0.000 & 11.130 & 1.00 & -36454.00 \\
Chow-Liu tree        & 0.210 & 0.196 & 0.195 & 0.082 & 23.296 & 2.09 & -36933.15 \\
PC                   & 0.073 & 0.129 & 0.088 & 0.031 & 26.815 & 2.41 & -36784.06 \\
GES                  & 0.146 & 0.220 & 0.169 & 0.043 & 27.685 & 2.49 & -36960.92 \\
Granger, BIC         & 0.164 & 0.351 & 0.217 & 0.341 & 29.389 & 2.64 & -37517.01 \\
Hill climbing, BIC   & 0.139 & 0.288 & 0.183 & 0.043 & 30.574 & 2.75 & -37452.14 \\
Exhaustive, BIC      & 0.139 & 0.288 & 0.183 & 0.043 & 30.574 & 2.75 & -37452.19 \\
DYNOTEARS            & 0.155 & 0.327 & 0.203 & 0.363 & 30.759 & 2.76 & -37556.48 \\
Hill climbing, BDeu  & 0.136 & 0.288 & 0.181 & 0.096 & 33.241 & 2.99 & -37454.06 \\
\bottomrule
\end{tabular}
\end{table}

All eight land between $2.1$ and $3.0$ times worse than reporting nothing.
Their $F_1$ at the middle scale runs from $0.088$ to $0.217$ against $0.677$
for the local score, and no fit of any of the eight beats the empty graph on
total distance in any of the $54$ pairs. The local score is the only entry
below the empty graph, by a mean of $4.870$ edge errors with $50$ wins, $4$ ties
and no losses.

The coarsest scale is where the separation is widest, and it is a column this
comparison could not carry until the grid placed a true edge there. Six of the
eight score below $0.100$ at scale $3$, and the two that do not, DYNOTEARS at
$0.363$ and the Granger test at $0.341$, still reach less than four tenths of
the local score's $0.918$. Those two are the only entries with an explicitly
temporal search. The other six run from $0.031$ to $0.096$ at that scale, near
enough to the empty graph's $0.000$ that what they mostly contribute at the
hourly scale is false positives.

\paragraph{The mechanism is a lagged self edge.}
The generator drives every label of a scenario from a shared regime $S_t^{(k)}$
that persists across epochs, so every label it drives is autocorrelated whether
or not any lagged edge points into it. A selection criterion that conditions on
nothing sees that autocorrelation and has one place to put it, namely an edge
from the label to itself at lag one. Splitting recall at scale $2$ by whether
the true edge points from a label to itself or from one label to another states
this directly rather than by inference. The Granger test recovers $1.000$ of
the true self edges and $0.210$ of the true cross edges; DYNOTEARS recovers
$0.986$ against $0.190$, and both hill climbing variants and exhaustive search
$0.957$ against $0.167$; PC recovers $0.729$ against $0.002$ and GES $0.543$
against $0.165$. The Chow-Liu tree is the one exception and is weak on both,
$0.429$ and $0.167$. The local score recovers $0.929$ of the self edges and
$0.524$ of the cross edges, and is the only entry of the nine above a quarter on
the cross edges. At the coarsest scale it is $0.944$ and $0.850$, so once the
regime is conditioned on the two kinds of edge are recovered at nearly the same
rate.

The proposed edge sets say the same thing from the other side. Across the $54$
(scenario, replicate) pairs of this experiment the generator places $391$ true
edges at scale $2$, of which $61$ point from a label to itself, a true
self-edge share of $0.156$. The proposed share is $0.985$ for PC, $0.886$ for
DYNOTEARS, $0.883$ for hill climbing with BIC and for exhaustive search,
$0.879$ for hill climbing with BDeu, $0.838$ for the Granger test, $0.743$ for
GES and $0.634$ for the Chow-Liu tree, against $0.280$ for the local score. Six
of the eight place more than four fifths of their scale-$2$ budget on a
construct that is about a sixth of the truth. Conditioning the score
on $q(\gS,\gH)$ removes that autocorrelation before an edge is priced, which is
what the precision column of Table~\ref{tab:results:comparison} records.

\paragraph{Held-out likelihood does not order the structures.}
Every one of the nine non-empty rows of
Table~\ref{tab:results:comparison} improves on the empty graph in validation
loss, by between $330.1$ and $1360.8$ nats, and the ordering it induces is close
to unrelated to the structural one. Spearman rank correlation between
mean total distance and mean validation loss across the ten rows is $-0.333$.
The empty graph has the worst validation loss of the ten and
the second-best distance; DYNOTEARS has the second-best validation loss and the
ninth-best distance. The local score leads both columns here, which is a
property of that row rather than of the criterion, since the other nine rows
are ordered almost independently by the two. Held-out likelihood is a real
quantity and it is not a proxy for structure recovery on this grid, so a paper
that reports only one of the two can support a conclusion the other refutes.

\subsection{Parameter recovery on a sample-size ladder}
\label{sec:results:params}

Structure recovery says which edges came back and nothing about the conditional
distributions attached to them. The parameter experiment
(\texttt{results/exp7/exp7\_results.csv}, $81$ rows) scores the learned
$\Theta$ against the generating $\Theta$ directly, in identified form, so that
two parameter sets inducing the same distribution score zero. It holds the
hierarchy shape fixed and varies only $T_1$, the number of finest-scale epochs,
over $10{,}080$, $20{,}160$ and $40{,}320$, chosen so that the coarsest scale
receives $168$, $336$ and $672$ epochs and the structural benchmark's own
horizon sits in the middle. Two starting points are separated. The
\emph{oracle start} begins the estimator at the generating parameters and
measures the estimator alone; the \emph{data start} builds its initial
parameters from moments of the data and measures the estimator together with
its initializer, which is the only setting that transfers to a log with no
ground truth.

The criterion was set before the study ran. Parameter recovery is a consistency
claim, so the error has to fall as $T_1$ grows; an error that does not fall
with sample size is a bias rather than an estimate.

\begin{table}[t]
\centering
\caption{Mean identified parameter error against $T_1$, over three generating
families and three replicates. The ladder is set from the coarsest scale rather
than the finest, so these three rungs give it $168$, $336$ and $672$ epochs;
the middle rung is the horizon the structural benchmark itself runs at. The third stored method installs the generating
parameters and fits nothing, so its error is zero at every rung and it is not
shown. The error averages whichever scale-free blocks a family identifies,
which is $5$ for family I, $6$ for family II and $8$ for family III, so a level
is comparable along a row and not down the family rows; only the per-family rows
carry a criterion. Source \texttt{results/exp7/exp7\_results.csv}.}
\label{tab:results:ladder}
\begin{tabular}{lrrr}
\toprule
Start & $T_1 = 10080$ & $T_1 = 20160$ & $T_1 = 40320$ \\
\midrule
Oracle start (estimator alone) & 0.0725 & 0.0695 & \textbf{0.0514} \\
Data start (estimator and initializer) & 0.2822 & 0.2595 & \textbf{0.2504} \\
\midrule
Family I, data start   & 0.3366 & 0.2735 & 0.3001 \\
Family II, data start  & 0.2145 & 0.2119 & \textbf{0.1924} \\
Family III, data start & 0.2956 & 0.2932 & \textbf{0.2588} \\
\bottomrule
\end{tabular}
\end{table}

Both starts meet the criterion, and the reading of this table changed when the
benchmark did. From the oracle start the mean identified parameter error goes
$0.0725$, $0.0695$ and $0.0514$, falling at every rung. From the data start it
goes $0.2822$, $0.2595$ and $0.2504$, also falling at every rung. Per family the
data start falls at every rung for family II, $0.2145$, $0.2119$ and $0.1924$,
and for family III, $0.2956$, $0.2932$ and $0.2588$. Family I is the one that
does not, $0.3366$, $0.2735$ and $0.3001$, and it carries neither a regime nor a
factor, so its reported error is the coarsest-scale emission block alone.

An earlier version of this grid placed a regime at all three scales and reported
a data start rising at every rung, $0.4261$, $0.4533$ and $0.4678$, with family
II alone at $0.6461$, $0.7933$ and $0.8445$. That rise is what the screen of
\cref{sec:benchmark_identifiability} removed. Three of family II's scored blocks
were the finest-scale transition table and the two allocation shares, all driven
by a regime whose own oracle could not beat a constant predictor, so the quantity
that grew with the sample size was the noise in an unidentified block rather than
a fit converging elsewhere. With the regime placed only where it is recoverable,
family II falls to about a quarter of its former level and becomes the best of
the three families rather than twice the worst. Families I and III are unchanged
to the digit, because a scenario carrying one regime at every scale is untouched
by where a regime is placed, which is the control this comparison has that a
rerun on a freshly randomized grid would not.

What remains is the gap between the two starts rather than a failure of either.
At the longest rung the data start sits at $0.2504$ against $0.0514$ from the
oracle start, so the initializer, not the estimator, is what the remaining error
measures. The coarsest transition table is recovered closely from either start,
at $0.0242$, $0.0417$ and $0.0167$ in total variation from the data start, and
the coarsest lagged support is recovered at $4.00$, $4.22$ and $5.11$ selected
keys against a true $4.00$, so the over-selection that remains is at the longest
rung and is small.

\paragraph{The objective and the parameter error no longer disagree.}
Comparing the two starts on their final objective rather than on their parameter
error separates a fit that found a worse optimum from one that found an equally
good optimum at different parameters, and the two have different remedies. Over
the $27$ paired fits the data start reaches a strictly worse objective in $3$ of
the $9$ family I fits, $2$ of $9$ for family II and none for family III, and in
every one of those it is worse on the parameter error as well. There is no fit
in which a worse objective comes with better parameters, where an earlier version
of this grid had $4$ of $27$. That pathology was the sign that the error was
being scored partly on quantities the likelihood does not see, and it is gone
with them. Two conditions attach: equality is equality at a relative tolerance,
not proven equality, and the objective compared is the monitored one, which
Section~\ref{sec:lim_meanfield} records as contained rather than exact.

\paragraph{The coarsest scale is recoverable from a data start, and only with
the level correction.}
Counting the lagged coefficient keys each fit ends with separates a coefficient
that is wrong from one that was never estimated. Against a true $4.00$ keys per
fit at the coarsest scale, a data start recovers $4.00$, $4.22$ and $5.11$
across the three rungs, so the support is recovered and the over-selection that
remains is at the longest rung. Removing two things restores the earlier behavior
exactly: initializing the coarsest level from the observed mean, which already
contains the contribution of the lagged parents, and letting the parameter step
refit that level against the empty graph before the structure search runs. With
both removed the same fits recover $0.00$, $0.00$ and $0.67$. The level is
therefore not a nuisance parameter here. Fitted before the search and against a
graph with no edges in it, it absorbs exactly what a true lagged edge would
explain, and the search then meets a model that already accounts for the data.

\paragraph{The short coarse series was the likeliest explanation, and it is
ruled out.}
The ladder holds the hierarchy shape fixed and varies $T_1$ alone, so the
coarsest scale receives $168$, $336$ and $672$ epochs across the three rungs, of
which a fit sees $117$, $235$ and $470$ after the training split, while the
finest receives $10{,}080$, $20{,}160$ and $40{,}320$. That scale now carries lagged
coefficients of its own, which made it the natural suspect. Reading the error
per scale and per parameter block removes it from suspicion on three counts.
The lagged coefficients of any scale are excluded from the reported error by
construction, since it averages the scale-free blocks alone, and perturbing the
coarsest lagged coefficients of a set of true parameters leaves it unchanged to
machine precision. Every coarsest-scale block falls at every rung from the
oracle start, the relative emission error $0.603$, $0.562$, $0.350$ and the
transition distance $0.224$, $0.074$, $0.068$. And recomputing the reported
error over subsets of the scales separates the readings directly: the coarsest
scale alone falls, $0.490$, $0.448$, $0.271$, while dropping it leaves the
ladder broken, $0.416$, $0.529$, $0.248$. A ladder that grew every scale at the
same rate would therefore not repair this one.

What the breakdown puts in its place is narrower and less comfortable. From the
oracle start the non-monotonicity is confined to family II, whose level is set
by the relative emission error at the \emph{finest} scale, reaching $7.9$ and
$9.5$ on its worst fits against $0.2$ at the coarsest, while families I and III
fall at every rung. That is the binary event-start scale on which the leak of
\cref{sec:lim_leak} sits and which the screen of
\cref{sec:benchmark_identifiability} places below its identifiability floor. The
residual therefore concentrates where this benchmark cannot separate a fitting
failure from a scale it was never able to identify, so the criterion stands as
failed and the cause is attributed no further than that.

\paragraph{The leak also shows up in the coefficients.}
The same rows record the absolute learned lag-coefficient mass placed on lagged
keys outside the true support, which is the parameter-space counterpart of a
false-positive edge. That mass is largest at the longest rung from either start
instead of shrinking. From the oracle start it grows at every step, from
$0.281$ at $T_1 = 1440$ to $1.194$ at $T_1 = 2880$ and $1.585$ at
$T_1 = 5760$; from the data start it runs $0.278$, $0.860$ and $1.376$ over the
same three rungs. The share of fits carrying any off-support mass runs $0.111$,
$0.444$, $0.556$ from the oracle start and $0.111$, $0.333$, $0.444$ from the
data start. It is not confined to one family. At the longest rung from the
oracle start the mass per fit is $2.384$ on family I, $1.458$ on family II and
$0.913$ on family III. It is confined to one scale. Of the $1.585$ overall at
that rung, $1.358$ sits at scale $1$, which is the scale the generator leaves
empty. The co-movement a persistent regime induces between
labels is being absorbed into
lagged coefficients, and a longer series makes that absorption more confident,
not less. The structural score suppresses it as an edge; the parameter step
still pays for it as a coefficient.

\subsection{External ground truth}
\label{sec:results:external}

Every number to this point is scored against data this benchmark generated,
which is the weakest evidence base for a claim about a learner. The external
experiment (\texttt{results/exp8/exp8\_structure.csv} and
\texttt{results/exp8/exp8\_cpt.csv}) repeats the structure comparison against
released networks. The bnlearn repository \citep{scutari2010bnlearn}
distributes fully parameterized networks with an exact directed graph and exact
conditional probability tables, and pgmpy \citep{ankan2024pgmpy} vendors them
inside its wheel, so the experiment runs with no network access. A static
network becomes a two-slice temporal network by mapping each edge onto a lag-1
edge, the standard construction with published precedent
\citep{trabelsi2013benchmarking, ouyang2023divideconquer}. Five networks (asia,
cancer, child, earthquake and sachs, from $5$ to $20$ nodes and $4$ to $25$
edges), three replicates and $4000$ epochs give $120$ structure fits including
the empty-graph reference, with no errors. The largest network of the earlier
version of this experiment, insurance at $27$ nodes, is not included, because
selecting the free quantities per dataset multiplies the number of fits by
eleven and that run did not complete. The omission is of size rather than of an
inconvenient result.

A tiled network is flat, so a two-scale hierarchy is imposed by summing blocks
of consecutive epochs and only the finest scale is scored, where the candidate
set is exactly the tiled ground truth. The block length is declared as the
square root of the series length, for the reason given in
\cref{sec:results:selection}, and the regime count, the factor dimension and
the edge penalty are selected per dataset on held-out log evidence.

Structure and parameters are scored on the same graph. A learned parent set
that differs from the true one produces a table of a different shape, so tables
cannot be compared entry by entry across methods. Both the learned model and
the released network do answer one common question, namely the distribution of
a node at epoch $t$ given the whole previous epoch, so that conditional is what
is scored: the table is estimated from the sampled data on the learned parents
alone, the released tables are read at the true parent configuration, and the
two are compared in total variation, averaged over epochs. An empty parent set
is a valid answer under this score and is charged the distance from the
marginal to the truth.

\begin{table}[t]
\centering
\caption{Structure and conditional probability recovery against released
bnlearn networks tiled into two-slice temporal networks, means over five
networks and three replicates. Same structural metric code as
Table~\ref{tab:results:comparison}, external ground truth. The last column is
the mean total variation distance of the conditional each method's own learned
parent set implies, so structure and parameters are scored on one graph.
Source \texttt{results/exp8/exp8\_structure.csv}.}
\label{tab:results:external}
\begin{tabular}{lrrrrr}
\toprule
Method & SHD & Precision & Recall & $F_1$ & CPT TV \\
\midrule
Hill climbing, BIC     & \textbf{1.13} & 0.987 & 0.920 & \textbf{0.950} & 0.0110 \\
Exhaustive, BIC        & \textbf{1.13} & 0.987 & 0.920 & \textbf{0.950} & 0.0110 \\
Granger, BIC           & 2.40 & 0.942 & 0.895 & 0.913 & 0.0172 \\
DYNOTEARS              & 2.80 & \textbf{1.000} & 0.778 & 0.866 & 0.0284 \\
n$^2$DBN, selected     & 3.27 & 0.869 & 0.762 & 0.806 & 0.0239 \\
Hill climbing, BDeu    & 4.60 & 0.682 & \textbf{0.967} & 0.778 & 0.0112 \\
Chow-Liu tree          & 7.20 & 0.776 & 0.519 & 0.610 & 0.0582 \\
Empty graph            & 11.60 & 0.000 & 0.000 & 0.000 & 0.1054 \\
\bottomrule
\end{tabular}
\end{table}

The ordering is not the same as on the synthetic grid, and the reason is a
property of the two benchmarks rather than a contradiction. A tiled static
network has no persistent latent state, so the confound that
\cref{sec:results:comparison} isolates is absent and the classical criteria
have nothing to be confounded by. Hill climbing under BIC leads at $F_1 =
0.950$. The local score, with its free quantities selected rather than
declared, reaches $0.806$, up from $0.532$ under the declared setting of the
first version of this experiment.

The per-network breakdown locates what is left. On asia the local score matches
hill climbing exactly, $F_1 = 0.933$ for both, and on earthquake both recover
the network exactly, $F_1 = 1.000$ at zero structural distance. On child, the
largest network here at $20$ nodes, it reaches $0.870$ against $0.933$ at a
higher precision, $0.958$ against $0.933$. On sachs it reaches $0.812$ against
$0.980$. The one clear loss is cancer, $0.413$ against $0.905$, a five-node
network with four edges, where the finest-scale allocation has very little to
condition on. Only three true edges separate the two learners across both
networks; \cref{sec:lim_recall} names them and reports what is known about why,
having tested and withdrawn the block-structure explanation we first proposed.

Conditional probability recovery follows structure recovery and separates the
methods less. The local score reaches a mean total variation distance of
$0.0239$ per conditioning row against $0.0110$ for hill climbing and $0.1054$
for the empty graph, so the tables it implies are far closer to the released
ones than reporting no dependence, at about twice the distance of the best
structural entry. Table~\ref{tab:results:cpt} reports the floor the sample
size allows, estimating each table on the true graph.

\begin{table}[t]
\centering
\caption{Conditional probability table recovery against the released tables
given the true structure, which is the floor the sample size allows rather than
a method score. Mean total variation distance per conditioning row. A node fit
is one node on one replicate. Source \texttt{results/exp8/exp8\_cpt.csv}.}
\label{tab:results:cpt}
\begin{tabular}{lrrrr}
\toprule
Network & Node fits & Mean & Median & Max \\
\midrule
asia       & 24 & 0.0198 & 0.0078 & 0.1317 \\
cancer     & 15 & 0.0122 & 0.0057 & 0.0671 \\
child      & 60 & 0.0194 & 0.0176 & 0.0457 \\
earthquake & 15 & 0.0315 & 0.0278 & 0.1160 \\
sachs      & 33 & 0.0240 & 0.0245 & 0.0465 \\
\midrule
All        & 147 & 0.0210 & 0.0171 & 0.1317 \\
\bottomrule
\end{tabular}
\end{table}

Given the true structure, $94.6$ percent of the $147$ node fits fall below a
total variation distance of $0.05$ and $98.6$ percent below $0.10$, so the
sampled data at this length identifies the released tables well and the
differences in the last column of Table~\ref{tab:results:external} are
attributable to structure rather than to estimation noise.

\subsection{Selecting the free quantities of the model}
\label{sec:results:selection}

A latent-variable model carries free quantities that a latent-free learner does
not. The classical learners in Table~\ref{tab:results:external} each choose
their own penalty through their own score, while the number of regimes $M_k$,
the factor dimension $r_k$ and the edge penalty $\lambda_k$ were held at
declared values. A table produced that way compares declarations, not criteria.
All three are therefore selected on held-out log evidence, computed on the
final quarter of the training series and never on the evaluation data, with the
graph budget and the block lengths of the hierarchy left as declarations
because no held-out score supplies a question about which lags and which
timescales are of interest.

Table~\ref{tab:results:frontier} sweeps $\lambda_k$ on both benchmarks and
reports $F_1$ beside the held-out negative log evidence.

\begin{table}[t]
\centering
\caption{Edge-penalty frontier on both benchmarks, with the held-out negative
log evidence of the same fit. External columns are four tiled networks (asia,
cancer, earthquake, sachs) at $T = 4000$ and block length $63$; synthetic
columns are all $22$ scenarios of the grid at the true hierarchy. Every cell
is read back from \texttt{results/exp9/exp9\_measurements.csv}. The minimum of
the evidence column and the maximum of every $F_1$ column fall on the same row
on the external networks and on different rows on the grid. The grid carries
true edges at the middle and the coarsest scale, so both are reported, together
with the mean over all three scales.}
\label{tab:results:frontier}
\small
\begin{tabular}{lrrrrrr}
\toprule
& \multicolumn{2}{c}{External} & \multicolumn{4}{c}{Synthetic} \\
\cmidrule(lr){2-3}\cmidrule(lr){4-7}
$\lambda_k$ & $F_1$ & held-out NLL & $F_{1,k=2}$ & $F_{1,k=3}$ & $F_{1,\mathrm{mean}}$ & held-out NLL \\
\midrule
$2.0$  & 0.666 & $-68061.4$ & \textbf{0.685} & 0.833 & \textbf{0.597} & $-37890.7$ \\
$1.0$  & 0.725 & $-68065.5$ & 0.669 & 0.940 & 0.582 & $-37923.6$ \\
$0.5$  & 0.812 & $-68067.6$ & 0.602 & \textbf{0.955} & 0.519 & $\mathbf{-37958.6}$ \\
$0.25$ & \textbf{0.869} & $\mathbf{-68070.9}$ & 0.549 & 0.857 & 0.468 & $-37923.8$ \\
$0.1$  & 0.791 & $-68068.9$ & 0.589 & 0.731 & 0.440 & $-37898.8$ \\
\bottomrule
\end{tabular}
\end{table}

Held-out evidence selects the maximizing penalty on the external networks and
does not on the grid. Externally it picks $\lambda_k = 0.25$, which is also
where $F_1$ peaks, and the evidence turns back up on exactly the row where
$F_1$ turns back down. On the grid it picks $0.5$ while $F_{1,k=2}$ peaks at
$2.0$, and the miss is not a tie: $F_{1,k=2}$ is $0.602$ at the selected
penalty against $0.685$ at the maximizing one. It is worse than that on the
columns $F_{1,k=2}$ hides. The strictest setting, $\lambda_k = 2$, reaches a
scale-$2$ precision of $0.864$ and a total structural distance of $6.77$ against
$0.520$ and $12.82$ at the selected $0.5$. Held-out evidence therefore prefers a
graph that is clearly worse at the middle scale.

The frontier was previously scored at the middle scale alone, which raised the
possibility that the disagreement was an artifact of reading a two-scale truth
through a one-scale metric. Rescoring it across both truth-bearing scales
removes that possibility rather than narrowing it to one scale. Held-out
evidence picks $0.5$; $F_{1,k=3}$ peaks at $0.5$ as well, while $F_{1,k=2}$ and
$F_{1,\mathrm{mean}}$ peak at $2.0$. So the selector agrees with the coarsest
scale and disagrees with the middle one, which is a sharper statement than the
blanket disagreement an earlier version of this grid produced and points at the
same conclusion.

That conclusion is the one the wider reading forces. The two truth-bearing
scales do not agree with each other: the middle scale is best at the strictest
penalty and the coarsest at an intermediate one, so no single $\lambda$
maximizes both and a scalar penalty cannot be optimal at every scale on this
grid, whatever selects it. And $F_{1,\mathrm{mean}}$ is not the neutral arbiter
it looks like, since the finest scale is scored for silence and its contribution
recovers as the penalty rises, so the mean rewards a strict penalty partly for
suppressing false positives rather than for finding true edges. The external
half of the table is unaffected throughout, since those networks are tiled flat
and scored at one scale by construction.

The reason is visible elsewhere in these results and is not special pleading.
Held-out likelihood does not order the methods the way structural distance does
on this grid, at a Spearman correlation of $-0.333$ between the two orderings.
A spurious lagged self edge is a genuinely good predictor, because the
autocorrelation it absorbs is real, so a criterion that rewards prediction
prefers denser graphs than the truth wherever a persistent latent state drives
the observations. The tiled networks have no such state, which is why the same
rule works there. The honest statement is therefore conditional: selecting the
penalty on held-out evidence is a principled rule that uses no ground truth,
and it recovers the maximizing setting when prediction and structure agree,
which is the case without a latent confound and not with one. That is a
limitation of held-out evidence as a selector rather than of the score being
selected, and it means an analyst who wants precision on data with a latent
regime should declare a strict penalty and disclose it rather than select one.

\paragraph{The imposed block length is not incidental.}
A tiled network is flat, so the second scale is the analyst's construction. The
finest-scale term of the score is a within-block allocation, so the evidence it
carries about a lagged effect grows with the block length while the number of
blocks falls. Measured at $\lambda_k = 0.5$ on the four networks, $F_1$ rises
from $0.511$ at block length $2$ through $0.601$ and $0.763$ at lengths $5$ and
$10$ to $0.812$ at length $20$, and is unchanged at $40$ and $80$. The block length is declared as the square root of the series length,
which lands inside that plateau and is fixed before any graph is scored.

\paragraph{Subsample stability is measured and not adopted.}
Selection over five contiguous subsamples in the sense of
\citet{meinshausen2010stability} raises external precision to $1.000$ at every
threshold tried, $0.6$, $0.8$ and $1.0$, but costs recall, giving $F_1$ of
$0.831$, $0.812$ and $0.802$ against $0.869$ for the single fit. Held-out
evidence agrees with $F_1$ here too, preferring the single fit at $-68070.7$
against $-68066.7$ for every threshold, so the extra layer is reported and left
out.

\subsection{The search is not what limits recovery}
\label{sec:results:search}

Greedy forward-backward accepts one edge at a time, so a pair of edges that
pays only jointly is unreachable from the incumbent set. Because the score
decomposes per target and every candidate is lagged
(Proposition~\ref{prop:learn_decomposition}), the per-target optimum can be certified
by enumeration whenever the in-degree bound keeps the subset count small, which
it does on both benchmarks: $232$ parent sets for an eleven-node network at
in-degree three, $2626$ for the grid. An exhaustive search and a genetic search
over parent sets were both implemented and run in place of the greedy one.

All three return the same parent sets on every cell of both benchmarks, so
every structural column agrees exactly: $F_1$ of $0.725$ and $0.812$ externally
at $\lambda_k$ of $1.0$ and $0.5$, and $F_{1,k=2}$ of $0.647$ and $0.546$ on
the grid at $2.0$ and $1.0$. The grid figures here run over the four-scenario
subset this comparison uses, which is why they sit below the corresponding rows
of \cref{tab:results:frontier}, an average over every scenario. The held-out losses differ in the fourth
significant figure, since the evolutionary search draws its own randomness and
the fitted coefficients are not bit-identical, but the graphs are. The greedy answer is therefore already
the certified per-target optimum on these data, and no search can improve it.
What limits recovery is the score and the protocol around it, not the
optimizer, which is why the preceding subsection is where the gain is.

\subsection{Aggregation consistency and the recorded objective}
\label{sec:results:invariants}

Two properties are checked on every fit rather than sampled. The first is that
the learned model reconciles across scales, so that for every coarse epoch $c$
at scale $k+1$ and every label $\ell$,
\begin{equation}
X_c^{(k+1,\ell)} \;=\; \sum_{t \in \child_k(c)} X_t^{(k,\ell)},
\qquad
Y_c^{(k+1,\ell)} \;=\; \sum_{t \in \child_k(c)} Y_t^{(k,\ell)} .
\label{eq:results:reconcile}
\end{equation}
The reported reconciliation error is the absolute violation of
Equation~\ref{eq:results:reconcile} summed over coarse epochs, labels and
adjacent scale pairs, so any discrepancy anywhere makes it positive. Across all
$1176$ scored fits of the five grid experiments, and separately on the observed
log below, it is exactly $0$ for both counts and durations, on every row, for
every method including the ablations and the eight external learners. The
equality holds by construction of the allocation and is verified rather than
assumed.

The second is that the recorded objective never decreases. Over the same five
experiments, $1056$ fits recorded a history of at least two outer sweeps,
totaling $4534$ successive steps. Not one of them is a decrease, and the
largest recorded step-to-step change in the wrong direction is exactly $0$. The
rejected-sweep counter, which records a proposal the objective refused, averages
between $0.94$ and $1.72$ per fit across the five experiments and is positive on
$826$ of the $1176$ fits, so the guard is active rather than vacuous.

\paragraph{An observed log.}
The model was also fitted once to an observed event log with no ground truth
(\texttt{results/exp6/exp6\_results.json}), $1601$ events over three labels, one
minute per finest slot, branching $10$ then $6$, so $T_1 = 1440$, $T_2 = 144$
and $T_3 = 24$. Both reconciliation errors are exactly $0$, the recorded
objective rises over all five sweeps, from $7002.16$ to $25044.83$, and one
sweep was rejected. The structure search selected no edge at any scale. The
subsample stability check, which refits on four resamples and records how often
each edge is reselected, returns one candidate at scale $3$, the lag-one edge
from the second label to the third, selected on all four resamples but not on
the full fit. That disagreement between the full fit and the resamples is
reported as it stands; on a log of this length it is a signal to collect more
data, not an edge.

\subsection{An observed log with no ground truth}
\label{sec:results:arm1}

The three preceding groups of tables score recovery against a truth the
generator wrote or a manifest a benchmark declared. This subsection reports the
one arm with neither. The log is LifeSnaps
(\texttt{data/observed\_lifesnaps\_event\_log.csv}), a wearable and
self-report study whose participants enter on their own dates. The panel is
each participant's first eight complete calendar weeks measured from that
participant's own first Monday, which all $51$ of the kept cohort clear, with a
mean observed share of $0.985$ and a minimum of $0.88$. Six behavior labels are
targets, nine columns are exogenous, and the four scales carry $80\,640$,
$1344$, $56$ and $8$ epochs. Three survey columns collapse to one node, since
\texttt{notification\_read} equals the sum of the two survey columns on
$0.971$ of participant-days and no within-group edge is expressible after the
collapse. No claim below is causal. What is claimed is that a graph priced on
this log reaches every scale, that the log drawn back out of the fitted model
reproduces the panel's marginal rates and its coarse-scale non-stationarity,
and that the empty graph does not.

Two stored runs carry the arm. The fit run
(\texttt{results/exp15/exp15\_results.json}) searches one graph over the panel
and grades it on a falsification ladder. The round-trip run
(\texttt{results/exp16/exp16\_results.json}) refits the same search, draws five
synthetic logs per leg from the fitted model, writes them in the raw event
schema, reads them back through the same panel reader and compares them with
the panel on an enumerated grid of non-stationarity forms. The two runs accept
the same $23$ edges, edge for edge, so the logs are drawn from the graph the
ladder grades.

\paragraph{Both directions of the pass are used.}
\Cref{prop:model_block_invariance} splits every per-scale parameter block into
the component a scale identifies and the component it cannot, and the two
directions of the learning pass carry one component each. A search that only
ever bought at one scale would leave that split untested on real data.
Table~\ref{tab:results:arm1_scales} reports where the $23$ accepted edges sit.

\begin{table}[t]
\centering
\caption{The accepted graph by scale and by source, on $180$ priced
candidates at an accepted price of $21.31$ nats against a permutation null
whose median is $0.51$ and whose maximum is $1604.0$. A self edge reads a
label's own past; an exogenous edge reads a calendar column or the scheduled
notification; a cross-label edge reads a different behavior label. The price
strata are self against everything else, so the first two columns together are
the $12$ the run reports as cross. Source
\texttt{results/exp15/exp15\_results.json}.}
\label{tab:results:arm1_scales}
\small
\begin{tabular}{lrrrr}
\toprule
Scale & Exogenous & Self & Cross-label & Total \\
\midrule
$k=1$, minute & $0$ & $0$ & $0$ & $0$ \\
$k=2$, hour   & $12$ & $4$ & $0$ & $16$ \\
$k=3$, day    & $0$ & $4$ & $0$ & $4$ \\
$k=4$, week   & $0$ & $3$ & $0$ & $3$ \\
\midrule
All           & $12$ & $11$ & $0$ & $23$ \\
\bottomrule
\end{tabular}
\end{table}

The search reaches three of the four scales and the split between them is
clean. Every exogenous relation is bought at the hour scale and every edge
bought at the day and week scales reads a label's own past. That is the pattern
the block decomposition predicts. A calendar node is a deterministic function of
the epoch, so by \cref{cor:model_calendar_home} it is identified only where it
varies inside a parent block, and the hour-of-day harmonics have their home
scale where a day is the parent. The scheduled notification is exogenous for a
different reason, that the study delivers it on a schedule the participant does
not set so no behavior precedes it, and it enters at the hour scale because that
is where it varies inside a parent block. What the day and week scales carry
is persistence, which is the component the fine-to-coarse direction supplies and
which a fit at the hour scale alone cannot express. So the two directions of the
pass each buy something on this log, and neither buys what the other does.

The third column of Table~\ref{tab:results:arm1_scales} is the result the
section is least able to soften. \textbf{No edge between two distinct behavior
labels is bought at any scale.} The candidate space carries them and the price
does not exclude them, and none survives the search. The declared anchor set
cannot adjudicate this,
because every positive it declares is the scheduled notification into the survey
node, which is exogenous into a target and not a relation between two behaviors.
Two readings are available and this panel does not separate them. Eight weeks of
$51$ participants may not price a between-label relation above a null whose
maximum is $1604$ nats when the same variance is available from the clock and
from the label's own past, which is where
\cref{sec:results:comparison} shows a latent-free criterion spending its budget.
Or these six labels are conditionally independent given the clock, the
participant level and their own histories, which is a substantive claim about
this cohort and one this arm has no instrument for. The section reports the
count and declines to choose.

The finest scale is expressible and priced and this log does not buy it. A
lag-one self column at the minute scale is identically zero on every scorable
row for four of the six labels, because the risk set excludes exactly the slots
where it would be one, and it is nonzero on $10$ rows of $3\,492\,956$ for a
fifth. Those columns are refused at pool construction on their measured
information, so the price is read on the pool the search actually walks.
A read of one label by another inside the same epoch is admitted at the finest
scale, where one epoch is one slot, and none is bought at $51$ participants,
which repeats what a smaller panel showed. Every lag-zero edge in the tables
above has a calendar source, which is a deterministic function of the epoch and
is never drawn. A same-epoch association is symmetric in its two orientations, so the
target order alone decides which way such an edge is written, and any claim
about direction has to come from outside this fit.

\paragraph{The falsification ladder.}
With no truth to score against, the graph is graded by what it declines to
select. Table~\ref{tab:results:arm1_rungs} reports the rungs, each with its own
denominator printed, since a rate on six is not a performance number.

\begin{table}[t]
\centering
\caption{The falsification rungs of the fit run. R1 and R2 are gates the graph
must not trip; R3 and R9 are read as counts. The declared anchor set is frozen
and hashed before the first scoring run. Source
\texttt{results/exp15/exp15\_results.json}.}
\label{tab:results:arm1_rungs}
\small
\begin{tabular}{llrl}
\toprule
Rung & What it counts & Value & Reading \\
\midrule
R1 & Declared protocol negatives selected & $0/6$ & gate held \\
R2 & Sham-label candidates selected & $1/362$ & $0.003$ false positive rate \\
R3 & Declared positives recovered & $2/6$ & $2$ of the $3$ at the hour scale \\
R9 & Edges reselected on $\ge 0.75$ of $16$ half-panels & $8/23$ & bound of $0.66$ on false ones \\
R7 & Invariant parent subsets for \texttt{gym} at the hour scale & $0/1$ & model rejected \\
\bottomrule
\end{tabular}
\end{table}

R1 holds. Of the $78$ declared protocol negatives, $72$ are structural, because
the scheduled label is an exogenous source and the axis cannot express an edge
into it, so the rung stands on the remaining six and none is selected. R2 is
the informative gate. A sham label is connected to nothing by construction and
one of its $362$ candidates is bought, a false positive rate of $0.003$ inside
the behavioral block. R3 recovers two of the six declared positives. Three of the
six are declared at the minute scale and three at the hour scale, at lags one to
three, and the two recovered are the hour-scale lags one and two; the run also
buys the hour-scale lag zero, which the protocol does not declare. On six the
binomial standard error at one half is $0.204$, so R3 is a gate with its
denominator printed and never a performance number.

R9 is what separates the receipt table from the accepted set. Refitting on
$16$ complementary half-panels, eight of the $23$ edges are reselected on at
least three quarters of them, where the complementary-pairs bound puts at most
$0.66$ low-probability candidates above that threshold. The three
edges with the largest standardized effects in the fit sit among the five with
the lowest reselection rates. Two lag-two self edges into \texttt{nap} at the
hour scale carry standardized effects of $19.2$ and $-145.0$ and are reselected
on one of $16$ half-panels each; their columns clear the support floor and do
not carry enough information to determine what they are given. A coefficient of that size
is a report that its column is nearly empty, and the stability rung is what
removes it. R7 rejects the model for \texttt{gym}. That target carries one
accepted parent at the hour scale, and no subset of it leaves a residual
distribution invariant across participant waves, so the intersection the rung
reports is empty. On one column the rung is a weak test and it is reported with
its width printed.

\paragraph{A calendar node above its home scale buys nothing.}
\Cref{cor:model_calendar_home} is a statement about the kernel and it has a
direct empirical test. Table~\ref{tab:results:arm1_calendar} runs the search
three times on the same panel, with the calendar nodes withheld, offered at
their home scale, and offered at every scale.

\begin{table}[t]
\centering
\caption{The calendar-node ablation, three pre-declared arms on one panel.
The objective is the penalized log-likelihood the search maximizes, so a
larger value is better. Source \texttt{results/exp15/exp15\_results.json}.}
\label{tab:results:arm1_calendar}
\small
\begin{tabular}{lrrrr}
\toprule
Calendar nodes & Candidates & Edges & Objective & vs.\ home \\
\midrule
Withheld          & $172$ & $21$ & $-62\,217.47$ & $-2295.64$ \\
At their home scale & $180$ & $23$ & $\mathbf{-59\,921.83}$ & $0.00$ \\
At every scale    & $204$ & $25$ & $-59\,955.73$ & $-33.90$ \\
\bottomrule
\end{tabular}
\end{table}

Offering the calendar nodes at their home scale is worth $2295.64$ nats over
withholding them. Offering them at every scale costs $33.90$ nats against
offering them at their home scale, while buying two more edges. Both halves of
the corollary hold on this log: the nodes are worth their price where they are
identified, and the price paid above their home scale is wasted. The corollary
turns what is usually a modeling convention into a consequence of the kernel,
and this is the measurement that says the convention is the right one.

\paragraph{The drawn log reproduces the panel's marginal rates.}
The round trip refits the search, draws from the fitted model over the same
roster, mask and window, writes raw events and reads them back. Reconciliation
between counts and durations at adjacent scales is exactly $0$ on every leg and
every replicate, as it is on the synthetic grid, and every drawn log reads back
onto the same $51$ participants and the same observed slots it was drawn over.
Table~\ref{tab:results:arm1_volume} reports the quantity the grid does not
read. Every marginal form of the grid but one divides by the series mean, and
the one exception reads the share of empty epochs, so a leg drawing a label at
half its real rate can still pass every cell it is tested at.

\begin{table}[t]
\centering
\caption{Episode counts on the real panel against the mean over five drawn
replicates, restricted to the same observation mask. The drawn column is the
fitted leg; the last three columns are the drawn-to-real ratio on the fitted
leg, on the leg with the calendar nodes withheld, and on the empty graph. The
summary row is the mean absolute log ratio over the four labels the panel
supports well, which is symmetric in the direction of a miss. Source
\texttt{results/exp16/exp16\_results.json}.}
\label{tab:results:arm1_volume}
\small
\begin{tabular}{lrrrrr}
\toprule
 & & & \multicolumn{3}{c}{Drawn-to-real ratio} \\
\cmidrule(lr){4-6}
Label & Real & Drawn & Fitted & No calendar & Empty \\
\midrule
\texttt{gym} & $506$ & $498.6$ & $0.99$ & $1.00$ & $1.14$ \\
\texttt{nap} & $214$ & $218.4$ & $1.02$ & $1.12$ & $1.24$ \\
\texttt{sleep} & $1841$ & $1919.2$ & $1.04$ & $1.08$ & $1.12$ \\
\texttt{walk} & $1243$ & $1321.0$ & $1.06$ & $0.95$ & $1.10$ \\
\texttt{mindfulness} & $68$ & $132.2$ & $1.94$ & $1.80$ & $1.18$ \\
\texttt{survey\_response} & $10\,162$ & $5418.4$ & $0.53$ & $0.57$ & $0.56$ \\
\midrule
Mean $|\log$ ratio$|$, top four & & & $\mathbf{0.034}$ & $0.061$ & $0.137$ \\
\bottomrule
\end{tabular}
\end{table}

Four of the six labels land within $6$ percent of their real volume. The two
that do not are the two the panel supports least well, for different reasons.
\texttt{mindfulness} occurs $68$ times over $2688$ observed participant-days,
which is one episode per $40$ days, and the draw doubles it. The
\texttt{survey\_response} node carries $72$ percent of the panel's behavioral
episodes and $98.9$ percent of them are zero-duration, so it is a count-only
label whose duration block is uninformative, and the draw reaches $53$ percent
of its rate. Both misses are reported as they stand.

Read as one number, the mean absolute log of the drawn-to-real ratio, which
charges a miss the same in either direction, the volumes are where the graph
separates from the empty one. Over the four well-supported labels the fitted
leg sits at $0.034$, withholding the calendar nodes at $0.061$ and the empty
graph at $0.137$, so the graph is four times closer to the panel than an
intercept and a participant level, and the calendar nodes are half of that
gain. Over all six labels the three legs are $0.239$, $0.232$ and $0.216$, and
the ordering reverses, because the empty graph does not overshoot the label it
has no information about. The two readings are reported together because the
second is the reason for stating which labels the first is read over.

The remaining departure is the share of days left empty. The real panel leaves
$456$ of its $2688$ observed days with no behavioral episode, $17.0$ percent;
the fitted leg leaves $142.8$ per replicate, $5.3$ percent, the no-calendar leg
$133.6$ and the empty graph $89.8$. The ordering is the same one the volumes
give, and no row of the grid reads this quantity.

\paragraph{Non-stationarity is preserved at the coarse scales.}
The comparison is an enumerated grid of $21$ forms over six labels and three
calendar scales, $228$ cells, each an equivalence test at its own declared
tolerance with the participant as the sampling unit. A cell is put to a test
only when the real panel supports one, and $182$ of the $228$ are descriptive
for a declared reason, $127$ of them because the tolerance is inside the
precision the panel can resolve. Table~\ref{tab:results:arm1_grid} reports the
four legs.

\begin{table}[t]
\centering
\caption{The non-stationarity grid, four legs of five replicates each. The
fitted leg carries the searched graph and an hour-scale regime chain; the
no-regime leg is the same search with the chain given one state; the
no-calendar leg is the same search with the clock nodes withheld; the empty
leg carries an intercept and a participant level. Source
\texttt{results/exp16/exp16\_results.json}.}
\label{tab:results:arm1_grid}
\small
\begin{tabular}{lrrrr}
\toprule
Leg & Edges & Tested cells & Failing cells & Failing rows \\
\midrule
Fitted        & $23$ & $46$ & $12$ & $5$ \\
No regime     & $23$ & $46$ & $13$ & $6$ \\
No calendar   & $21$ & $46$ & $12$ & $5$ \\
Empty graph   & $0$  & $46$ & $14$ & $6$ \\
\bottomrule
\end{tabular}
\end{table}

The panel fails the empty graph, which is the control that says it can fail at
all. It also fails the fitted leg, since the verdict is a conjunction over every
tested cell and twelve of the $46$ do not pass.

The row-level reading is the informative one and it needs the positive control
first. Grading one drawn replicate of the fitted leg against the other four,
which nothing separates, five of the $11$ tested rows fail: the
seasonality-amplitude drift, the dispersion, the factor share, the same-epoch
dependence drift and the lagged dependence drift. A row that fails its own
replicate noise has a tolerance narrower than the spread of the draw, so the
fitted leg failing that same row says nothing about the fit. Six rows carry
evidence once those five are removed, and \textbf{the fitted leg passes four of
them}: the normalized trend, the unit-root pair, the share of Haar detail energy
the coarsest level holds, and the windowed maximum mean discrepancy. It fails
the within-day harmonic profile and the drift in persistence.

The same reading bounds what the comparison against the empty graph can carry.
The fitted leg fails two cells fewer of $46$, and the one row that separates
them is the drift in same-epoch dependence, which is one of the five that fail
their own noise. \textbf{On the six rows that carry evidence the two legs are
identical}, passing the same four and failing the same two. So this grid, on
this roster, does not separate a fitted graph from an intercept and a
participant level, and the quantity that does is the volume comparison above,
which the grid does not read.

Table~\ref{tab:results:arm1_cuts} reads the same question as a segmentation.
Both logs are cut independently at their own changepoints under a Poisson cost
and the two cut sets are compared.

\begin{table}[t]
\centering
\caption{Changepoint agreement between the real panel and the first drawn
replicate of the fitted leg, per label and per scale, under a Poisson cost.
Each participant's two series are cut independently and matched, and both
columns are means over the $51$ participants. $F_1$ is scored inside a margin
of three epochs; annotation error is the absolute difference in the number of
cuts. Source \texttt{results/exp16/exp16\_results.json}.}
\label{tab:results:arm1_cuts}
\small
\begin{tabular}{lrrrrrr}
\toprule
 & \multicolumn{2}{c}{Week} & \multicolumn{2}{c}{Day} & \multicolumn{2}{c}{Hour} \\
\cmidrule(lr){2-3}\cmidrule(lr){4-5}\cmidrule(lr){6-7}
Label & $F_1$ & Ann.\ err. & $F_1$ & Ann.\ err. & $F_1$ & Ann.\ err. \\
\midrule
\texttt{gym}             & $0.961$ & $0.06$ & $0.784$ & $0.29$ & $0.765$ & $0.41$ \\
\texttt{mindfulness}     & $0.954$ & $0.06$ & $0.915$ & $0.10$ & $0.890$ & $0.16$ \\
\texttt{nap}             & $0.961$ & $0.04$ & $0.922$ & $0.08$ & $0.902$ & $0.12$ \\
\texttt{sleep}           & $0.784$ & $0.24$ & $0.667$ & $0.41$ & $0.627$ & $0.51$ \\
\texttt{walk}            & $0.706$ & $0.29$ & $0.588$ & $0.49$ & $0.471$ & $0.76$ \\
\texttt{survey\_response} & $0.601$ & $0.53$ & $0.355$ & $1.41$ & $0.054$ & $6.27$ \\
\bottomrule
\end{tabular}
\end{table}

Agreement is highest at the week scale and falls monotonically toward the hour
scale for all six labels. At the week scale five of the six reach $F_1$ above
$0.70$, three of them above $0.95$, and those five miss the panel's cut count by
under a third of a cut. At the hour scale three labels hold above $0.75$ and
\texttt{survey\_response} collapses to $0.054$ at an annotation error of $6.27$
cuts. So the departure grows toward the fine scale for every label, and it is
larger by an order of magnitude for the one whose duration block carries
nothing. That agrees with the grid, whose two evidence-bearing failures are the
within-day harmonic profile and the drift in persistence.

Between-participant dispersion is read as the coefficient of quartile
variation over weekly totals, with a percentile bootstrap over participants.
Three of the six labels are comparable, \texttt{nap} at $0.333$ against
$0.333$, \texttt{walk} at $0.385$ against $0.385$ and \texttt{survey\_response}
at $0.121$ against $0.243$ on overlapping intervals. \texttt{sleep} separates
by $0.021$. \texttt{gym} has overlapping intervals and is called incomparable
because the two panels define the measure for shares of the roster $0.388$
apart. \texttt{mindfulness} is refused: every participant defining a value for
it is pinned at $1.0$ by a zero lower quartile, so the label carries no
dispersion to compare.

\paragraph{What the controls establish and what they do not.}
Four negative controls are declared against the grid. The no-regime control
does its job, breaking every row it is declared against with nothing spurious.
The other three do not. Withholding the calendar nodes leaves the
seasonality-amplitude drift unbroken; shuffling labels across the cross-section
leaves the same-epoch dependence drift unbroken; and a within-participant time
shuffle, which has to break every marginal row, leaves five of them unbroken
and eight untested. A control that fails to break a row it was declared against
says the grid cannot reach that row on this roster, and the honest reading is
that $46$ tested cells of $228$ is the sensitivity this panel supports. That
bound, and not the fitted model, is what the three failing controls measure.
It is also why the row-level reading above is reported net of the positive
control and why the volume table is reported at all.


\subsection{Cost}
\label{sec:results:cost}

Wall-clock cost is reported for the fits as run, one process per fit, with
thread-level parallelism pinned to one so that the numbers are comparable
across methods.

On the synthetic grid a full fit averages $72.1$ seconds with a median of
$60.8$ and a maximum of $203.0$, over a mean of $5.4$ outer sweeps. The empty
graph, which runs the same inference on a fixed graph, averages $14.3$ seconds,
and the oracle-parameter row, which runs inference once and fits nothing,
averages $1.6$. The five grid experiments together account for $96\,886$
seconds of fitting, of which the ten-method comparison of
Section~\ref{sec:results:comparison} is $68\,361$.

The cost is dominated by the structure search rather than by inference; on the
grid the empty-graph row runs the same inference on a fixed graph for $14.3$
seconds against $72.1$ for the full fit. On the external benchmark the local
score costs $174.7$ seconds per fit on average and $461.7$ at most, against
$6.6$ for exhaustive search, $0.3$ for either hill climbing variant, $0.02$
for Granger and $0.003$ for DYNOTEARS. Most of that difference is the selection of
\cref{sec:results:selection}, which fits ten candidate configurations before
the one it keeps, so a single fit at a declared setting costs about a tenth of
the reported figure. Even so the local score is between twenty six times and
five orders of magnitude more expensive than the classical entries, exhaustive
search being the nearest of them and DYNOTEARS the furthest, and its position
in Table~\ref{tab:results:external} does not depend on that. Cost is not the
reason to prefer any row in these tables, and it is reported so that the
comparison is not read as free.

\section{Limitations and conclusion}
\label{sec:limitations}

\subsection{Same-epoch co-movement leaks into finest-scale lagged edges}
\label{sec:lim_leak}

The clearest failure of the model is that same-epoch structure is partly reported
as lagged structure at the finest scale, and the reported edges are lagged self
edges. The benchmark grid places no true edge at the finest scale, so every edge
returned there is a false positive and the rate is directly measurable. Over the
66 fits of the ideal-recovery grid, the local score returns 2.212 spurious
finest-scale edges per fit, 3.244 across the fits that return any at all, and
leaves the finest scale entirely clean in 21 of them. Over the 54 fits of the
algorithm comparison the same leak is a saturation of 0.293 of the one-parent
budget per label. That is far better than every classical entry, five of which
use the entire budget for every label in every fit and three of which reach
0.770 to 0.967, and it is not zero.

Two qualifications belong with those rates. The first is that the grid was
reshuffled between this run and the one reported before it, which moved true
edges onto the coarsest scale, so the earlier finest-scale rate of 1.106 edges
per fit and the present 2.212 are not a like-for-like pair. What can be said is
that the direction is against the model. The second is that the finest scale is
now known to be unable to carry a findable edge at all. Run at a nonzero
in-degree on every scale, the screen of \cref{sec:benchmark_identifiability}
puts the median gain of a true finest-scale edge below the floor it has to clear
in all three families, by 3.40, 9.46 and 6.27 nats, with three of twenty-one
scenarios clearing, because the finest scale is a binary event-start process and
a lagged edge there can only move which minute of a block begins an episode. A
method that returns an edge there is reporting structure the data does not
contain, and no measurement on this grid separates that from correctly finding a
finest-scale edge, since none can be placed.

Four candidate causes were tested. Three were excluded by measurement and the
fourth is what survives. The first candidate was the latent factor, on the
reading that the leak is co-movement a factor would otherwise absorb. It is
excluded, because the scenario families that carry a factor and force a single
regime leak at 0.289 in saturation, essentially the 0.283 of the family that
carries no latent state at all. The second candidate was the regime, whose
persistence autocorrelates every label it drives. The same measurement excludes
it, since the families that carry a regime and no factor leak at 0.296, the same
rate again, and the stationary family, which has neither a regime nor a factor,
still leaks at 0.283 over its 12 fits, 6 of which return a finest-scale edge. On
the grid used before this one the regime carried a visible share of the
difference between the families and was excluded only as the sole cause; on this
one the families do not differ, so there is no difference left for it to carry.
The third candidate was the finest-scale cap that constrains counts to be binary,
on the reading that the allocation likelihood does not model the redistribution
the cap forces. It is excluded, because a probe with the cap removed leaks more
rather than less.

What survives all three exclusions is within-block competition for a fixed parent
total. The children of a parent epoch share a total that does not change, so a
child whose sibling has already taken counts has fewer left to take, and a lagged
feature of a sibling in the same block is a proxy for how much of the block total
is already used up. Three properties of the measurement identify this mechanism
rather than any of the excluded ones. The effect is present within a block and
absent across block boundaries; the coefficient the search fits for the spurious
edge is negative, which is the sign competition predicts and the sign a genuine
excitatory dependence does not; and the effect survives in the family with no
latent state at all, because the block total is fixed there too. Adding an
explicit depletion covariate to the allocation score, an exclusive cumulative sum
of the counts already taken by earlier children of the same block, captures a
large amount of real structure and improves held-out fit, and it does not remove
the finest-scale false positives of the benchmark. The gains those surviving
edges carry sit an order of magnitude above anything a block-preserving surrogate
of the same feature earns, which places their source in the part of the
finest-scale regime path the posterior cannot identify rather than in
competition.

The refactorization that looks like the obvious closure is not one. The
Dirichlet-multinomial log-likelihood of one label in one block, with parent total
\(N\), child counts \(x_t\) summing to \(N\), allocation probabilities
\(\omega_t\) summing to one over \(t \in \child_k(c)\), and concentration
\(\kappa\), can be written as a sequential product of beta-binomial kernels, one
per child given its realized predecessors, which is exactly the form an exact
per-child normalizer would need. Ordering the children arbitrarily and setting
\(a_t = \kappa \omega_t\),
\(b_t = \kappa \sum_{s>t} \omega_s\)
and \(N_t = N - \sum_{s<t} x_s\), the identity is
\begin{align}
&\sum_{t \in \child_k(c)}
\Bigl[
\log \Gamma\!\left(x_t + \kappa\omega_t\right)
- \log \Gamma\!\left(\kappa\omega_t\right)
\Bigr]
+ \log \Gamma(\kappa)
- \log \Gamma(N + \kappa)
\nonumber\\
&\qquad =
\sum_{t \in \child_k(c)}
\Bigl[
\log B\!\left(x_t + a_t,\, N_t - x_t + b_t\right)
- \log B\!\left(a_t, b_t\right)
\Bigr],
\label{eq:lim_betabinomial}
\end{align}
where the last child in the chosen order has \(b_t = 0\) and \(x_t = N_t\), so
the two divergent beta terms it contributes cancel and its summand is zero, which
is the statement that the last child takes whatever the block has left. The
identity holds for every \(\omega\), including an \(\omega\) that depends on the
data, and the combinatorial factors the two sides omit are equal to each other,
since \(\prod_t \binom{N_t}{x_t} = N! / \prod_t x_t!\). Both sides are therefore
the same function of the parameters, with the same gradient, the same maximizer
and the same edge gain. A direct evaluation of the two sides agrees to
floating-point rounding. Refactorizing the block changes no
score, so it is not a route to an exact normalizer, and an earlier claim in this
project that it was such a route is withdrawn.

\subsection{The allocation normalizer is evaluated under a mean field}
\label{sec:lim_meanfield}

The allocation probabilities of a block couple all of its children through a
single normalizer, and the implementation evaluates that normalizer by scoring
each child at its own per-regime score while holding its block-mates at the
regime-averaged score. The consequence is precise. The parameter step accepts a
block when the objective computed at the regime posterior it was given does not
fall, and the next latent-state inference step recomputes that objective at the
new posterior. The two are not the same function, so a block that raised the
first can lower the second, and the second is what the outer loop monitors.

The remedy shipped here is an outer guard rather than an exact normalizer. A
sweep that ends below the best sweep so far is rejected, the state reverts to the
best sweep, and the sweep is retried with a halved parameter learning rate up to
a bounded number of times, after which the fit stops at its best state. The guard
works as specified. Over the 1056 fits of the reported runs that execute more
than one sweep, the recorded objective decreases in none, against a decrease in
81.1 percent of the corresponding fits of the run made before the guard existed.
The residual is contained and not removed, and the count of rejected sweeps is
the measure of how much of it there is. The guard rejected at least one sweep in
826 of the 1176 reported fits, at most three in any fit. Removing the residual
needs a normalizer that is exact under the variational family, and the exact
expectation of the block log-normalizer under the joint posterior over the
children's regimes has a number of terms exponential in the block size, so it is
not available at the branching factors used here.

\subsection{Recall on external data}
\label{sec:lim_recall}

On released Bayesian networks tiled into two-slice temporal networks, hill
climbing under the Bayesian information criterion still recovers more edges than
the local score does. Most of the gap the first version of this experiment
showed was the protocol rather than the criterion, and the part that survives
correcting the protocol is smaller and has a structural explanation.

Three protocol choices were wrong and all three were ours. Two regimes were
being fitted to data generated with none, so the regime chain was free to absorb
variation the emission owns. The edge penalty was inherited from the synthetic
grid, where a persistent latent state makes a non-edge expensive to reject,
while a tiled network has no such null level. And the block length of the
imposed hierarchy was set to ten, in the flat part of no curve. Selecting the
regime count, the factor dimension and the penalty on held-out log evidence, and
declaring the block length as the square root of the series length, takes $F_1$
on five tiled networks from $0.532$ to $0.806$, precision from $0.750$ to
$0.869$ and recall from $0.414$ to $0.762$. Two of the five are then recovered
as well as hill climbing recovers them, exactly $F_1 = 0.933$ on asia and
exactly $1.000$ on earthquake.

What remains is concentrated rather than uniform. On child, the largest network
scored, the local score reaches $F_1 = 0.870$ against $0.933$ at the higher
precision of the two, $0.958$ against $0.933$. On sachs it reaches $0.812$
against $0.980$; on cancer it reaches $0.413$ against $0.905$.

Averaged over the five networks and three replicates the deficit is about eight
true edges, spread over three networks: $3.33$ on child, $2.67$ on sachs and
$2.00$ on cancer, with asia and earthquake exactly tied. On the single
replicate we examined in detail the deficit is three edges, and we report those
three rather than a mechanism we cannot support. On sachs the two are
\texttt{PKA} to \texttt{Mek} and \texttt{Erk} to \texttt{Akt}. The first is a
budget substitution: the in-degree was saturated at three of three and the slot
went to \texttt{Plcg}, whose conditional gain is $-13.88$. The second is not:
the in-degree was two of three with a slot free, and the score declined to add
an edge whose conditional gain is $+17.36$. On cancer the score returns no
parents at all for \texttt{Cancer}, although \texttt{Smoker} to \texttt{Cancer}
has a sole-parent gain of $+20.91$ and an odds ratio of $15$. A positively
scoring edge that is not taken with room to take it is a defect in the local
score or in what the latent posterior does to it, and we state it as an open
question rather than explain it.

We had proposed that the finest-scale term, being a within-block allocation
conditioned on the block total, cannot see a parent effect that moves the total.
That explanation is wrong and we withdraw it. On the three edges at issue the
within-block association is the whole association: the ratio of within-block to
raw association averages $1.007$ for recovered edges against $0.999$ for missed
on sachs, and $1.003$ against $0.989$ on cancer; conditioning exactly on block
totals retains $99.4$ against $98.7$ percent of the pooled association on sachs
and $97.7$ against $93.5$ percent on cancer; and between-block variation is
$1.1$ to $2.0$ percent of each target's fine-scale variance against $1.6$
percent expected from sampling alone. There is no between-block variation for
an effect to hide in.

Two candidate explanations were tested and rejected. The search is not the
cause, since exhaustive per-target enumeration and a genetic search return the
same parent sets as the greedy one on every cell of both benchmarks. And
error-rate-controlling selection is not the remedy, since subsample stability
raises precision to $1.000$ but costs enough recall to lower $F_1$, and held-out
evidence prefers the single fit.

\subsection{Parameter recovery fails the ladder from either start}
\label{sec:lim_params}

Parameter recovery is a consistency claim, so the test that was stated before
the experiment ran is a sample-size ladder rather than any single value, and the
criterion was that the error must fall at every step. The ladder is set from the
coarsest scale rather than the finest, so its three rungs give that scale $168$,
$336$ and $672$ epochs, and the middle rung is the horizon the structural
benchmark itself runs at. Neither fitted start meets the criterion. From a
ground-truth start the mean parameter error runs $0.2899$, $0.3494$ and
$0.3208$; from a data-driven start $0.4261$, $0.4533$ and $0.4678$, rising at
every rung.

The failure is confined to one generating family. Family III falls at every
rung, $0.2956$, $0.2932$ and $0.2588$, and meets the criterion outright; family
I is close to flat; family II rises at every rung, $0.6461$, $0.7933$ and
$0.8445$, and sets the level. An error that rises with the sample size is a fit
converging to a different solution and describing that better, rather than a
slow one, and family II is the regime-carrying family whose finer latent paths
are unidentified even with the generating parameters installed.

One part of the model's behavior at the coarsest scale has a mechanism and a
correction. A level initialized from the observed mean already contains the
contribution of the lagged parents, and the parameter step refits that level
against a graph with no edges in it before the structure search runs, so the
level absorbs exactly what a true lagged edge would explain and the search meets
a model that already accounts for the data. Correcting both, by initializing the
level net of what its own lag-1 values explain and by searching before the level
is refitted, takes coarsest-scale recovery from $0.00$, $0.00$ and $0.67$ lagged
keys per fit to $4.67$, $4.44$ and $4.67$ against a true $4.44$. Measured against
the same fits without the correction, parameter error improves in seven of the
nine family-by-rung cells and worsens in one, by $0.044$. Family II's rise is
present in both, so it is not caused by the correction.

The practical consequence is unchanged. Fitting an observed log supports
prediction, exact reconciliation and stability selection, and does not yet
support a claim about recovered parameters, because the criterion set in advance
is not met from the only start a real log admits.

\subsection{What no external benchmark can test}
\label{sec:lim_scope}

The scope limit follows from \cref{sec:related_benchmarks} and cannot be
engineered away within this paper. No released benchmark combines event counts,
exact cross-scale aggregation and a latent regime, so the property the model is
built around cannot be scored against an external ground truth. The tiled
networks test the structure learner, the local score and the table estimator on
flat discrete data, and the two-scale hierarchy imposed on them by summing blocks
of epochs is an artifact of the experiment rather than a property of the data.
Every multiscale number in this work therefore comes from a generator written for
it, and the defense of that generator is the identifiability screen rather than
an appeal to realism.

What a real log can test is prediction, exact reconciliation, the stability of a
selected edge across subsamples, and the shift confirmation used inside the
structure step, since none of these needs a known graph. The self-compatibility
check of \citet{faller2024selfcompatibility} adds a further falsification test
that requires no ground truth, and it is necessary rather than sufficient.
Recovery of a graph, of a regime path or of a parameter cannot be measured on
real data at all, and no claim of that kind is made here from real data.

\subsection{What the observational arm does not settle}
\label{sec:lim_arm1}

The observational arm of \cref{sec:results:arm1} is the one place a real log is
put to the model, and four of its results are bounded in ways a reader has to
carry.

The grid's reach is a property of the roster. Of $228$ cells, $46$ are put to a
test and $127$ of the rest are descriptive because the declared tolerance is
inside the precision $51$ participants over eight weeks can resolve. Three of
the four declared negative controls fail to break at least one row they are
declared against. A control that cannot break a row is measuring what the grid
can reach on this panel, so $46$ tested cells is the sensitivity of the
comparison and not a property of the fitted model. A larger roster would move
that number and might move the verdicts with it.

Five of the $11$ tested rows fail their own replicate noise, which is why the
row-level reading of \cref{sec:results:arm1} is reported net of the positive
control and the raw count of $12$ failing cells of $46$ is not the figure to
read. A tolerance narrower than the spread between two draws of the same fitted
model grades the draw's variance and says nothing about the fit. The same
reading bounds the comparison against the empty graph. The fitted leg fails two
cells fewer of $46$, and the single row that separates the two is one of the
five that fail their own noise, so on the six rows that carry evidence the
fitted graph and an intercept with a participant level pass and fail exactly the
same ones. The grid is not what separates them on this roster.

Two of the six labels are drawn at rates the panel cannot pin down.
\texttt{mindfulness} occurs once per $40$ observed participant-days and is
drawn at $1.94$ times its rate; \texttt{survey\_response} is a count-only node
whose duration block is uninformative and is drawn at $0.53$. The drawn logs
also leave $5.3$ percent of observed days with no episode against the panel's
$17.0$ percent, a marginal quantity no row of the grid reads. These are
reported as measured and none of them is repaired by a correction inside the
model as it stands.

No edge between two distinct behavior labels is bought at any scale, so the
graph this arm reports is a set of clock and scheduling effects and a set of
per-label persistences. That is a negative result about this panel and it is
reported as one. Nothing in the declared anchor set adjudicates it, since every
declared positive is the scheduled notification into the survey node, and the
arm supplies no instrument that separates a between-label relation this cohort
does not carry from one eight weeks of $51$ participants cannot price above the
null.

The graph the arm reports carries two edges whose coefficients cannot be
interpreted. Two lag-two self edges into \texttt{nap} at the hour scale reach
standardized effects of $19.2$ and $-145.0$ on columns that clear the support
floor and carry far too little information to determine what they are given.
They are reselected on one of $16$ half-panels each, so the stability rung
removes them, and the support floor is the wrong instrument for this case. The
invariance rung rejects the model for \texttt{gym}, on the one accepted parent
that target carries at the hour scale, so it is a weak test and it is reported
with its width printed. Neither result is repaired here, and both are reported
beside the graph they qualify.

\subsection{Conclusion}
\label{sec:conclusion}

This work establishes five things.

A single model can meet the four requirements of
\cref{sec:intro_requirements} at once. Exact aggregation between a fine count and
its coarse summary holds by construction rather than by repair, and the measured
reconciliation error is zero for counts and durations in every reported fit,
including a fit on an observed log. Lagged dependence is learned per scale under
a penalized local score rather than assumed. Persistent drift and same-epoch
co-movement are carried by a regime and a factor respectively, so that
co-movement is not reported as a directed dependence. And the finest scale maps
back deterministically to a dated event trace.

Scoring a candidate parent set against a model that already contains the latent
variables changes what the search finds. On a benchmark screened so that its
scenarios contain what the metrics claim to measure, over 540 fits spanning 18
scenarios at three replicates and ten methods, the latent-conditional score
reaches \(F_1 = 0.690\) at the ten-minute scale and \(0.892\) at the hourly
scale, the two that carry true edges, against 0.109 to 0.251 and 0.021 to 0.346
for eight published learners, and it is the only entry whose total structural
Hamming distance, 5.963, beats the 10.981 of reporting nothing. Every classical
entry lands between 1.98 and 2.92 times that reference. The mechanism behind the
gap is a confound any latent-free criterion shares, and recall split by edge kind
states it directly rather than leaving it to be inferred. At the ten-minute scale
seven of the eight latent-free entries recover between 0.679 and 1.000 of the
true lagged self edges, and none of the eight recovers more than 0.245 of the
cross-label ones, while the latent-conditional score recovers 0.919 and 0.523 and
is the only entry above a quarter on cross edges.

On external ground truth the ordering depends on how the two model classes are
configured, which is itself a finding about the comparison rather than about
either criterion. A latent-variable model carries free quantities that a
latent-free learner does not, and holding them at declared values while the
classical learners choose their own penalty through their own score measures
the declarations. Selecting the regime count, the factor dimension and the edge
penalty on held-out log evidence from the training split raises $F_1$ on the
tiled networks from $0.532$ to $0.806$ when each dataset selects for itself,
and to $0.869$ at a precision of $0.983$ when one penalty is chosen for the
benchmark. On the synthetic grid the same rule does not select the maximizing
penalty. It picks $0.5$, where $F_1$ at the ten-minute scale is $0.610$, while
that $F_1$ peaks at the strictest setting of $2.0$ with $0.691$, which is also
the best setting on precision, $0.900$, and on total distance, $6.82$. Held-out
likelihood does not order structure where a persistent latent state
makes a spurious edge a good predictor, which is a limitation of the selector
rather than of the score it selects for. The frontier was once scored on the
ten-minute scale alone, which left open whether the disagreement was an artifact
of reading a two-scale truth through a one-scale metric. Rescoring at the
coarsest scale and on the mean closes that: both disagree with the selector as
well, at $1.0$ and $2.0$ against the selected $0.5$. What the wider reading adds
is that the two truth-bearing scales do not agree with each other, so no scalar
penalty is best at both. Hill climbing under BIC still leads on the
tiled networks, by a margin in recall, so the statement is that the gap is
smaller and one-sided rather than that the two benchmarks contradict each other.

A model fitted to a real log draws a log a panel of non-stationarity tests
cannot easily tell from it. On $51$ LifeSnaps participants over eight complete
calendar weeks each, the fitted graph draws four of six labels within $6$
percent of their real episode volume at exactly zero reconciliation error and
four times closer than an intercept with a participant level draws them; net of
a positive control that names the rows failing their own replicate noise, it
passes four of the six grid rows that carry evidence. On those six rows the
empty graph passes and fails exactly the same ones, so the grid at this roster
size bounds what it can resolve and the volume comparison is what separates the
two. Changepoint agreement between the two
logs is monotone in the calendar scale for all six labels, from $0.95$ to
$0.96$ at the week scale for three of them down to $0.05$ at the hour scale for
the one count-only label, which places the departure at the within-day shape
and nowhere coarser. The same panel confirms
\cref{cor:model_calendar_home} in both of its directions: offering the calendar
nodes at their home scale is worth $2295.64$ nats over withholding them and
offering them at every scale costs $33.90$ nats against offering them at their
home scale. On that log every cross-label edge the search buys sits at the hour
scale and every edge it buys at the day and week scales reads a label's own
past, which is the block decomposition of
\cref{prop:model_block_invariance} appearing as a property of a fit.

Finally, the failures are located rather than described. The finest-scale leak
has three candidate causes excluded by measurement and one algebraic dead end
disproven; the mean-field normalizer is contained by a guard whose residual is
counted rather than assumed away; the residual external recall gap is
quantified after selection rather than argued away; and parameter recovery is
reported as failing a criterion stated in advance from either fitted start, with
its candidate cause named as a hypothesis to test rather than as a conclusion.
On the observed log the search buys no edge between two distinct behavior
labels at any scale, and the declared anchor set contains no such relation to
adjudicate it, so that count is reported with the two readings it admits and
with the statement that this arm carries no instrument to choose between them.


\bibliography{n2dbn_refs}

\end{document}
